\documentclass[10pt,a4paper]{article}
\usepackage[a4paper,margin=1.8cm]{geometry}
\usepackage{fontspec}
\usepackage{polyglossia}
\setmainlanguage{russian}
\setotherlanguage{english}
\setmainfont{CMU Serif}
\setsansfont{CMU Sans Serif}
\setmonofont{CMU Typewriter Text}
\usepackage{amsmath}
\usepackage[table]{xcolor}
\usepackage{booktabs,longtable,array,graphicx,pdflscape,hyperref,enumitem,fancyhdr,titlesec,needspace}
\definecolor{hm}{RGB}{31,111,235}
\definecolor{yes}{RGB}{26,127,55}
\definecolor{part}{RGB}{210,153,34}
\hypersetup{colorlinks=true,linkcolor=blue!50!black,urlcolor=blue!50!black,citecolor=blue!50!black}
\urlstyle{tt}
\setlist{itemsep=1pt,topsep=3pt}
\setlength{\parindent}{1.0cm}
\setlength{\parskip}{0pt}
\renewcommand{\arraystretch}{1.12}
\sloppy\emergencystretch=3em
\setcounter{tocdepth}{2}
\setcounter{secnumdepth}{3}
\titleformat{\subsubsection}{\normalfont\normalsize\bfseries}{\thesubsubsection}{0.6em}{}
\pagestyle{fancy}\fancyhf{}
\fancyhead[L]{\small Обзор инструментов и сценариев}
\fancyfoot[C]{\thepage}
\renewcommand{\headrulewidth}{0.3pt}
\newcommand{\lab}[1]{\par\noindent\textbf{#1.}\ }
\title{Подробный обзор инструментов оценивания конечных смесей\\и прикладных сценариев}
\author{}
\date{}
\begin{document}
\maketitle

\begin{abstract}\noindent
В документе подробно разобрана эмпирическая база обзора: 27 инструментов оценивания конечных смесей, 86 требований к библиотеке
и 10 прикладных сценариев. Документ не повторяет основной текст работы, а дополняет его: содержит матрицы покрытия по каждой категории требований, сопоставление разнородных средств,
анализ чувствительности оценок, сходство инструментов по профилю возможностей, расчёт приоритетов развития \texttt{pysatl-mpest}. Описания отдельных инструментов, требований и сценариев,
обоснование и ссылка для каждой из 2322 оценок, а также связанные публикации приведены в дашбордах и здесь не дублируются. Все числа вычислены из таблиц базы, на которой построены два дашборда.
Интерактивные версии базы: дашборд инструментов и требований (\texttt{dashboard.html}) и дашборд сценариев (\texttt{usecases\_dashboard.html}).
\end{abstract}
\tableofcontents
\clearpage

\section*{Краткие выводы}
\addcontentsline{toc}{section}{Краткие выводы}
\begin{enumerate}
\item \textbf{Нет инструмента, закрывающего потребности сценариев целиком.} Для каждого из 10 сценариев определены критические требования
(31 различных требований); ни один из 27 инструментов не выполняет полностью критические требования ни одного сценария.
Лучший инструмент для сценария в среднем закрывает 78\,\% его критических требований.
\item \textbf{Разнородные смеси полностью поддерживают 6 средств из 27:}
\texttt{pomegranate}, \texttt{mixem}, \texttt{FlexMix}, \texttt{brms}, \texttt{Distributions.jl}, \texttt{PySATL mpest}; ещё 3 --- частично (\texttt{StepMix}, \texttt{Rmixmod (mixmod)}, \texttt{mixdist}).
Схему ECME среди них реализует только \texttt{PySATL mpest}; блочная схема ECM есть у \texttt{FlexMix}, \texttt{PySATL mpest}.
\item \textbf{12 из 86 требований не выполнено полностью ни одним инструментом}, ещё 17 выполнены не более чем двумя;
базовый слой (EM, нормальное семейство, логарифмическая устойчивость, сборка и тесты) есть почти у всех.
\item \textbf{Положение \texttt{pysatl-mpest}.} Обязательные требования выполнены на 81\,\% (лучший из остальных --- \texttt{FlexMix}, 72\,\%);
по всем требованиям 6-е место из 27. В разнородных сценариях UC-07, UC-08 он лидирует; в однородных гауссовых уступает зрелым пакетам.
\item \textbf{Приоритеты развития} следуют из двух списков: обязательные требования, не выполненные библиотекой
(кластеризация данных; рандомизированный мультистарт; декомпозиция плотности смеси; явные параметрические границы; параметрический бутстреп отношения правдоподобий (LRT)), и расчёт выигрыша по сценариям (раздел~\ref{sec:roadmap}).
\item \textbf{Ограничения:} оценки получены по документации и коду без запуска программ и характеризуют заявленные возможности, а не качество; сценарии --- классы задач, а не полная выборка.
\end{enumerate}

\part{Данные и методика}
\section{Эмпирическая база}
\subsection{Состав}
База состоит из связанных таблиц; ссылочная целостность между ними поддерживается.
\begin{center}\small
\begin{tabular}{@{}l p{9.2cm} r@{}}
\toprule
Таблица базы & Содержание & Записей \\
\midrule
Инструменты & язык, репозиторий, лицензия, публикация, описание & 27 \\
Требования & категория, название, описание, источник, приоритет & 86 \\
Оценки & статус, обоснование и ссылка для каждой пары <<инструмент --- требование>> & 2322 \\
Связи требований с публикациями & роли публикаций: основа, метод, эталонные данные & 576 \\
Публикации & авторы, год, цитирования, ссылки, воспроизводимость & 445 \\
Сценарии & описание и характеристики данных и смеси & 10 \\
Требования сценариев & критические и желательные требования с обоснованием & 214 \\
Публикации и наборы данных сценариев & публикации и наборы данных, на которых основан сценарий & 48 \\
\bottomrule
\end{tabular}
\end{center}
Матрица оценок полна: 27 $\times$ 86 $=$ 2322 строк. Распределение статусов: выполнено --- 624 (27\,\%),
частично --- 289 (12\,\%), не выполнено --- 1409 (61\,\%).
К 2104 из 2322 оценок приложена ссылка на документацию или код.

\subsection{Как собиралась база}
Каталог требований составлен по инструментам: сначала разобраны 27 инструментов и их возможности оформлены как атомарные требования; затем из статей
добавлены специфичные методы, которых в инструментах нет. Приоритеты: обязательные --- базовые архитектурные и математические требования ядра (34),
желательные --- важные конкурентные возможности (32), факультативные --- исследовательские и специализированные методы из статей (20).
Сценарии выбирались по публикациям; затем проверялось, какие требования нужны сценарию и какие инструменты их выполняют.
Из 445 публикаций реестра 325 связаны с требованиями или сценариями (324 --- с требованиями, 45 --- со сценариями).

\subsection{Границы предметной области}
В область работы входят только конечные смеси одномерных непрерывных распределений, однородные и разнородные. Вне области: многомерные смеси,
дискретные смеси и смеси с точечной массой, регрессионные смеси с ковариатами, модели латентных классов, скрытые марковские модели, модели дискретного выбора,
многоуровневые модели. Оценки инструментов выставлены по одним и тем же требованиям независимо от этих границ, поэтому инструменты, рассчитанные прежде всего
на многомерные данные (\texttt{scikit-learn (sklearn.mixture)}, \texttt{mixem}, \texttt{student-mixture}, \texttt{EMCluster}, \texttt{teigen}, \texttt{sn (Skew-Normal/t)}, \texttt{GaussianMixtures.jl}), оцениваются по тому, что они дают в одномерном случае.

\section{Правила подсчёта}
\subsection{Покрытие}
Покрытие группы требований инструментом --- среднее значение весов: выполнено $=1$, частично $=\tfrac12$, не выполнено $=0$. Показатель отражает заявленные возможности по
документации и коду; программы не запускались, поэтому он не измеряет качество оценивания. В разделах про сценарии покрытие считается по критическим требованиям сценария.

\subsection{Чувствительность к весу частичных оценок}
Вес $\tfrac12$ для частично выполненных требований --- условность. В таблице показано, как меняется место инструмента при весах $0$ (частично не засчитывается),
$\tfrac12$ и $1$ (частично приравнивается к выполненному), а также при взвешивании по приоритетам (обязательные $\times 3$, желательные $\times 2$, факультативные $\times 1$).
{\footnotesize
\begin{longtable}{@{}l r r r r r r r r@{}}
\toprule
Инструмент & $w{=}0$ & $w{=}\tfrac12$ & $w{=}1$ & По приоритетам & Обяз. & Жел. & Факульт. & Разброс \\
\midrule
\endhead
brms & 1 & 1 & 1 & 1 & 3 & 1 & 1 & 0 \\
FlexMix & 2 & 2 & 2 & 2 & 2 & 7 & 4 & 0 \\
mclust (v5+) & 7 & 3 & 4 & 6 & 10 & 2 & 5 & 4 \\
bnpy & 9 & 4 & 3 & 11 & 12 & 8 & 2 & 8 \\
mixR & 4 & 5 & 6 & 4 & 6 & 6 & 9 & 2 \\
PySATL mpest & 3 & 6 & 11 & 3 & 1 & 21 & 27 & 8 \\
StepMix & 6 & 7 & 8 & 5 & 4 & 10 & 12 & 3 \\
Rmixmod (mixmod) & 5 & 8 & 9 & 7 & 9 & 5 & 6 & 4 \\
mixtools & 11 & 9 & 5 & 8 & 8 & 12 & 3 & 6 \\
rebmix & 8 & 10 & 13 & 12 & 13 & 3 & 20 & 5 \\
pomegranate & 13 & 11 & 7 & 9 & 7 & 19 & 10 & 6 \\
Distributions.jl & 10 & 12 & 12 & 10 & 5 & 14 & 19 & 2 \\
distfit & 19 & 13 & 10 & 14 & 18 & 4 & 21 & 9 \\
mixdist & 12 & 14 & 14 & 15 & 14 & 11 & 13 & 3 \\
scikit-learn (sklearn.mixture) & 14 & 15 & 15 & 13 & 11 & 20 & 11 & 2 \\
sn (Skew-Normal/t) & 15 & 16 & 18 & 16 & 15 & 9 & 22 & 3 \\
EMCluster & 16 & 17 & 17 & 17 & 16 & 15 & 8 & 1 \\
teigen & 20 & 18 & 16 & 18 & 23 & 16 & 7 & 4 \\
CAMAN & 17 & 19 & 19 & 19 & 17 & 13 & 18 & 2 \\
GaussianMixtures.jl & 21 & 20 & 21 & 20 & 21 & 18 & 23 & 1 \\
fpc & 24 & 21 & 20 & 21 & 24 & 22 & 16 & 4 \\
ClusterR & 18 & 22 & 24 & 23 & 27 & 17 & 15 & 6 \\
mlpack (mlpack::gmm) & 22 & 23 & 22 & 22 & 20 & 25 & 17 & 1 \\
student-mixture & 26 & 24 & 23 & 25 & 22 & 26 & 14 & 3 \\
kasakh Mixture-Models & 23 & 25 & 25 & 24 & 25 & 23 & 25 & 2 \\
Armadillo C++ & 25 & 26 & 26 & 26 & 26 & 24 & 24 & 1 \\
mixem & 27 & 27 & 27 & 27 & 19 & 27 & 26 & 0 \\
\bottomrule
\end{longtable}}
Места определяются по убыванию покрытия; у равных значений порядок произвольный. Наиболее устойчивые лидеры по всем схемам:
\texttt{brms}, \texttt{FlexMix}.
Положение \texttt{pysatl-mpest}: 3-е при $w{=}0$, 6-е при $w{=}\tfrac12$, 11-е при $w{=}1$, 3-е с весами приоритетов;
по обязательным требованиям 1-е, по желательным 21-е, по факультативным 27-е.

\subsection{Ограничения методики}
\begin{itemize}
\item Оценки выставлены по документации и коду; поведение программ не измерялось. Вывод об отсутствии возможности относится к рассмотренным интерфейсам и версиям.
\item Равные веса требований внутри приоритета не отражают разную трудоёмкость реализации.
\item Набор из 10 сценариев иллюстрирует классы задач и не исчерпывает их; сценарии выбирались по публикациям с открытыми данными.
\item Статус \texttt{pysatl-mpest} соответствует состоянию на дату оценки.
\end{itemize}

\part{Инструменты}
\section{Каталог}
Список инструментов в порядке реестра. Статус: <<рассм.>> --- инструмент разобран по документации и коду, <<канд.>> --- кандидат, оценён по документации.
{\footnotesize
\begin{longtable}{@{}>{\raggedright\arraybackslash}p{4.2cm} >{\raggedright\arraybackslash}p{1.8cm} >{\raggedright\arraybackslash}p{2.3cm} l >{\raggedright\arraybackslash}p{6.2cm}@{}}
\toprule
Инструмент & Язык & Лицензия & Статус & Репозиторий \\
\midrule
\endhead
PySATL mpest & Python & MIT & рассм. & \url{https://github.com/PySATL/rework-pysatl-mpest} \\
scikit-learn (sklearn.mixture) & Python & BSD-3-Clause & рассм. & \url{https://github.com/scikit-learn/scikit-learn} \\
pomegranate & Python & MIT & рассм. & \url{https://github.com/jmschrei/pomegranate} \\
mixem & Python & MIT & рассм. & \url{https://github.com/sseemayer/mixem} \\
StepMix & Python & MIT & рассм. & \url{https://github.com/Labo-Lacourse/stepmix} \\
student-mixture & Python & MIT & рассм. & \url{https://github.com/omritomer/student-mixture} \\
kasakh Mixture-Models & Python & MIT & рассм. & \url{https://github.com/kasakh/Mixture-Models} \\
distfit & Python & MIT & рассм. & \url{https://github.com/erdogant/distfit} \\
mclust (v5+) & R & GPL-2 | GPL-3 & рассм. & \url{https://github.com/cran/mclust} \\
mixtools & R & GPL-2 | GPL-3 & рассм. & \url{https://github.com/cran/mixtools} \\
FlexMix & R & GPL-2 | GPL-3 & рассм. & \url{https://github.com/cran/flexmix} \\
mixR & R & GPL-2 | GPL-3 & рассм. & \url{https://github.com/GaryBAYLOR/mixR} \\
rebmix & R / C++ & GPL-2 | GPL-3 & рассм. & \url{https://github.com/cran/rebmix} \\
Rmixmod (mixmod) & R / C++ & GPL-3 & рассм. & \url{https://github.com/mixmod/mixmod} \\
ClusterR & R / C++ & GPL-3 & рассм. & \url{https://github.com/mlampros/ClusterR} \\
mixdist & R & GPL-2 & рассм. & \url{https://github.com/cran/mixdist} \\
EMCluster & R / C & GPL-2 & рассм. & \url{https://github.com/snoweye/EMCluster} \\
teigen & R & GPL-2 | GPL-3 & рассм. & \url{https://github.com/cran/teigen} \\
fpc & R & GPL-2 | GPL-3 & рассм. & \url{https://github.com/cran/fpc} \\
sn (Skew-Normal/t) & R & GPL-2 | GPL-3 & рассм. & \url{https://github.com/cran/sn} \\
brms & R / C++ & GPL-2 & рассм. & \url{https://github.com/paul-buerkner/brms} \\
GaussianMixtures.jl & Julia & MIT & рассм. & \url{https://github.com/davidavdav/GaussianMixtures.jl} \\
Distributions.jl & Julia & MIT & рассм. & \url{https://github.com/JuliaStats/Distributions.jl} \\
mlpack (mlpack::gmm) & C++ & BSD-3-Clause & рассм. & \url{https://github.com/mlpack/mlpack} \\
Armadillo C++ & C++ & Apache-2.0 & рассм. & \url{https://gitlab.com/conradsnicta/armadillo-code} \\
bnpy & Python & BSD-3-Clause & канд. & \url{https://github.com/bnpy/bnpy} \\
CAMAN & R / C & GPL-2 & канд. & \url{https://github.com/cran/CAMAN} \\
\bottomrule
\end{longtable}}

\subsection{Распределение по языкам и лицензиям}
\begin{center}\small
\begin{tabular}{@{}l r l r@{}}
\toprule
Язык & Инструментов & Лицензия & Инструментов \\
\midrule
Python & 9 & MIT & 9 \\
R & 8 & GPL-2 | GPL-3 & 8 \\
R / C++ & 4 & GPL-2 & 4 \\
R / C & 2 & BSD-3-Clause & 3 \\
Julia & 2 & GPL-3 & 2 \\
C++ & 2 & Apache-2.0 & 1 \\
\bottomrule
\end{tabular}
\end{center}

\begin{landscape}
\section{Сводная таблица покрытия}
Инструменты упорядочены по общему покрытию (вес частичных $\tfrac12$). Числа в цветных ячейках --- проценты; чем темнее ячейка, тем выше покрытие.
Обозначения: <<Да>>, <<Част.>>, <<Нет>> --- число требований с соответствующим статусом (из 86).
{\scriptsize\setlength{\tabcolsep}{3pt}
\begin{longtable}{@{}r l l r r r r r r r *{10}{r}@{}}
\toprule
\# & Инструмент & Язык & \rotatebox{90}{\scriptsize Да} & \rotatebox{90}{\scriptsize Част.} & \rotatebox{90}{\scriptsize Нет} & \rotatebox{90}{\scriptsize Все, \%} & \rotatebox{90}{\scriptsize Обяз., \%} & \rotatebox{90}{\scriptsize Жел., \%} & \rotatebox{90}{\scriptsize Факульт., \%} & \rotatebox{90}{\scriptsize Оценивание} & \rotatebox{90}{\scriptsize Структура смеси} & \rotatebox{90}{\scriptsize Инициализация} & \rotatebox{90}{\scriptsize Архитектура} & \rotatebox{90}{\scriptsize Инфраструктура} & \rotatebox{90}{\scriptsize Визуализация} & \rotatebox{90}{\scriptsize Специальные методы} & \rotatebox{90}{\scriptsize Представление смеси} & \rotatebox{90}{\scriptsize Производительность} & \rotatebox{90}{\scriptsize Диагностика} \\
\midrule
\endhead
1 & \textbf{brms} & R / C++ & 41 & 11 & 34 & \cellcolor{hm!43}54 & \cellcolor{hm!51}66 & \cellcolor{hm!38}47 & \cellcolor{hm!36}45 & \cellcolor{hm!30}35 & \cellcolor{hm!38}48 & \cellcolor{hm!40}50 & \cellcolor{hm!75}100 & \cellcolor{hm!58}75 & \cellcolor{hm!61}80 & \cellcolor{hm!33}40 & \cellcolor{hm!49}62 & \cellcolor{hm!66}88 & \cellcolor{hm!49}62 \\
2 & \textbf{FlexMix} & R & 33 & 14 & 39 & \cellcolor{hm!38}47 & \cellcolor{hm!55}72 & \cellcolor{hm!28}33 & \cellcolor{hm!22}25 & \cellcolor{hm!28}33 & \cellcolor{hm!40}50 & \cellcolor{hm!40}50 & \cellcolor{hm!61}80 & \cellcolor{hm!58}75 & \cellcolor{hm!54}70 & \cellcolor{hm!26}30 & \cellcolor{hm!22}25 & \cellcolor{hm!31}38 & \cellcolor{hm!31}38 \\
3 & \textbf{mclust (v5+)} & R & 27 & 17 & 42 & \cellcolor{hm!34}41 & \cellcolor{hm!42}53 & \cellcolor{hm!32}39 & \cellcolor{hm!22}25 & \cellcolor{hm!24}27 & \cellcolor{hm!26}30 & \cellcolor{hm!46}58 & \cellcolor{hm!40}50 & \cellcolor{hm!58}75 & \cellcolor{hm!47}60 & \cellcolor{hm!33}40 & \cellcolor{hm!40}50 & \cellcolor{hm!58}75 & \cellcolor{hm!31}38 \\
4 & \textbf{bnpy} & Python & 26 & 18 & 42 & \cellcolor{hm!33}41 & \cellcolor{hm!40}50 & \cellcolor{hm!26}30 & \cellcolor{hm!35}42 & \cellcolor{hm!28}33 & \cellcolor{hm!28}33 & \cellcolor{hm!52}67 & \cellcolor{hm!40}50 & \cellcolor{hm!28}33 & \cellcolor{hm!40}50 & \cellcolor{hm!19}20 & \cellcolor{hm!40}50 & \cellcolor{hm!75}100 & \cellcolor{hm!31}38 \\
5 & \textbf{mixR} & R & 30 & 10 & 46 & \cellcolor{hm!33}41 & \cellcolor{hm!47}60 & \cellcolor{hm!29}34 & \cellcolor{hm!17}18 & \cellcolor{hm!27}31 & \cellcolor{hm!25}28 & \cellcolor{hm!75}100 & \cellcolor{hm!33}40 & \cellcolor{hm!58}75 & \cellcolor{hm!47}60 & \cellcolor{hm!19}20 & \cellcolor{hm!49}62 & \cellcolor{hm!40}50 & \cellcolor{hm!5}0 \\
6 & \textbf{PySATL mpest} (цель) & Python & 31 & 6 & 49 & \cellcolor{hm!33}40 & \cellcolor{hm!62}81 & \cellcolor{hm!19}20 & \cellcolor{hm!5}0 & \cellcolor{hm!31}38 & \cellcolor{hm!34}41 & \cellcolor{hm!17}17 & \cellcolor{hm!61}80 & \cellcolor{hm!75}100 & \cellcolor{hm!5}0 & \cellcolor{hm!5}0 & \cellcolor{hm!49}62 & \cellcolor{hm!40}50 & \cellcolor{hm!5}0 \\
7 & \textbf{StepMix} & Python & 28 & 11 & 47 & \cellcolor{hm!32}39 & \cellcolor{hm!48}62 & \cellcolor{hm!25}28 & \cellcolor{hm!17}18 & \cellcolor{hm!25}29 & \cellcolor{hm!25}28 & \cellcolor{hm!46}58 & \cellcolor{hm!47}60 & \cellcolor{hm!58}75 & \cellcolor{hm!47}60 & \cellcolor{hm!26}30 & \cellcolor{hm!22}25 & \cellcolor{hm!40}50 & \cellcolor{hm!31}38 \\
8 & \textbf{Rmixmod (mixmod)} & R / C++ & 28 & 10 & 48 & \cellcolor{hm!32}38 & \cellcolor{hm!42}53 & \cellcolor{hm!29}34 & \cellcolor{hm!19}20 & \cellcolor{hm!22}25 & \cellcolor{hm!22}24 & \cellcolor{hm!46}58 & \cellcolor{hm!61}80 & \cellcolor{hm!58}75 & \cellcolor{hm!47}60 & \cellcolor{hm!19}20 & \cellcolor{hm!40}50 & \cellcolor{hm!66}88 & \cellcolor{hm!5}0 \\
9 & \textbf{mixtools} & R & 24 & 18 & 44 & \cellcolor{hm!32}38 & \cellcolor{hm!44}56 & \cellcolor{hm!24}27 & \cellcolor{hm!24}28 & \cellcolor{hm!31}38 & \cellcolor{hm!29}35 & \cellcolor{hm!40}50 & \cellcolor{hm!26}30 & \cellcolor{hm!52}67 & \cellcolor{hm!47}60 & \cellcolor{hm!26}30 & \cellcolor{hm!22}25 & \cellcolor{hm!14}12 & \cellcolor{hm!31}38 \\
10 & \textbf{rebmix} & R / C++ & 27 & 9 & 50 & \cellcolor{hm!31}37 & \cellcolor{hm!40}50 & \cellcolor{hm!32}39 & \cellcolor{hm!12}10 & \cellcolor{hm!21}23 & \cellcolor{hm!38}48 & \cellcolor{hm!5}0 & \cellcolor{hm!47}60 & \cellcolor{hm!58}75 & \cellcolor{hm!26}30 & \cellcolor{hm!5}0 & \cellcolor{hm!40}50 & \cellcolor{hm!58}75 & \cellcolor{hm!22}25 \\
11 & \textbf{pomegranate} & Python & 23 & 16 & 47 & \cellcolor{hm!30}36 & \cellcolor{hm!47}60 & \cellcolor{hm!20}22 & \cellcolor{hm!17}18 & \cellcolor{hm!27}31 & \cellcolor{hm!29}35 & \cellcolor{hm!40}50 & \cellcolor{hm!47}60 & \cellcolor{hm!46}58 & \cellcolor{hm!5}0 & \cellcolor{hm!19}20 & \cellcolor{hm!49}62 & \cellcolor{hm!49}62 & \cellcolor{hm!5}0 \\
12 & \textbf{Distributions.jl} & Julia & 24 & 12 & 50 & \cellcolor{hm!29}35 & \cellcolor{hm!47}60 & \cellcolor{hm!21}23 & \cellcolor{hm!12}10 & \cellcolor{hm!18}19 & \cellcolor{hm!29}35 & \cellcolor{hm!17}17 & \cellcolor{hm!54}70 & \cellcolor{hm!69}92 & \cellcolor{hm!12}10 & \cellcolor{hm!5}0 & \cellcolor{hm!66}88 & \cellcolor{hm!58}75 & \cellcolor{hm!14}12 \\
13 & \textbf{distfit} & Python & 19 & 18 & 49 & \cellcolor{hm!28}33 & \cellcolor{hm!36}44 & \cellcolor{hm!30}36 & \cellcolor{hm!10}8 & \cellcolor{hm!15}15 & \cellcolor{hm!32}39 & \cellcolor{hm!5}0 & \cellcolor{hm!26}30 & \cellcolor{hm!58}75 & \cellcolor{hm!47}60 & \cellcolor{hm!12}10 & \cellcolor{hm!66}88 & \cellcolor{hm!31}38 & \cellcolor{hm!22}25 \\
14 & \textbf{mixdist} & R & 23 & 10 & 53 & \cellcolor{hm!28}33 & \cellcolor{hm!39}49 & \cellcolor{hm!24}27 & \cellcolor{hm!16}15 & \cellcolor{hm!31}38 & \cellcolor{hm!25}28 & \cellcolor{hm!28}33 & \cellcolor{hm!26}30 & \cellcolor{hm!52}67 & \cellcolor{hm!19}20 & \cellcolor{hm!19}20 & \cellcolor{hm!5}0 & \cellcolor{hm!22}25 & \cellcolor{hm!40}50 \\
15 & \textbf{scikit-learn (sklearn.mixture)} & Python & 23 & 10 & 53 & \cellcolor{hm!28}33 & \cellcolor{hm!41}51 & \cellcolor{hm!20}22 & \cellcolor{hm!17}18 & \cellcolor{hm!20}21 & \cellcolor{hm!20}22 & \cellcolor{hm!46}58 & \cellcolor{hm!40}50 & \cellcolor{hm!69}92 & \cellcolor{hm!5}0 & \cellcolor{hm!26}30 & \cellcolor{hm!31}38 & \cellcolor{hm!66}88 & \cellcolor{hm!5}0 \\
16 & \textbf{sn (Skew-Normal/t)} & R & 23 & 8 & 55 & \cellcolor{hm!27}31 & \cellcolor{hm!39}49 & \cellcolor{hm!25}28 & \cellcolor{hm!10}8 & \cellcolor{hm!25}29 & \cellcolor{hm!22}24 & \cellcolor{hm!28}33 & \cellcolor{hm!26}30 & \cellcolor{hm!58}75 & \cellcolor{hm!19}20 & \cellcolor{hm!12}10 & \cellcolor{hm!58}75 & \cellcolor{hm!22}25 & \cellcolor{hm!22}25 \\
17 & \textbf{EMCluster} & R / C & 22 & 9 & 55 & \cellcolor{hm!27}31 & \cellcolor{hm!37}46 & \cellcolor{hm!21}23 & \cellcolor{hm!17}18 & \cellcolor{hm!15}15 & \cellcolor{hm!22}24 & \cellcolor{hm!46}58 & \cellcolor{hm!47}60 & \cellcolor{hm!58}75 & \cellcolor{hm!19}20 & \cellcolor{hm!19}20 & \cellcolor{hm!22}25 & \cellcolor{hm!40}50 & \cellcolor{hm!31}38 \\
18 & \textbf{teigen} & R & 19 & 14 & 53 & \cellcolor{hm!26}30 & \cellcolor{hm!35}43 & \cellcolor{hm!21}23 & \cellcolor{hm!19}20 & \cellcolor{hm!21}23 & \cellcolor{hm!23}26 & \cellcolor{hm!34}42 & \cellcolor{hm!33}40 & \cellcolor{hm!52}67 & \cellcolor{hm!47}60 & \cellcolor{hm!12}10 & \cellcolor{hm!5}0 & \cellcolor{hm!40}50 & \cellcolor{hm!14}12 \\
19 & \textbf{CAMAN} & R / C & 20 & 9 & 57 & \cellcolor{hm!25}28 & \cellcolor{hm!36}44 & \cellcolor{hm!21}23 & \cellcolor{hm!12}10 & \cellcolor{hm!20}21 & \cellcolor{hm!28}33 & \cellcolor{hm!28}33 & \cellcolor{hm!19}20 & \cellcolor{hm!28}33 & \cellcolor{hm!47}60 & \cellcolor{hm!12}10 & \cellcolor{hm!31}38 & \cellcolor{hm!22}25 & \cellcolor{hm!22}25 \\
20 & \textbf{GaussianMixtures.jl} & Julia & 18 & 9 & 59 & \cellcolor{hm!23}26 & \cellcolor{hm!35}43 & \cellcolor{hm!20}22 & \cellcolor{hm!8}5 & \cellcolor{hm!15}15 & \cellcolor{hm!17}17 & \cellcolor{hm!40}50 & \cellcolor{hm!33}40 & \cellcolor{hm!58}75 & \cellcolor{hm!5}0 & \cellcolor{hm!5}0 & \cellcolor{hm!31}38 & \cellcolor{hm!75}100 & \cellcolor{hm!5}0 \\
21 & \textbf{fpc} & R & 16 & 13 & 57 & \cellcolor{hm!23}26 & \cellcolor{hm!34}41 & \cellcolor{hm!18}19 & \cellcolor{hm!14}12 & \cellcolor{hm!17}17 & \cellcolor{hm!25}28 & \cellcolor{hm!28}33 & \cellcolor{hm!40}50 & \cellcolor{hm!52}67 & \cellcolor{hm!19}20 & \cellcolor{hm!12}10 & \cellcolor{hm!5}0 & \cellcolor{hm!14}12 & \cellcolor{hm!31}38 \\
22 & \textbf{ClusterR} & R / C++ & 20 & 4 & 62 & \cellcolor{hm!23}26 & \cellcolor{hm!31}37 & \cellcolor{hm!20}22 & \cellcolor{hm!14}12 & \cellcolor{hm!17}17 & \cellcolor{hm!14}13 & \cellcolor{hm!28}33 & \cellcolor{hm!47}60 & \cellcolor{hm!58}75 & \cellcolor{hm!19}20 & \cellcolor{hm!12}10 & \cellcolor{hm!5}0 & \cellcolor{hm!58}75 & \cellcolor{hm!22}25 \\
23 & \textbf{mlpack (mlpack::gmm)} & C++ & 18 & 7 & 61 & \cellcolor{hm!22}25 & \cellcolor{hm!36}44 & \cellcolor{hm!14}12 & \cellcolor{hm!14}12 & \cellcolor{hm!14}12 & \cellcolor{hm!11}9 & \cellcolor{hm!34}42 & \cellcolor{hm!40}50 & \cellcolor{hm!63}83 & \cellcolor{hm!5}0 & \cellcolor{hm!12}10 & \cellcolor{hm!40}50 & \cellcolor{hm!58}75 & \cellcolor{hm!22}25 \\
24 & \textbf{student-mixture} & Python & 15 & 10 & 61 & \cellcolor{hm!21}23 & \cellcolor{hm!35}43 & \cellcolor{hm!10}8 & \cellcolor{hm!16}15 & \cellcolor{hm!24}27 & \cellcolor{hm!20}22 & \cellcolor{hm!46}58 & \cellcolor{hm!40}50 & \cellcolor{hm!17}17 & \cellcolor{hm!5}0 & \cellcolor{hm!12}10 & \cellcolor{hm!22}25 & \cellcolor{hm!5}0 & \cellcolor{hm!5}0 \\
25 & \textbf{kasakh Mixture-Models} & Python & 17 & 5 & 64 & \cellcolor{hm!21}23 & \cellcolor{hm!33}40 & \cellcolor{hm!18}19 & \cellcolor{hm!5}0 & \cellcolor{hm!18}19 & \cellcolor{hm!20}22 & \cellcolor{hm!40}50 & \cellcolor{hm!33}40 & \cellcolor{hm!34}42 & \cellcolor{hm!5}0 & \cellcolor{hm!5}0 & \cellcolor{hm!22}25 & \cellcolor{hm!31}38 & \cellcolor{hm!5}0 \\
26 & \textbf{Armadillo C++} & C++ & 15 & 6 & 65 & \cellcolor{hm!20}21 & \cellcolor{hm!32}38 & \cellcolor{hm!14}12 & \cellcolor{hm!8}5 & \cellcolor{hm!15}15 & \cellcolor{hm!11}9 & \cellcolor{hm!34}42 & \cellcolor{hm!19}20 & \cellcolor{hm!58}75 & \cellcolor{hm!5}0 & \cellcolor{hm!5}0 & \cellcolor{hm!31}38 & \cellcolor{hm!58}75 & \cellcolor{hm!5}0 \\
27 & \textbf{mixem} & Python & 14 & 5 & 67 & \cellcolor{hm!18}19 & \cellcolor{hm!36}44 & \cellcolor{hm!8}5 & \cellcolor{hm!5}0 & \cellcolor{hm!15}15 & \cellcolor{hm!20}22 & \cellcolor{hm!17}17 & \cellcolor{hm!54}70 & \cellcolor{hm!40}50 & \cellcolor{hm!5}0 & \cellcolor{hm!5}0 & \cellcolor{hm!14}12 & \cellcolor{hm!5}0 & \cellcolor{hm!5}0 \\
\bottomrule
\end{longtable}}
Спектр узкий: наибольшее покрытие достигает 54\,\% (\texttt{brms}),
медиана составляет 33\,\%; у 26 из 27 инструментов покрытие ниже половины.

\end{landscape}
\section{Сходство инструментов по профилю возможностей}
Для каждого инструмента найдены три ближайших по мере Жаккара над множествами полностью выполненных требований. Значение $1$ означает совпадение профилей,
$0$ --- отсутствие общих выполненных требований. Это показывает, какие инструменты заменяют друг друга и какие занимают отдельные ниши.
{\footnotesize
\begin{longtable}{@{}>{\raggedright\arraybackslash}p{3.8cm} >{\raggedright\arraybackslash}p{12.4cm}@{}}
\toprule
Инструмент & Ближайшие (мера Жаккара) \\
\midrule
\endhead
brms & FlexMix (0.51), Rmixmod (mixmod) (0.44), pomegranate (0.42) \\
FlexMix & StepMix (0.61), Rmixmod (mixmod) (0.53), brms (0.51) \\
mclust (v5+) & StepMix (0.72), EMCluster (0.63), Rmixmod (mixmod) (0.62) \\
bnpy & scikit-learn (sklearn.mixture) (0.48), GaussianMixtures.jl (0.47), pomegranate (0.40) \\
mixR & mclust (v5+) (0.54), EMCluster (0.53), Rmixmod (mixmod) (0.53) \\
PySATL mpest & Distributions.jl (0.57), pomegranate (0.46), FlexMix (0.45) \\
StepMix & mclust (v5+) (0.72), EMCluster (0.61), FlexMix (0.61) \\
Rmixmod (mixmod) & EMCluster (0.72), ClusterR (0.66), mclust (v5+) (0.62) \\
mixtools & FlexMix (0.50), mixdist (0.47), CAMAN (0.47) \\
rebmix & EMCluster (0.53), mclust (v5+) (0.50), mixR (0.46) \\
pomegranate & Distributions.jl (0.62), mixem (0.54), Armadillo C++ (0.52) \\
Distributions.jl & pomegranate (0.62), mixem (0.58), PySATL mpest (0.57) \\
distfit & Distributions.jl (0.34), mixR (0.32), mixdist (0.31) \\
mixdist & mixtools (0.47), sn (Skew-Normal/t) (0.44), FlexMix (0.44) \\
scikit-learn (sklearn.mixture) & ClusterR (0.65), Armadillo C++ (0.65), EMCluster (0.61) \\
sn (Skew-Normal/t) & Distributions.jl (0.52), mixdist (0.44), mixR (0.43) \\
EMCluster & Rmixmod (mixmod) (0.72), mclust (v5+) (0.63), ClusterR (0.62) \\
teigen & mclust (v5+) (0.53), EMCluster (0.52), Rmixmod (mixmod) (0.47) \\
CAMAN & mixtools (0.47), FlexMix (0.43), mclust (v5+) (0.42) \\
GaussianMixtures.jl & Armadillo C++ (0.83), mlpack (mlpack::gmm) (0.64), ClusterR (0.58) \\
fpc & Armadillo C++ (0.41), EMCluster (0.41), mclust (v5+) (0.39) \\
ClusterR & Armadillo C++ (0.67), Rmixmod (mixmod) (0.66), scikit-learn (sklearn.mixture) (0.65) \\
mlpack (mlpack::gmm) & Armadillo C++ (0.74), GaussianMixtures.jl (0.64), ClusterR (0.58) \\
student-mixture & scikit-learn (sklearn.mixture) (0.46), StepMix (0.39), EMCluster (0.37) \\
kasakh Mixture-Models & GaussianMixtures.jl (0.40), EMCluster (0.39), Armadillo C++ (0.39) \\
Armadillo C++ & GaussianMixtures.jl (0.83), mlpack (mlpack::gmm) (0.74), ClusterR (0.67) \\
mixem & Distributions.jl (0.58), pomegranate (0.54), FlexMix (0.42) \\
\bottomrule
\end{longtable}}
Наиболее похожая пара: \texttt{GaussianMixtures.jl} и \texttt{Armadillo C++}
(0.83). Ближайший к \texttt{pysatl-mpest} инструмент: \texttt{Distributions.jl} (0.57);
профиль \texttt{pysatl-mpest} не совпадает ни с одним из рассмотренных (близок к Distributions.jl).

\part{Требования}
\section{Общая картина}
\begin{center}\small
\begin{tabular}{@{}l r r r r r r r@{}}
\toprule
Категория & Всего & Обяз. & Жел. & Фак. & \shortstack[r]{\scriptsize Среднее покрытие\\\scriptsize инструментами, \%} & \shortstack[r]{\scriptsize Выполнено\\\scriptsize никем} & \shortstack[r]{\scriptsize Покрытие\\\scriptsize \texttt{mpest}, \%} \\
\midrule
Оценивание & 24 & 9 & 7 & 8 & \cellcolor{hm!22}24 & 5 & \cellcolor{hm!31}38 \\
Структура смеси & 23 & 10 & 11 & 2 & \cellcolor{hm!25}28 & 4 & \cellcolor{hm!34}41 \\
Инициализация & 6 & 3 & 3 & 0 & \cellcolor{hm!36}44 & 0 & \cellcolor{hm!17}17 \\
Архитектура & 5 & 3 & 2 & 0 & \cellcolor{hm!41}51 & 0 & \cellcolor{hm!61}80 \\
Инфраструктура & 6 & 6 & 0 & 0 & \cellcolor{hm!52}67 & 0 & \cellcolor{hm!75}100 \\
Визуализация & 5 & 1 & 2 & 2 & \cellcolor{hm!27}32 & 0 & \cellcolor{hm!5}0 \\
Специальные методы & 5 & 0 & 0 & 5 & \cellcolor{hm!16}15 & 3 & \cellcolor{hm!5}0 \\
Представление смеси & 4 & 2 & 2 & 0 & \cellcolor{hm!32}38 & 0 & \cellcolor{hm!49}62 \\
Производительность & 4 & 0 & 3 & 1 & \cellcolor{hm!42}53 & 0 & \cellcolor{hm!40}50 \\
Диагностика & 4 & 0 & 2 & 2 & \cellcolor{hm!20}21 & 0 & \cellcolor{hm!5}0 \\
\bottomrule
\end{tabular}
\end{center}
Покрытие считается по 26 инструментам без \texttt{pysatl-mpest}.

\subsection{Распределение требований по числу инструментов}
Сколько требований выполняют полностью ровно $k$ инструментов (из 26, без \texttt{pysatl-mpest}).
{\small
\begin{center}
\begin{tabular}{@{}l rrrrrrrrr@{}}
\toprule
Число инструментов & 0 & 1 & 2 & 3--4 & 5--7 & 8--11 & 12--15 & 16--19 & 20+ \\
\midrule
Требований & 15 & 10 & 4 & 16 & 11 & 11 & 5 & 5 & 9 \\
\bottomrule
\end{tabular}
\end{center}}
Распределение сильно скошено: 9 требований выполняют не менее 20 инструментов, а 29 --- не более двух.

\subsection{Что умеют почти все}
Не менее 20 из 26 инструментов полностью выполняют: \texttt{EST-01} Каркас и фасад классического EM-алгоритма; \texttt{EST-05} Стратегия суррогатной Q-функции на M-шаге; \texttt{STRUC-01} Однородные смеси распределений; \texttt{STRUC-03} Нормальное семейство распределений; \texttt{TECH-03} Логарифмическая стабильность вычислений; \texttt{INFRA-01} Декларативное управление зависимостями и сборкой; \texttt{INFRA-05} Автоматизированная непрерывная интеграция; \texttt{INFRA-06} Автоматическая генерация документации с математикой; \texttt{REPR-02} Сэмплирование.
Это базовый слой, который любая новая библиотека обязана воспроизвести, но отличие на нём не строится.

\subsection{Что не умеет никто}
Ни один из 27 инструментов полностью не выполняет 12 требований: \texttt{EST-08} Стратегия метода L-моментов на M-шаге; \texttt{EST-09} Замкнутые аналитические формулы параметров на M-шаге; \texttt{STRUC-14} Жадное инкрементальное обучение порядка смеси; \texttt{STRUC-18} Предсказание порядка смеси через мета-обучение (Meta-Learning); \texttt{STRUC-19} Классификация параметрических семейств с помощью нейросетей; \texttt{STRUC-20} Амортизированный нейросетевой вывод параметров смеси (Deep Sets); \texttt{EST-19} Ускорение сходимости EM-алгоритма (Aitken, SQUAREM, PX-EM); \texttt{EST-21} Оценивание методом минимального расстояния (Hellinger, Wasserstein); \texttt{PROP-02} Скорректированные тесты отношения правдоподобий (VLMR, LMR, MLRT); \texttt{PROP-04} Символьно-аналитическое оценивание смесей высокой точности; \texttt{PROP-05} Автоматический эвристический перебор аналитических форм компонент; \texttt{EST-22} Покомпонентный EM-алгоритм (CW-EM, SAGE).
11 из них реализуются хотя бы частично (для отдельных семейств или как расширение);
1 не представлено даже частично: символьно-аналитическое оценивание смесей высокой точности.

\subsection{Что умеют единицы}
Не более двух инструментов, включая \texttt{pysatl-mpest}, полностью выполняют 17 требований: \texttt{EST-03} Каркас и фасад ECME-алгоритма; \texttt{EST-06} Стратегия наблюдаемого правдоподобия на M-шаге; \texttt{EST-07} Стратегия взвешенного метода моментов на M-шаге; \texttt{EST-11} Алгоритм Rough-Enhanced-Bayes; \texttt{INIT-05} Инициализация методом моментов по сегментам; \texttt{STRUC-08} Оценивание распределений с неизвестным порогом; \texttt{STRUC-13} Сравнение и отбор конфигураций семейств; \texttt{STRUC-16} Объединение перекрывающихся компонент по критерию унимодальности (DIP-тест); \texttt{STRUC-17} Непараметрическое градиентное нахождение числа компонент (NPMLE); \texttt{INFRA-03} Property-based тестирование математических инвариантов; \texttt{REPR-03} Аналитический расчет моментов смеси; \texttt{DIAG-01} Статистические критерии согласия (KS, AD, Chi2); \texttt{EST-18} Оценивание по цензурированным и усеченным данным; \texttt{EST-20} Рекурсивный потоковый EM-алгоритм; \texttt{DIAG-04} Оценка отделимости и перекрытия компонент; \texttt{EST-23} Робастное оценивание по усеченному правдоподобию (TCLUST, C-шаг); \texttt{EST-24} Байесовское сэмплирование методом расширения данных (Gibbs, MCMC).

\begin{landscape}
\section{Матрицы покрытия по категориям}
Для каждой категории требований показаны статусы всех инструментов. Условные обозначения: \colorbox{yes!35}{+} --- выполнено,
\colorbox{part!45}{$\pm$} --- частично, \textcolor{gray}{·} --- не выполнено. Столбцы обозначены идентификаторами требований; расшифровка идёт под таблицей.

\subsection{Оценивание}
{\scriptsize\setlength{\tabcolsep}{2.2pt}
\begin{longtable}{@{}l *{24}{c}@{}}
\toprule
Инструмент & \rotatebox{90}{\scriptsize EST-01} & \rotatebox{90}{\scriptsize EST-02} & \rotatebox{90}{\scriptsize EST-03} & \rotatebox{90}{\scriptsize EST-04} & \rotatebox{90}{\scriptsize EST-05} & \rotatebox{90}{\scriptsize EST-06} & \rotatebox{90}{\scriptsize EST-07} & \rotatebox{90}{\scriptsize EST-08} & \rotatebox{90}{\scriptsize EST-09} & \rotatebox{90}{\scriptsize EST-10} & \rotatebox{90}{\scriptsize EST-11} & \rotatebox{90}{\scriptsize EST-12} & \rotatebox{90}{\scriptsize EST-13} & \rotatebox{90}{\scriptsize EST-14} & \rotatebox{90}{\scriptsize EST-15} & \rotatebox{90}{\scriptsize EST-16} & \rotatebox{90}{\scriptsize EST-17} & \rotatebox{90}{\scriptsize EST-18} & \rotatebox{90}{\scriptsize EST-19} & \rotatebox{90}{\scriptsize EST-20} & \rotatebox{90}{\scriptsize EST-21} & \rotatebox{90}{\scriptsize EST-22} & \rotatebox{90}{\scriptsize EST-23} & \rotatebox{90}{\scriptsize EST-24} \\
\midrule
\endhead
brms & \textcolor{gray}{·} & \textcolor{gray}{·} & \textcolor{gray}{·} & \cellcolor{yes!35}+ & \textcolor{gray}{·} & \textcolor{gray}{·} & \textcolor{gray}{·} & \textcolor{gray}{·} & \textcolor{gray}{·} & \textcolor{gray}{·} & \textcolor{gray}{·} & \cellcolor{yes!35}+ & \cellcolor{yes!35}+ & \cellcolor{yes!35}+ & \cellcolor{yes!35}+ & \cellcolor{part!45}$\pm$ & \cellcolor{yes!35}+ & \cellcolor{yes!35}+ & \textcolor{gray}{·} & \textcolor{gray}{·} & \textcolor{gray}{·} & \textcolor{gray}{·} & \textcolor{gray}{·} & \cellcolor{yes!35}+ \\
FlexMix & \cellcolor{yes!35}+ & \cellcolor{yes!35}+ & \textcolor{gray}{·} & \textcolor{gray}{·} & \cellcolor{yes!35}+ & \textcolor{gray}{·} & \textcolor{gray}{·} & \textcolor{gray}{·} & \cellcolor{part!45}$\pm$ & \cellcolor{yes!35}+ & \textcolor{gray}{·} & \cellcolor{yes!35}+ & \textcolor{gray}{·} & \cellcolor{part!45}$\pm$ & \cellcolor{part!45}$\pm$ & \textcolor{gray}{·} & \cellcolor{yes!35}+ & \cellcolor{part!45}$\pm$ & \textcolor{gray}{·} & \textcolor{gray}{·} & \textcolor{gray}{·} & \textcolor{gray}{·} & \textcolor{gray}{·} & \textcolor{gray}{·} \\
mclust (v5+) & \cellcolor{yes!35}+ & \textcolor{gray}{·} & \textcolor{gray}{·} & \textcolor{gray}{·} & \cellcolor{yes!35}+ & \textcolor{gray}{·} & \textcolor{gray}{·} & \textcolor{gray}{·} & \cellcolor{part!45}$\pm$ & \textcolor{gray}{·} & \textcolor{gray}{·} & \cellcolor{part!45}$\pm$ & \textcolor{gray}{·} & \cellcolor{yes!35}+ & \cellcolor{part!45}$\pm$ & \textcolor{gray}{·} & \cellcolor{yes!35}+ & \textcolor{gray}{·} & \cellcolor{part!45}$\pm$ & \textcolor{gray}{·} & \textcolor{gray}{·} & \textcolor{gray}{·} & \cellcolor{part!45}$\pm$ & \textcolor{gray}{·} \\
bnpy & \cellcolor{yes!35}+ & \textcolor{gray}{·} & \textcolor{gray}{·} & \cellcolor{yes!35}+ & \cellcolor{yes!35}+ & \textcolor{gray}{·} & \textcolor{gray}{·} & \textcolor{gray}{·} & \cellcolor{part!45}$\pm$ & \textcolor{gray}{·} & \textcolor{gray}{·} & \cellcolor{yes!35}+ & \textcolor{gray}{·} & \textcolor{gray}{·} & \textcolor{gray}{·} & \textcolor{gray}{·} & \cellcolor{yes!35}+ & \textcolor{gray}{·} & \cellcolor{part!45}$\pm$ & \cellcolor{yes!35}+ & \textcolor{gray}{·} & \cellcolor{part!45}$\pm$ & \textcolor{gray}{·} & \cellcolor{part!45}$\pm$ \\
mixR & \cellcolor{yes!35}+ & \textcolor{gray}{·} & \textcolor{gray}{·} & \cellcolor{part!45}$\pm$ & \cellcolor{yes!35}+ & \textcolor{gray}{·} & \textcolor{gray}{·} & \textcolor{gray}{·} & \cellcolor{part!45}$\pm$ & \textcolor{gray}{·} & \textcolor{gray}{·} & \textcolor{gray}{·} & \cellcolor{yes!35}+ & \textcolor{gray}{·} & \cellcolor{part!45}$\pm$ & \cellcolor{yes!35}+ & \cellcolor{yes!35}+ & \cellcolor{yes!35}+ & \textcolor{gray}{·} & \textcolor{gray}{·} & \textcolor{gray}{·} & \textcolor{gray}{·} & \textcolor{gray}{·} & \textcolor{gray}{·} \\
PySATL mpest & \cellcolor{part!45}$\pm$ & \cellcolor{yes!35}+ & \cellcolor{yes!35}+ & \textcolor{gray}{·} & \cellcolor{yes!35}+ & \cellcolor{yes!35}+ & \cellcolor{part!45}$\pm$ & \cellcolor{part!45}$\pm$ & \cellcolor{part!45}$\pm$ & \textcolor{gray}{·} & \textcolor{gray}{·} & \cellcolor{yes!35}+ & \textcolor{gray}{·} & \cellcolor{yes!35}+ & \cellcolor{yes!35}+ & \textcolor{gray}{·} & \textcolor{gray}{·} & \textcolor{gray}{·} & \textcolor{gray}{·} & \textcolor{gray}{·} & \textcolor{gray}{·} & \textcolor{gray}{·} & \textcolor{gray}{·} & \textcolor{gray}{·} \\
StepMix & \cellcolor{yes!35}+ & \cellcolor{yes!35}+ & \textcolor{gray}{·} & \textcolor{gray}{·} & \cellcolor{yes!35}+ & \textcolor{gray}{·} & \textcolor{gray}{·} & \textcolor{gray}{·} & \cellcolor{part!45}$\pm$ & \textcolor{gray}{·} & \textcolor{gray}{·} & \cellcolor{yes!35}+ & \textcolor{gray}{·} & \cellcolor{yes!35}+ & \cellcolor{part!45}$\pm$ & \textcolor{gray}{·} & \cellcolor{yes!35}+ & \textcolor{gray}{·} & \textcolor{gray}{·} & \textcolor{gray}{·} & \textcolor{gray}{·} & \textcolor{gray}{·} & \textcolor{gray}{·} & \textcolor{gray}{·} \\
Rmixmod (mixmod) & \cellcolor{yes!35}+ & \textcolor{gray}{·} & \textcolor{gray}{·} & \textcolor{gray}{·} & \cellcolor{yes!35}+ & \textcolor{gray}{·} & \textcolor{gray}{·} & \textcolor{gray}{·} & \cellcolor{part!45}$\pm$ & \cellcolor{yes!35}+ & \textcolor{gray}{·} & \cellcolor{part!45}$\pm$ & \textcolor{gray}{·} & \cellcolor{part!45}$\pm$ & \cellcolor{part!45}$\pm$ & \textcolor{gray}{·} & \cellcolor{yes!35}+ & \textcolor{gray}{·} & \textcolor{gray}{·} & \textcolor{gray}{·} & \textcolor{gray}{·} & \textcolor{gray}{·} & \textcolor{gray}{·} & \textcolor{gray}{·} \\
mixtools & \cellcolor{yes!35}+ & \cellcolor{yes!35}+ & \cellcolor{yes!35}+ & \textcolor{gray}{·} & \cellcolor{yes!35}+ & \textcolor{gray}{·} & \textcolor{gray}{·} & \textcolor{gray}{·} & \cellcolor{part!45}$\pm$ & \cellcolor{yes!35}+ & \textcolor{gray}{·} & \cellcolor{part!45}$\pm$ & \textcolor{gray}{·} & \cellcolor{part!45}$\pm$ & \cellcolor{yes!35}+ & \textcolor{gray}{·} & \cellcolor{yes!35}+ & \textcolor{gray}{·} & \textcolor{gray}{·} & \textcolor{gray}{·} & \textcolor{gray}{·} & \textcolor{gray}{·} & \textcolor{gray}{·} & \cellcolor{part!45}$\pm$ \\
rebmix & \cellcolor{yes!35}+ & \textcolor{gray}{·} & \textcolor{gray}{·} & \textcolor{gray}{·} & \textcolor{gray}{·} & \textcolor{gray}{·} & \textcolor{gray}{·} & \textcolor{gray}{·} & \textcolor{gray}{·} & \textcolor{gray}{·} & \cellcolor{yes!35}+ & \textcolor{gray}{·} & \textcolor{gray}{·} & \cellcolor{yes!35}+ & \cellcolor{part!45}$\pm$ & \cellcolor{yes!35}+ & \cellcolor{yes!35}+ & \textcolor{gray}{·} & \textcolor{gray}{·} & \textcolor{gray}{·} & \textcolor{gray}{·} & \textcolor{gray}{·} & \textcolor{gray}{·} & \textcolor{gray}{·} \\
pomegranate & \cellcolor{yes!35}+ & \textcolor{gray}{·} & \textcolor{gray}{·} & \textcolor{gray}{·} & \cellcolor{yes!35}+ & \textcolor{gray}{·} & \textcolor{gray}{·} & \textcolor{gray}{·} & \cellcolor{part!45}$\pm$ & \cellcolor{part!45}$\pm$ & \textcolor{gray}{·} & \cellcolor{yes!35}+ & \cellcolor{part!45}$\pm$ & \cellcolor{yes!35}+ & \cellcolor{part!45}$\pm$ & \textcolor{gray}{·} & \cellcolor{yes!35}+ & \textcolor{gray}{·} & \textcolor{gray}{·} & \cellcolor{part!45}$\pm$ & \textcolor{gray}{·} & \textcolor{gray}{·} & \textcolor{gray}{·} & \textcolor{gray}{·} \\
Distributions.jl & \cellcolor{yes!35}+ & \textcolor{gray}{·} & \textcolor{gray}{·} & \cellcolor{part!45}$\pm$ & \cellcolor{yes!35}+ & \textcolor{gray}{·} & \textcolor{gray}{·} & \textcolor{gray}{·} & \cellcolor{part!45}$\pm$ & \textcolor{gray}{·} & \textcolor{gray}{·} & \textcolor{gray}{·} & \textcolor{gray}{·} & \textcolor{gray}{·} & \textcolor{gray}{·} & \textcolor{gray}{·} & \cellcolor{yes!35}+ & \cellcolor{part!45}$\pm$ & \textcolor{gray}{·} & \textcolor{gray}{·} & \textcolor{gray}{·} & \textcolor{gray}{·} & \textcolor{gray}{·} & \textcolor{gray}{·} \\
distfit & \textcolor{gray}{·} & \textcolor{gray}{·} & \textcolor{gray}{·} & \cellcolor{yes!35}+ & \textcolor{gray}{·} & \textcolor{gray}{·} & \cellcolor{part!45}$\pm$ & \textcolor{gray}{·} & \textcolor{gray}{·} & \textcolor{gray}{·} & \textcolor{gray}{·} & \cellcolor{part!45}$\pm$ & \cellcolor{yes!35}+ & \textcolor{gray}{·} & \cellcolor{part!45}$\pm$ & \textcolor{gray}{·} & \textcolor{gray}{·} & \textcolor{gray}{·} & \textcolor{gray}{·} & \textcolor{gray}{·} & \textcolor{gray}{·} & \textcolor{gray}{·} & \textcolor{gray}{·} & \textcolor{gray}{·} \\
mixdist & \cellcolor{yes!35}+ & \textcolor{gray}{·} & \textcolor{gray}{·} & \cellcolor{yes!35}+ & \cellcolor{yes!35}+ & \textcolor{gray}{·} & \textcolor{gray}{·} & \textcolor{gray}{·} & \cellcolor{part!45}$\pm$ & \textcolor{gray}{·} & \textcolor{gray}{·} & \cellcolor{yes!35}+ & \cellcolor{yes!35}+ & \textcolor{gray}{·} & \cellcolor{part!45}$\pm$ & \cellcolor{yes!35}+ & \cellcolor{yes!35}+ & \textcolor{gray}{·} & \cellcolor{part!45}$\pm$ & \textcolor{gray}{·} & \cellcolor{part!45}$\pm$ & \textcolor{gray}{·} & \textcolor{gray}{·} & \textcolor{gray}{·} \\
scikit-learn (sklearn.mixture) & \cellcolor{yes!35}+ & \textcolor{gray}{·} & \textcolor{gray}{·} & \textcolor{gray}{·} & \cellcolor{yes!35}+ & \textcolor{gray}{·} & \textcolor{gray}{·} & \textcolor{gray}{·} & \cellcolor{part!45}$\pm$ & \textcolor{gray}{·} & \textcolor{gray}{·} & \textcolor{gray}{·} & \textcolor{gray}{·} & \cellcolor{yes!35}+ & \cellcolor{part!45}$\pm$ & \textcolor{gray}{·} & \cellcolor{yes!35}+ & \textcolor{gray}{·} & \textcolor{gray}{·} & \textcolor{gray}{·} & \textcolor{gray}{·} & \textcolor{gray}{·} & \textcolor{gray}{·} & \textcolor{gray}{·} \\
sn (Skew-Normal/t) & \cellcolor{yes!35}+ & \textcolor{gray}{·} & \textcolor{gray}{·} & \cellcolor{yes!35}+ & \cellcolor{yes!35}+ & \textcolor{gray}{·} & \cellcolor{yes!35}+ & \textcolor{gray}{·} & \cellcolor{part!45}$\pm$ & \textcolor{gray}{·} & \textcolor{gray}{·} & \cellcolor{yes!35}+ & \textcolor{gray}{·} & \textcolor{gray}{·} & \cellcolor{part!45}$\pm$ & \textcolor{gray}{·} & \cellcolor{yes!35}+ & \textcolor{gray}{·} & \textcolor{gray}{·} & \textcolor{gray}{·} & \textcolor{gray}{·} & \textcolor{gray}{·} & \textcolor{gray}{·} & \textcolor{gray}{·} \\
EMCluster & \cellcolor{yes!35}+ & \textcolor{gray}{·} & \textcolor{gray}{·} & \textcolor{gray}{·} & \cellcolor{yes!35}+ & \textcolor{gray}{·} & \textcolor{gray}{·} & \textcolor{gray}{·} & \cellcolor{part!45}$\pm$ & \cellcolor{part!45}$\pm$ & \textcolor{gray}{·} & \textcolor{gray}{·} & \textcolor{gray}{·} & \textcolor{gray}{·} & \textcolor{gray}{·} & \textcolor{gray}{·} & \textcolor{gray}{·} & \textcolor{gray}{·} & \textcolor{gray}{·} & \textcolor{gray}{·} & \textcolor{gray}{·} & \textcolor{gray}{·} & \cellcolor{part!45}$\pm$ & \textcolor{gray}{·} \\
teigen & \cellcolor{part!45}$\pm$ & \cellcolor{yes!35}+ & \textcolor{gray}{·} & \textcolor{gray}{·} & \cellcolor{yes!35}+ & \textcolor{gray}{·} & \textcolor{gray}{·} & \textcolor{gray}{·} & \cellcolor{part!45}$\pm$ & \textcolor{gray}{·} & \textcolor{gray}{·} & \cellcolor{part!45}$\pm$ & \textcolor{gray}{·} & \cellcolor{part!45}$\pm$ & \cellcolor{part!45}$\pm$ & \textcolor{gray}{·} & \textcolor{gray}{·} & \textcolor{gray}{·} & \cellcolor{part!45}$\pm$ & \textcolor{gray}{·} & \textcolor{gray}{·} & \textcolor{gray}{·} & \cellcolor{part!45}$\pm$ & \textcolor{gray}{·} \\
CAMAN & \cellcolor{yes!35}+ & \textcolor{gray}{·} & \textcolor{gray}{·} & \textcolor{gray}{·} & \cellcolor{yes!35}+ & \textcolor{gray}{·} & \textcolor{gray}{·} & \textcolor{gray}{·} & \cellcolor{part!45}$\pm$ & \textcolor{gray}{·} & \textcolor{gray}{·} & \cellcolor{yes!35}+ & \textcolor{gray}{·} & \textcolor{gray}{·} & \textcolor{gray}{·} & \textcolor{gray}{·} & \cellcolor{yes!35}+ & \textcolor{gray}{·} & \cellcolor{part!45}$\pm$ & \textcolor{gray}{·} & \textcolor{gray}{·} & \textcolor{gray}{·} & \textcolor{gray}{·} & \textcolor{gray}{·} \\
GaussianMixtures.jl & \cellcolor{yes!35}+ & \textcolor{gray}{·} & \textcolor{gray}{·} & \textcolor{gray}{·} & \cellcolor{yes!35}+ & \textcolor{gray}{·} & \textcolor{gray}{·} & \textcolor{gray}{·} & \cellcolor{part!45}$\pm$ & \textcolor{gray}{·} & \textcolor{gray}{·} & \textcolor{gray}{·} & \textcolor{gray}{·} & \cellcolor{part!45}$\pm$ & \cellcolor{part!45}$\pm$ & \textcolor{gray}{·} & \textcolor{gray}{·} & \textcolor{gray}{·} & \textcolor{gray}{·} & \textcolor{gray}{·} & \textcolor{gray}{·} & \textcolor{gray}{·} & \textcolor{gray}{·} & \textcolor{gray}{·} \\
fpc & \cellcolor{yes!35}+ & \textcolor{gray}{·} & \textcolor{gray}{·} & \textcolor{gray}{·} & \cellcolor{part!45}$\pm$ & \textcolor{gray}{·} & \textcolor{gray}{·} & \textcolor{gray}{·} & \cellcolor{part!45}$\pm$ & \textcolor{gray}{·} & \textcolor{gray}{·} & \textcolor{gray}{·} & \textcolor{gray}{·} & \textcolor{gray}{·} & \cellcolor{part!45}$\pm$ & \textcolor{gray}{·} & \cellcolor{part!45}$\pm$ & \textcolor{gray}{·} & \textcolor{gray}{·} & \textcolor{gray}{·} & \textcolor{gray}{·} & \textcolor{gray}{·} & \cellcolor{yes!35}+ & \textcolor{gray}{·} \\
ClusterR & \cellcolor{yes!35}+ & \textcolor{gray}{·} & \textcolor{gray}{·} & \textcolor{gray}{·} & \cellcolor{yes!35}+ & \textcolor{gray}{·} & \textcolor{gray}{·} & \textcolor{gray}{·} & \cellcolor{part!45}$\pm$ & \textcolor{gray}{·} & \textcolor{gray}{·} & \textcolor{gray}{·} & \textcolor{gray}{·} & \cellcolor{part!45}$\pm$ & \textcolor{gray}{·} & \textcolor{gray}{·} & \cellcolor{yes!35}+ & \textcolor{gray}{·} & \textcolor{gray}{·} & \textcolor{gray}{·} & \textcolor{gray}{·} & \textcolor{gray}{·} & \textcolor{gray}{·} & \textcolor{gray}{·} \\
mlpack (mlpack::gmm) & \cellcolor{yes!35}+ & \textcolor{gray}{·} & \textcolor{gray}{·} & \textcolor{gray}{·} & \cellcolor{yes!35}+ & \textcolor{gray}{·} & \textcolor{gray}{·} & \textcolor{gray}{·} & \cellcolor{part!45}$\pm$ & \textcolor{gray}{·} & \textcolor{gray}{·} & \textcolor{gray}{·} & \textcolor{gray}{·} & \textcolor{gray}{·} & \cellcolor{part!45}$\pm$ & \textcolor{gray}{·} & \textcolor{gray}{·} & \textcolor{gray}{·} & \textcolor{gray}{·} & \textcolor{gray}{·} & \textcolor{gray}{·} & \textcolor{gray}{·} & \textcolor{gray}{·} & \textcolor{gray}{·} \\
student-mixture & \cellcolor{yes!35}+ & \cellcolor{yes!35}+ & \textcolor{gray}{·} & \textcolor{gray}{·} & \cellcolor{yes!35}+ & \textcolor{gray}{·} & \textcolor{gray}{·} & \textcolor{gray}{·} & \cellcolor{part!45}$\pm$ & \textcolor{gray}{·} & \textcolor{gray}{·} & \cellcolor{part!45}$\pm$ & \cellcolor{part!45}$\pm$ & \cellcolor{yes!35}+ & \cellcolor{part!45}$\pm$ & \textcolor{gray}{·} & \textcolor{gray}{·} & \textcolor{gray}{·} & \textcolor{gray}{·} & \textcolor{gray}{·} & \textcolor{gray}{·} & \textcolor{gray}{·} & \cellcolor{part!45}$\pm$ & \textcolor{gray}{·} \\
kasakh Mixture-Models & \textcolor{gray}{·} & \textcolor{gray}{·} & \textcolor{gray}{·} & \cellcolor{yes!35}+ & \textcolor{gray}{·} & \textcolor{gray}{·} & \textcolor{gray}{·} & \textcolor{gray}{·} & \textcolor{gray}{·} & \textcolor{gray}{·} & \textcolor{gray}{·} & \cellcolor{part!45}$\pm$ & \cellcolor{yes!35}+ & \cellcolor{yes!35}+ & \cellcolor{yes!35}+ & \textcolor{gray}{·} & \textcolor{gray}{·} & \textcolor{gray}{·} & \textcolor{gray}{·} & \textcolor{gray}{·} & \textcolor{gray}{·} & \textcolor{gray}{·} & \textcolor{gray}{·} & \textcolor{gray}{·} \\
Armadillo C++ & \cellcolor{yes!35}+ & \textcolor{gray}{·} & \textcolor{gray}{·} & \textcolor{gray}{·} & \cellcolor{yes!35}+ & \textcolor{gray}{·} & \textcolor{gray}{·} & \textcolor{gray}{·} & \cellcolor{part!45}$\pm$ & \textcolor{gray}{·} & \textcolor{gray}{·} & \textcolor{gray}{·} & \textcolor{gray}{·} & \cellcolor{part!45}$\pm$ & \cellcolor{part!45}$\pm$ & \textcolor{gray}{·} & \textcolor{gray}{·} & \textcolor{gray}{·} & \textcolor{gray}{·} & \textcolor{gray}{·} & \textcolor{gray}{·} & \textcolor{gray}{·} & \textcolor{gray}{·} & \textcolor{gray}{·} \\
mixem & \cellcolor{yes!35}+ & \textcolor{gray}{·} & \textcolor{gray}{·} & \textcolor{gray}{·} & \cellcolor{yes!35}+ & \textcolor{gray}{·} & \textcolor{gray}{·} & \textcolor{gray}{·} & \cellcolor{part!45}$\pm$ & \textcolor{gray}{·} & \textcolor{gray}{·} & \textcolor{gray}{·} & \textcolor{gray}{·} & \cellcolor{part!45}$\pm$ & \cellcolor{part!45}$\pm$ & \textcolor{gray}{·} & \textcolor{gray}{·} & \textcolor{gray}{·} & \textcolor{gray}{·} & \textcolor{gray}{·} & \textcolor{gray}{·} & \textcolor{gray}{·} & \textcolor{gray}{·} & \textcolor{gray}{·} \\
\midrule \textit{Выполнено, шт.} & 22 & 6 & 2 & 6 & 22 & 1 & 1 & 0 & 0 & 3 & 1 & 9 & 5 & 9 & 4 & 3 & 16 & 2 & 0 & 1 & 0 & 0 & 1 & 1 \\
\bottomrule
\end{longtable}}
{\scriptsize\setlength{\tabcolsep}{3pt}
\begin{tabular}{@{}l >{\raggedright\arraybackslash}p{7.4cm} l l >{\raggedright\arraybackslash}p{7.4cm} l@{}}
\texttt{EST-01} & Каркас и фасад классического EM-алгоритма & обязательные & \texttt{EST-13} & Явные параметрические границы & обязательные \\
\texttt{EST-02} & Каркас и фасад ECM-алгоритма & обязательные & \texttt{EST-14} & Предотвращение вырождения масштаба и сингулярностей правдоподобия & обязательные \\
\texttt{EST-03} & Каркас и фасад ECME-алгоритма & обязательные & \texttt{EST-15} & Сохранение истории сходимости и времени работы & желательные \\
\texttt{EST-04} & Direct MLE & желательные & \texttt{EST-16} & Оценивание по сгруппированным и интервальным данным & желательные \\
\texttt{EST-05} & Стратегия суррогатной Q-функции на M-шаге & обязательные & \texttt{EST-17} & Взвешенные наблюдения и веса выборки & желательные \\
\texttt{EST-06} & Стратегия наблюдаемого правдоподобия на M-шаге & обязательные & \texttt{EST-18} & Оценивание по цензурированным и усеченным данным & факультативные \\
\texttt{EST-07} & Стратегия взвешенного метода моментов на M-шаге & желательные & \texttt{EST-19} & Ускорение сходимости EM-алгоритма (Aitken, SQUAREM, PX-EM) & факультативные \\
\texttt{EST-08} & Стратегия метода L-моментов на M-шаге & желательные & \texttt{EST-20} & Рекурсивный потоковый EM-алгоритм & факультативные \\
\texttt{EST-09} & Замкнутые аналитические формулы параметров на M-шаге & обязательные & \texttt{EST-21} & Оценивание методом минимального расстояния (Hellinger, Wasserstein) & факультативные \\
\texttt{EST-10} & Стохастический и классификационный EM-алгоритм & факультативные & \texttt{EST-22} & Покомпонентный EM-алгоритм (CW-EM, SAGE) & желательные \\
\texttt{EST-11} & Алгоритм Rough-Enhanced-Bayes & факультативные & \texttt{EST-23} & Робастное оценивание по усеченному правдоподобию (TCLUST, C-шаг) & факультативные \\
\texttt{EST-12} & Фиксация подмножеств параметров & обязательные & \texttt{EST-24} & Байесовское сэмплирование методом расширения данных (Gibbs, MCMC) & факультативные \\
\end{tabular}}

\subsection{Структура смеси}
{\scriptsize\setlength{\tabcolsep}{2.2pt}
\begin{longtable}{@{}l *{23}{c}@{}}
\toprule
Инструмент & \rotatebox{90}{\scriptsize STRUC-01} & \rotatebox{90}{\scriptsize STRUC-02} & \rotatebox{90}{\scriptsize STRUC-03} & \rotatebox{90}{\scriptsize STRUC-04} & \rotatebox{90}{\scriptsize STRUC-05} & \rotatebox{90}{\scriptsize STRUC-06} & \rotatebox{90}{\scriptsize STRUC-07} & \rotatebox{90}{\scriptsize STRUC-08} & \rotatebox{90}{\scriptsize STRUC-09} & \rotatebox{90}{\scriptsize STRUC-10} & \rotatebox{90}{\scriptsize STRUC-11} & \rotatebox{90}{\scriptsize STRUC-12} & \rotatebox{90}{\scriptsize STRUC-13} & \rotatebox{90}{\scriptsize STRUC-14} & \rotatebox{90}{\scriptsize STRUC-15} & \rotatebox{90}{\scriptsize STRUC-16} & \rotatebox{90}{\scriptsize STRUC-17} & \rotatebox{90}{\scriptsize STRUC-18} & \rotatebox{90}{\scriptsize STRUC-19} & \rotatebox{90}{\scriptsize STRUC-20} & \rotatebox{90}{\scriptsize STRUC-21} & \rotatebox{90}{\scriptsize STRUC-22} & \rotatebox{90}{\scriptsize STRUC-23} \\
\midrule
\endhead
brms & \cellcolor{yes!35}+ & \cellcolor{yes!35}+ & \cellcolor{yes!35}+ & \cellcolor{yes!35}+ & \cellcolor{yes!35}+ & \cellcolor{yes!35}+ & \cellcolor{yes!35}+ & \cellcolor{part!45}$\pm$ & \cellcolor{yes!35}+ & \textcolor{gray}{·} & \textcolor{gray}{·} & \textcolor{gray}{·} & \textcolor{gray}{·} & \textcolor{gray}{·} & \textcolor{gray}{·} & \textcolor{gray}{·} & \textcolor{gray}{·} & \textcolor{gray}{·} & \textcolor{gray}{·} & \textcolor{gray}{·} & \cellcolor{yes!35}+ & \cellcolor{yes!35}+ & \cellcolor{part!45}$\pm$ \\
FlexMix & \cellcolor{yes!35}+ & \cellcolor{yes!35}+ & \cellcolor{yes!35}+ & \cellcolor{yes!35}+ & \cellcolor{yes!35}+ & \cellcolor{part!45}$\pm$ & \cellcolor{yes!35}+ & \textcolor{gray}{·} & \cellcolor{yes!35}+ & \cellcolor{yes!35}+ & \cellcolor{part!45}$\pm$ & \cellcolor{yes!35}+ & \textcolor{gray}{·} & \textcolor{gray}{·} & \textcolor{gray}{·} & \textcolor{gray}{·} & \textcolor{gray}{·} & \textcolor{gray}{·} & \textcolor{gray}{·} & \textcolor{gray}{·} & \cellcolor{yes!35}+ & \cellcolor{part!45}$\pm$ & \textcolor{gray}{·} \\
mclust (v5+) & \cellcolor{yes!35}+ & \textcolor{gray}{·} & \cellcolor{yes!35}+ & \textcolor{gray}{·} & \textcolor{gray}{·} & \textcolor{gray}{·} & \cellcolor{part!45}$\pm$ & \textcolor{gray}{·} & \cellcolor{yes!35}+ & \cellcolor{yes!35}+ & \cellcolor{yes!35}+ & \textcolor{gray}{·} & \textcolor{gray}{·} & \textcolor{gray}{·} & \textcolor{gray}{·} & \cellcolor{part!45}$\pm$ & \textcolor{gray}{·} & \textcolor{gray}{·} & \textcolor{gray}{·} & \textcolor{gray}{·} & \cellcolor{yes!35}+ & \textcolor{gray}{·} & \textcolor{gray}{·} \\
bnpy & \cellcolor{yes!35}+ & \textcolor{gray}{·} & \cellcolor{yes!35}+ & \textcolor{gray}{·} & \textcolor{gray}{·} & \textcolor{gray}{·} & \textcolor{gray}{·} & \textcolor{gray}{·} & \cellcolor{yes!35}+ & \textcolor{gray}{·} & \textcolor{gray}{·} & \cellcolor{yes!35}+ & \textcolor{gray}{·} & \cellcolor{part!45}$\pm$ & \cellcolor{yes!35}+ & \textcolor{gray}{·} & \textcolor{gray}{·} & \textcolor{gray}{·} & \textcolor{gray}{·} & \cellcolor{part!45}$\pm$ & \textcolor{gray}{·} & \cellcolor{part!45}$\pm$ & \cellcolor{yes!35}+ \\
mixR & \cellcolor{yes!35}+ & \textcolor{gray}{·} & \cellcolor{yes!35}+ & \cellcolor{part!45}$\pm$ & \cellcolor{yes!35}+ & \textcolor{gray}{·} & \textcolor{gray}{·} & \textcolor{gray}{·} & \cellcolor{yes!35}+ & \textcolor{gray}{·} & \cellcolor{yes!35}+ & \textcolor{gray}{·} & \textcolor{gray}{·} & \textcolor{gray}{·} & \textcolor{gray}{·} & \textcolor{gray}{·} & \textcolor{gray}{·} & \textcolor{gray}{·} & \textcolor{gray}{·} & \textcolor{gray}{·} & \cellcolor{part!45}$\pm$ & \cellcolor{part!45}$\pm$ & \textcolor{gray}{·} \\
PySATL mpest & \cellcolor{yes!35}+ & \cellcolor{yes!35}+ & \cellcolor{yes!35}+ & \cellcolor{yes!35}+ & \cellcolor{yes!35}+ & \cellcolor{yes!35}+ & \cellcolor{yes!35}+ & \cellcolor{yes!35}+ & \cellcolor{part!45}$\pm$ & \textcolor{gray}{·} & \textcolor{gray}{·} & \cellcolor{yes!35}+ & \textcolor{gray}{·} & \textcolor{gray}{·} & \textcolor{gray}{·} & \textcolor{gray}{·} & \textcolor{gray}{·} & \textcolor{gray}{·} & \textcolor{gray}{·} & \textcolor{gray}{·} & \textcolor{gray}{·} & \textcolor{gray}{·} & \textcolor{gray}{·} \\
StepMix & \cellcolor{yes!35}+ & \cellcolor{part!45}$\pm$ & \cellcolor{yes!35}+ & \textcolor{gray}{·} & \textcolor{gray}{·} & \textcolor{gray}{·} & \textcolor{gray}{·} & \textcolor{gray}{·} & \cellcolor{yes!35}+ & \cellcolor{yes!35}+ & \cellcolor{yes!35}+ & \textcolor{gray}{·} & \textcolor{gray}{·} & \textcolor{gray}{·} & \textcolor{gray}{·} & \textcolor{gray}{·} & \textcolor{gray}{·} & \textcolor{gray}{·} & \textcolor{gray}{·} & \textcolor{gray}{·} & \cellcolor{yes!35}+ & \textcolor{gray}{·} & \textcolor{gray}{·} \\
Rmixmod (mixmod) & \cellcolor{yes!35}+ & \cellcolor{part!45}$\pm$ & \cellcolor{yes!35}+ & \textcolor{gray}{·} & \textcolor{gray}{·} & \textcolor{gray}{·} & \textcolor{gray}{·} & \textcolor{gray}{·} & \cellcolor{yes!35}+ & \cellcolor{yes!35}+ & \textcolor{gray}{·} & \textcolor{gray}{·} & \textcolor{gray}{·} & \textcolor{gray}{·} & \textcolor{gray}{·} & \textcolor{gray}{·} & \textcolor{gray}{·} & \textcolor{gray}{·} & \textcolor{gray}{·} & \textcolor{gray}{·} & \cellcolor{yes!35}+ & \textcolor{gray}{·} & \textcolor{gray}{·} \\
mixtools & \cellcolor{yes!35}+ & \textcolor{gray}{·} & \cellcolor{yes!35}+ & \cellcolor{yes!35}+ & \cellcolor{yes!35}+ & \textcolor{gray}{·} & \cellcolor{part!45}$\pm$ & \cellcolor{part!45}$\pm$ & \cellcolor{part!45}$\pm$ & \cellcolor{part!45}$\pm$ & \cellcolor{yes!35}+ & \textcolor{gray}{·} & \textcolor{gray}{·} & \textcolor{gray}{·} & \textcolor{gray}{·} & \textcolor{gray}{·} & \cellcolor{part!45}$\pm$ & \textcolor{gray}{·} & \textcolor{gray}{·} & \textcolor{gray}{·} & \cellcolor{part!45}$\pm$ & \textcolor{gray}{·} & \textcolor{gray}{·} \\
rebmix & \cellcolor{yes!35}+ & \textcolor{gray}{·} & \cellcolor{yes!35}+ & \cellcolor{yes!35}+ & \cellcolor{yes!35}+ & \textcolor{gray}{·} & \cellcolor{part!45}$\pm$ & \textcolor{gray}{·} & \cellcolor{yes!35}+ & \cellcolor{yes!35}+ & \cellcolor{yes!35}+ & \cellcolor{yes!35}+ & \textcolor{gray}{·} & \cellcolor{part!45}$\pm$ & \cellcolor{yes!35}+ & \cellcolor{part!45}$\pm$ & \cellcolor{part!45}$\pm$ & \textcolor{gray}{·} & \textcolor{gray}{·} & \textcolor{gray}{·} & \textcolor{gray}{·} & \textcolor{gray}{·} & \textcolor{gray}{·} \\
pomegranate & \cellcolor{yes!35}+ & \cellcolor{yes!35}+ & \cellcolor{yes!35}+ & \cellcolor{yes!35}+ & \cellcolor{yes!35}+ & \cellcolor{part!45}$\pm$ & \cellcolor{yes!35}+ & \textcolor{gray}{·} & \textcolor{gray}{·} & \textcolor{gray}{·} & \textcolor{gray}{·} & \cellcolor{part!45}$\pm$ & \textcolor{gray}{·} & \textcolor{gray}{·} & \textcolor{gray}{·} & \textcolor{gray}{·} & \textcolor{gray}{·} & \textcolor{gray}{·} & \textcolor{gray}{·} & \textcolor{gray}{·} & \cellcolor{part!45}$\pm$ & \cellcolor{part!45}$\pm$ & \textcolor{gray}{·} \\
Distributions.jl & \cellcolor{yes!35}+ & \cellcolor{yes!35}+ & \cellcolor{yes!35}+ & \cellcolor{yes!35}+ & \cellcolor{yes!35}+ & \cellcolor{yes!35}+ & \cellcolor{yes!35}+ & \textcolor{gray}{·} & \cellcolor{part!45}$\pm$ & \textcolor{gray}{·} & \textcolor{gray}{·} & \textcolor{gray}{·} & \textcolor{gray}{·} & \textcolor{gray}{·} & \textcolor{gray}{·} & \textcolor{gray}{·} & \textcolor{gray}{·} & \textcolor{gray}{·} & \textcolor{gray}{·} & \textcolor{gray}{·} & \textcolor{gray}{·} & \cellcolor{part!45}$\pm$ & \textcolor{gray}{·} \\
distfit & \textcolor{gray}{·} & \textcolor{gray}{·} & \cellcolor{yes!35}+ & \cellcolor{yes!35}+ & \cellcolor{yes!35}+ & \cellcolor{yes!35}+ & \cellcolor{yes!35}+ & \cellcolor{part!45}$\pm$ & \cellcolor{part!45}$\pm$ & \textcolor{gray}{·} & \textcolor{gray}{·} & \textcolor{gray}{·} & \cellcolor{yes!35}+ & \cellcolor{part!45}$\pm$ & \textcolor{gray}{·} & \cellcolor{part!45}$\pm$ & \textcolor{gray}{·} & \textcolor{gray}{·} & \cellcolor{part!45}$\pm$ & \textcolor{gray}{·} & \textcolor{gray}{·} & \cellcolor{part!45}$\pm$ & \textcolor{gray}{·} \\
mixdist & \cellcolor{yes!35}+ & \cellcolor{part!45}$\pm$ & \cellcolor{yes!35}+ & \cellcolor{yes!35}+ & \cellcolor{yes!35}+ & \textcolor{gray}{·} & \textcolor{gray}{·} & \cellcolor{part!45}$\pm$ & \cellcolor{part!45}$\pm$ & \textcolor{gray}{·} & \textcolor{gray}{·} & \textcolor{gray}{·} & \textcolor{gray}{·} & \textcolor{gray}{·} & \textcolor{gray}{·} & \textcolor{gray}{·} & \textcolor{gray}{·} & \textcolor{gray}{·} & \textcolor{gray}{·} & \textcolor{gray}{·} & \cellcolor{yes!35}+ & \textcolor{gray}{·} & \textcolor{gray}{·} \\
scikit-learn (sklearn.mixture) & \cellcolor{yes!35}+ & \textcolor{gray}{·} & \cellcolor{yes!35}+ & \textcolor{gray}{·} & \textcolor{gray}{·} & \textcolor{gray}{·} & \textcolor{gray}{·} & \textcolor{gray}{·} & \cellcolor{yes!35}+ & \textcolor{gray}{·} & \textcolor{gray}{·} & \cellcolor{part!45}$\pm$ & \textcolor{gray}{·} & \textcolor{gray}{·} & \textcolor{gray}{·} & \textcolor{gray}{·} & \textcolor{gray}{·} & \textcolor{gray}{·} & \textcolor{gray}{·} & \textcolor{gray}{·} & \cellcolor{part!45}$\pm$ & \textcolor{gray}{·} & \cellcolor{yes!35}+ \\
sn (Skew-Normal/t) & \cellcolor{yes!35}+ & \textcolor{gray}{·} & \cellcolor{yes!35}+ & \textcolor{gray}{·} & \cellcolor{part!45}$\pm$ & \cellcolor{yes!35}+ & \textcolor{gray}{·} & \textcolor{gray}{·} & \cellcolor{part!45}$\pm$ & \textcolor{gray}{·} & \textcolor{gray}{·} & \textcolor{gray}{·} & \textcolor{gray}{·} & \textcolor{gray}{·} & \textcolor{gray}{·} & \textcolor{gray}{·} & \textcolor{gray}{·} & \textcolor{gray}{·} & \textcolor{gray}{·} & \textcolor{gray}{·} & \cellcolor{part!45}$\pm$ & \cellcolor{yes!35}+ & \textcolor{gray}{·} \\
EMCluster & \cellcolor{yes!35}+ & \textcolor{gray}{·} & \cellcolor{yes!35}+ & \textcolor{gray}{·} & \textcolor{gray}{·} & \textcolor{gray}{·} & \cellcolor{part!45}$\pm$ & \textcolor{gray}{·} & \cellcolor{yes!35}+ & \cellcolor{yes!35}+ & \textcolor{gray}{·} & \textcolor{gray}{·} & \textcolor{gray}{·} & \textcolor{gray}{·} & \textcolor{gray}{·} & \textcolor{gray}{·} & \textcolor{gray}{·} & \textcolor{gray}{·} & \textcolor{gray}{·} & \textcolor{gray}{·} & \cellcolor{yes!35}+ & \textcolor{gray}{·} & \textcolor{gray}{·} \\
teigen & \cellcolor{yes!35}+ & \textcolor{gray}{·} & \textcolor{gray}{·} & \textcolor{gray}{·} & \textcolor{gray}{·} & \cellcolor{yes!35}+ & \textcolor{gray}{·} & \textcolor{gray}{·} & \cellcolor{yes!35}+ & \cellcolor{yes!35}+ & \textcolor{gray}{·} & \textcolor{gray}{·} & \textcolor{gray}{·} & \textcolor{gray}{·} & \textcolor{gray}{·} & \textcolor{gray}{·} & \textcolor{gray}{·} & \textcolor{gray}{·} & \textcolor{gray}{·} & \textcolor{gray}{·} & \cellcolor{yes!35}+ & \cellcolor{yes!35}+ & \textcolor{gray}{·} \\
CAMAN & \cellcolor{yes!35}+ & \textcolor{gray}{·} & \cellcolor{yes!35}+ & \cellcolor{yes!35}+ & \textcolor{gray}{·} & \textcolor{gray}{·} & \textcolor{gray}{·} & \textcolor{gray}{·} & \cellcolor{yes!35}+ & \textcolor{gray}{·} & \cellcolor{yes!35}+ & \cellcolor{yes!35}+ & \textcolor{gray}{·} & \cellcolor{part!45}$\pm$ & \textcolor{gray}{·} & \textcolor{gray}{·} & \cellcolor{yes!35}+ & \textcolor{gray}{·} & \textcolor{gray}{·} & \textcolor{gray}{·} & \textcolor{gray}{·} & \textcolor{gray}{·} & \textcolor{gray}{·} \\
GaussianMixtures.jl & \cellcolor{yes!35}+ & \textcolor{gray}{·} & \cellcolor{yes!35}+ & \textcolor{gray}{·} & \textcolor{gray}{·} & \textcolor{gray}{·} & \textcolor{gray}{·} & \textcolor{gray}{·} & \cellcolor{part!45}$\pm$ & \textcolor{gray}{·} & \textcolor{gray}{·} & \textcolor{gray}{·} & \textcolor{gray}{·} & \cellcolor{part!45}$\pm$ & \cellcolor{yes!35}+ & \textcolor{gray}{·} & \textcolor{gray}{·} & \textcolor{gray}{·} & \textcolor{gray}{·} & \textcolor{gray}{·} & \textcolor{gray}{·} & \textcolor{gray}{·} & \textcolor{gray}{·} \\
fpc & \cellcolor{yes!35}+ & \textcolor{gray}{·} & \cellcolor{yes!35}+ & \textcolor{gray}{·} & \textcolor{gray}{·} & \textcolor{gray}{·} & \textcolor{gray}{·} & \textcolor{gray}{·} & \cellcolor{part!45}$\pm$ & \textcolor{gray}{·} & \cellcolor{yes!35}+ & \cellcolor{yes!35}+ & \textcolor{gray}{·} & \textcolor{gray}{·} & \cellcolor{part!45}$\pm$ & \cellcolor{yes!35}+ & \textcolor{gray}{·} & \cellcolor{part!45}$\pm$ & \textcolor{gray}{·} & \textcolor{gray}{·} & \textcolor{gray}{·} & \textcolor{gray}{·} & \textcolor{gray}{·} \\
ClusterR & \cellcolor{yes!35}+ & \textcolor{gray}{·} & \cellcolor{yes!35}+ & \textcolor{gray}{·} & \textcolor{gray}{·} & \textcolor{gray}{·} & \textcolor{gray}{·} & \textcolor{gray}{·} & \cellcolor{yes!35}+ & \textcolor{gray}{·} & \textcolor{gray}{·} & \textcolor{gray}{·} & \textcolor{gray}{·} & \textcolor{gray}{·} & \textcolor{gray}{·} & \textcolor{gray}{·} & \textcolor{gray}{·} & \textcolor{gray}{·} & \textcolor{gray}{·} & \textcolor{gray}{·} & \textcolor{gray}{·} & \textcolor{gray}{·} & \textcolor{gray}{·} \\
mlpack (mlpack::gmm) & \cellcolor{yes!35}+ & \textcolor{gray}{·} & \cellcolor{yes!35}+ & \textcolor{gray}{·} & \textcolor{gray}{·} & \textcolor{gray}{·} & \textcolor{gray}{·} & \textcolor{gray}{·} & \textcolor{gray}{·} & \textcolor{gray}{·} & \textcolor{gray}{·} & \textcolor{gray}{·} & \textcolor{gray}{·} & \textcolor{gray}{·} & \textcolor{gray}{·} & \textcolor{gray}{·} & \textcolor{gray}{·} & \textcolor{gray}{·} & \textcolor{gray}{·} & \textcolor{gray}{·} & \textcolor{gray}{·} & \textcolor{gray}{·} & \textcolor{gray}{·} \\
student-mixture & \cellcolor{yes!35}+ & \textcolor{gray}{·} & \textcolor{gray}{·} & \textcolor{gray}{·} & \textcolor{gray}{·} & \cellcolor{yes!35}+ & \textcolor{gray}{·} & \textcolor{gray}{·} & \cellcolor{yes!35}+ & \textcolor{gray}{·} & \textcolor{gray}{·} & \textcolor{gray}{·} & \textcolor{gray}{·} & \textcolor{gray}{·} & \textcolor{gray}{·} & \textcolor{gray}{·} & \textcolor{gray}{·} & \textcolor{gray}{·} & \textcolor{gray}{·} & \textcolor{gray}{·} & \textcolor{gray}{·} & \cellcolor{yes!35}+ & \cellcolor{yes!35}+ \\
kasakh Mixture-Models & \cellcolor{yes!35}+ & \textcolor{gray}{·} & \cellcolor{yes!35}+ & \textcolor{gray}{·} & \textcolor{gray}{·} & \cellcolor{yes!35}+ & \textcolor{gray}{·} & \textcolor{gray}{·} & \cellcolor{yes!35}+ & \cellcolor{yes!35}+ & \textcolor{gray}{·} & \textcolor{gray}{·} & \textcolor{gray}{·} & \textcolor{gray}{·} & \textcolor{gray}{·} & \textcolor{gray}{·} & \textcolor{gray}{·} & \textcolor{gray}{·} & \textcolor{gray}{·} & \textcolor{gray}{·} & \textcolor{gray}{·} & \textcolor{gray}{·} & \textcolor{gray}{·} \\
Armadillo C++ & \cellcolor{yes!35}+ & \textcolor{gray}{·} & \cellcolor{yes!35}+ & \textcolor{gray}{·} & \textcolor{gray}{·} & \textcolor{gray}{·} & \textcolor{gray}{·} & \textcolor{gray}{·} & \textcolor{gray}{·} & \textcolor{gray}{·} & \textcolor{gray}{·} & \textcolor{gray}{·} & \textcolor{gray}{·} & \textcolor{gray}{·} & \textcolor{gray}{·} & \textcolor{gray}{·} & \textcolor{gray}{·} & \textcolor{gray}{·} & \textcolor{gray}{·} & \textcolor{gray}{·} & \textcolor{gray}{·} & \textcolor{gray}{·} & \textcolor{gray}{·} \\
mixem & \cellcolor{yes!35}+ & \cellcolor{yes!35}+ & \cellcolor{yes!35}+ & \cellcolor{yes!35}+ & \cellcolor{yes!35}+ & \textcolor{gray}{·} & \textcolor{gray}{·} & \textcolor{gray}{·} & \textcolor{gray}{·} & \textcolor{gray}{·} & \textcolor{gray}{·} & \textcolor{gray}{·} & \textcolor{gray}{·} & \textcolor{gray}{·} & \textcolor{gray}{·} & \textcolor{gray}{·} & \textcolor{gray}{·} & \textcolor{gray}{·} & \textcolor{gray}{·} & \textcolor{gray}{·} & \textcolor{gray}{·} & \textcolor{gray}{·} & \textcolor{gray}{·} \\
\midrule \textit{Выполнено, шт.} & 26 & 6 & 25 & 11 & 11 & 8 & 6 & 1 & 15 & 8 & 7 & 6 & 1 & 0 & 3 & 1 & 1 & 0 & 0 & 0 & 8 & 4 & 3 \\
\bottomrule
\end{longtable}}
{\scriptsize\setlength{\tabcolsep}{3pt}
\begin{tabular}{@{}l >{\raggedright\arraybackslash}p{7.4cm} l l >{\raggedright\arraybackslash}p{7.4cm} l@{}}
\texttt{STRUC-01} & Однородные смеси распределений & обязательные & \texttt{STRUC-13} & Сравнение и отбор конфигураций семейств & желательные \\
\texttt{STRUC-02} & Разнородные смеси распределений & обязательные & \texttt{STRUC-14} & Жадное инкрементальное обучение порядка смеси & желательные \\
\texttt{STRUC-03} & Нормальное семейство распределений & обязательные & \texttt{STRUC-15} & Динамические переходы слияния и расщепления (SMEM) & желательные \\
\texttt{STRUC-04} & Экспоненциальное семейство распределений & обязательные & \texttt{STRUC-16} & Объединение перекрывающихся компонент по критерию унимодальности (DIP-тест) & желательные \\
\texttt{STRUC-05} & Асимметричные положительные семейства (Weibull, Gamma, Lognormal) & обязательные & \texttt{STRUC-17} & Непараметрическое градиентное нахождение числа компонент (NPMLE) & желательные \\
\texttt{STRUC-06} & Тяжелохвостые семейства (Pareto, Cauchy) & обязательные & \texttt{STRUC-18} & Предсказание порядка смеси через мета-обучение (Meta-Learning) & желательные \\
\texttt{STRUC-07} & Семейства на ограниченном носителе (Beta, Uniform) & желательные & \texttt{STRUC-19} & Классификация параметрических семейств с помощью нейросетей & желательные \\
\texttt{STRUC-08} & Оценивание распределений с неизвестным порогом & обязательные & \texttt{STRUC-20} & Амортизированный нейросетевой вывод параметров смеси (Deep Sets) & желательные \\
\texttt{STRUC-09} & Информационные критерии выбора числа компонент (AIC, BIC, CAIC) & обязательные & \texttt{STRUC-21} & Равенство и линейные ограничения параметров компонент & желательные \\
\texttt{STRUC-10} & Критерий интегрированного правдоподобия классификации (ICL) & желательные & \texttt{STRUC-22} & Асимметричные и тяжелохвостые семейства (Skew-Normal, Student-t) & факультативные \\
\texttt{STRUC-11} & Параметрический бутстреп отношения правдоподобий (LRT) & обязательные & \texttt{STRUC-23} & Вариационный байесовский вывод (VBEM) & факультативные \\
\texttt{STRUC-12} & Динамическое удаление вырожденных компонент & обязательные &  & &  \\
\end{tabular}}

\subsection{Инициализация}
{\scriptsize\setlength{\tabcolsep}{2.2pt}
\begin{longtable}{@{}l *{6}{c}@{}}
\toprule
Инструмент & \rotatebox{90}{\scriptsize INIT-01} & \rotatebox{90}{\scriptsize INIT-02} & \rotatebox{90}{\scriptsize INIT-03} & \rotatebox{90}{\scriptsize INIT-04} & \rotatebox{90}{\scriptsize INIT-05} & \rotatebox{90}{\scriptsize INIT-06} \\
\midrule
\endhead
brms & \cellcolor{yes!35}+ & \cellcolor{yes!35}+ & \textcolor{gray}{·} & \textcolor{gray}{·} & \textcolor{gray}{·} & \cellcolor{yes!35}+ \\
FlexMix & \cellcolor{yes!35}+ & \cellcolor{yes!35}+ & \textcolor{gray}{·} & \textcolor{gray}{·} & \textcolor{gray}{·} & \cellcolor{yes!35}+ \\
mclust (v5+) & \cellcolor{part!45}$\pm$ & \textcolor{gray}{·} & \cellcolor{yes!35}+ & \cellcolor{yes!35}+ & \textcolor{gray}{·} & \cellcolor{yes!35}+ \\
bnpy & \cellcolor{yes!35}+ & \cellcolor{yes!35}+ & \cellcolor{yes!35}+ & \textcolor{gray}{·} & \textcolor{gray}{·} & \cellcolor{yes!35}+ \\
mixR & \cellcolor{yes!35}+ & \cellcolor{yes!35}+ & \cellcolor{yes!35}+ & \cellcolor{yes!35}+ & \cellcolor{yes!35}+ & \cellcolor{yes!35}+ \\
PySATL mpest & \cellcolor{yes!35}+ & \textcolor{gray}{·} & \textcolor{gray}{·} & \textcolor{gray}{·} & \textcolor{gray}{·} & \textcolor{gray}{·} \\
StepMix & \cellcolor{part!45}$\pm$ & \cellcolor{yes!35}+ & \cellcolor{yes!35}+ & \textcolor{gray}{·} & \textcolor{gray}{·} & \cellcolor{yes!35}+ \\
Rmixmod (mixmod) & \cellcolor{part!45}$\pm$ & \cellcolor{yes!35}+ & \cellcolor{yes!35}+ & \textcolor{gray}{·} & \textcolor{gray}{·} & \cellcolor{yes!35}+ \\
mixtools & \cellcolor{yes!35}+ & \textcolor{gray}{·} & \textcolor{gray}{·} & \cellcolor{yes!35}+ & \textcolor{gray}{·} & \cellcolor{yes!35}+ \\
rebmix & \textcolor{gray}{·} & \textcolor{gray}{·} & \textcolor{gray}{·} & \textcolor{gray}{·} & \textcolor{gray}{·} & \textcolor{gray}{·} \\
pomegranate & \cellcolor{yes!35}+ & \textcolor{gray}{·} & \cellcolor{yes!35}+ & \textcolor{gray}{·} & \textcolor{gray}{·} & \cellcolor{yes!35}+ \\
Distributions.jl & \cellcolor{yes!35}+ & \textcolor{gray}{·} & \textcolor{gray}{·} & \textcolor{gray}{·} & \textcolor{gray}{·} & \textcolor{gray}{·} \\
distfit & \textcolor{gray}{·} & \textcolor{gray}{·} & \textcolor{gray}{·} & \textcolor{gray}{·} & \textcolor{gray}{·} & \textcolor{gray}{·} \\
mixdist & \cellcolor{yes!35}+ & \textcolor{gray}{·} & \textcolor{gray}{·} & \cellcolor{yes!35}+ & \textcolor{gray}{·} & \textcolor{gray}{·} \\
scikit-learn (sklearn.mixture) & \cellcolor{part!45}$\pm$ & \cellcolor{yes!35}+ & \cellcolor{yes!35}+ & \textcolor{gray}{·} & \textcolor{gray}{·} & \cellcolor{yes!35}+ \\
sn (Skew-Normal/t) & \cellcolor{yes!35}+ & \textcolor{gray}{·} & \textcolor{gray}{·} & \textcolor{gray}{·} & \cellcolor{yes!35}+ & \textcolor{gray}{·} \\
EMCluster & \cellcolor{part!45}$\pm$ & \cellcolor{yes!35}+ & \cellcolor{yes!35}+ & \textcolor{gray}{·} & \textcolor{gray}{·} & \cellcolor{yes!35}+ \\
teigen & \cellcolor{part!45}$\pm$ & \textcolor{gray}{·} & \cellcolor{yes!35}+ & \textcolor{gray}{·} & \textcolor{gray}{·} & \cellcolor{yes!35}+ \\
CAMAN & \cellcolor{yes!35}+ & \textcolor{gray}{·} & \textcolor{gray}{·} & \textcolor{gray}{·} & \textcolor{gray}{·} & \cellcolor{yes!35}+ \\
GaussianMixtures.jl & \cellcolor{yes!35}+ & \textcolor{gray}{·} & \cellcolor{yes!35}+ & \textcolor{gray}{·} & \textcolor{gray}{·} & \cellcolor{yes!35}+ \\
fpc & \textcolor{gray}{·} & \textcolor{gray}{·} & \cellcolor{yes!35}+ & \textcolor{gray}{·} & \textcolor{gray}{·} & \cellcolor{yes!35}+ \\
ClusterR & \textcolor{gray}{·} & \textcolor{gray}{·} & \cellcolor{yes!35}+ & \textcolor{gray}{·} & \textcolor{gray}{·} & \cellcolor{yes!35}+ \\
mlpack (mlpack::gmm) & \cellcolor{part!45}$\pm$ & \cellcolor{part!45}$\pm$ & \cellcolor{yes!35}+ & \textcolor{gray}{·} & \textcolor{gray}{·} & \cellcolor{part!45}$\pm$ \\
student-mixture & \cellcolor{part!45}$\pm$ & \cellcolor{yes!35}+ & \cellcolor{yes!35}+ & \textcolor{gray}{·} & \textcolor{gray}{·} & \cellcolor{yes!35}+ \\
kasakh Mixture-Models & \cellcolor{yes!35}+ & \textcolor{gray}{·} & \cellcolor{yes!35}+ & \textcolor{gray}{·} & \textcolor{gray}{·} & \cellcolor{yes!35}+ \\
Armadillo C++ & \cellcolor{part!45}$\pm$ & \textcolor{gray}{·} & \cellcolor{yes!35}+ & \textcolor{gray}{·} & \textcolor{gray}{·} & \cellcolor{yes!35}+ \\
mixem & \cellcolor{yes!35}+ & \textcolor{gray}{·} & \textcolor{gray}{·} & \textcolor{gray}{·} & \textcolor{gray}{·} & \textcolor{gray}{·} \\
\midrule \textit{Выполнено, шт.} & 14 & 9 & 16 & 4 & 2 & 19 \\
\bottomrule
\end{longtable}}
{\scriptsize\setlength{\tabcolsep}{3pt}
\begin{tabular}{@{}l >{\raggedright\arraybackslash}p{7.4cm} l l >{\raggedright\arraybackslash}p{7.4cm} l@{}}
\texttt{INIT-01} & Пользовательская явная инициализация & обязательные & \texttt{INIT-04} & Квантильное детерминированное разбиение & желательные \\
\texttt{INIT-02} & Рандомизированный мультистарт & обязательные & \texttt{INIT-05} & Инициализация методом моментов по сегментам & желательные \\
\texttt{INIT-03} & Кластеризация данных & обязательные & \texttt{INIT-06} & Случайный выбор наблюдений & желательные \\
\end{tabular}}

\subsection{Архитектура}
{\scriptsize\setlength{\tabcolsep}{2.2pt}
\begin{longtable}{@{}l *{5}{c}@{}}
\toprule
Инструмент & \rotatebox{90}{\scriptsize TECH-01} & \rotatebox{90}{\scriptsize TECH-02} & \rotatebox{90}{\scriptsize TECH-03} & \rotatebox{90}{\scriptsize TECH-04} & \rotatebox{90}{\scriptsize TECH-05} \\
\midrule
\endhead
brms & \cellcolor{yes!35}+ & \cellcolor{yes!35}+ & \cellcolor{yes!35}+ & \cellcolor{yes!35}+ & \cellcolor{yes!35}+ \\
FlexMix & \cellcolor{yes!35}+ & \cellcolor{yes!35}+ & \cellcolor{yes!35}+ & \cellcolor{yes!35}+ & \textcolor{gray}{·} \\
mclust (v5+) & \cellcolor{part!45}$\pm$ & \textcolor{gray}{·} & \cellcolor{yes!35}+ & \cellcolor{yes!35}+ & \textcolor{gray}{·} \\
bnpy & \cellcolor{part!45}$\pm$ & \cellcolor{yes!35}+ & \cellcolor{yes!35}+ & \textcolor{gray}{·} & \textcolor{gray}{·} \\
mixR & \textcolor{gray}{·} & \textcolor{gray}{·} & \cellcolor{yes!35}+ & \cellcolor{yes!35}+ & \textcolor{gray}{·} \\
PySATL mpest & \cellcolor{yes!35}+ & \cellcolor{yes!35}+ & \cellcolor{yes!35}+ & \cellcolor{yes!35}+ & \textcolor{gray}{·} \\
StepMix & \cellcolor{part!45}$\pm$ & \cellcolor{part!45}$\pm$ & \cellcolor{yes!35}+ & \textcolor{gray}{·} & \cellcolor{yes!35}+ \\
Rmixmod (mixmod) & \cellcolor{yes!35}+ & \textcolor{gray}{·} & \cellcolor{yes!35}+ & \cellcolor{yes!35}+ & \cellcolor{yes!35}+ \\
mixtools & \textcolor{gray}{·} & \textcolor{gray}{·} & \cellcolor{part!45}$\pm$ & \cellcolor{yes!35}+ & \textcolor{gray}{·} \\
rebmix & \textcolor{gray}{·} & \textcolor{gray}{·} & \cellcolor{yes!35}+ & \cellcolor{yes!35}+ & \cellcolor{yes!35}+ \\
pomegranate & \cellcolor{part!45}$\pm$ & \cellcolor{yes!35}+ & \cellcolor{yes!35}+ & \textcolor{gray}{·} & \cellcolor{part!45}$\pm$ \\
Distributions.jl & \cellcolor{part!45}$\pm$ & \cellcolor{yes!35}+ & \cellcolor{yes!35}+ & \cellcolor{yes!35}+ & \textcolor{gray}{·} \\
distfit & \textcolor{gray}{·} & \cellcolor{part!45}$\pm$ & \cellcolor{part!45}$\pm$ & \textcolor{gray}{·} & \cellcolor{part!45}$\pm$ \\
mixdist & \textcolor{gray}{·} & \textcolor{gray}{·} & \cellcolor{part!45}$\pm$ & \cellcolor{yes!35}+ & \textcolor{gray}{·} \\
scikit-learn (sklearn.mixture) & \cellcolor{part!45}$\pm$ & \textcolor{gray}{·} & \cellcolor{yes!35}+ & \textcolor{gray}{·} & \cellcolor{yes!35}+ \\
sn (Skew-Normal/t) & \textcolor{gray}{·} & \textcolor{gray}{·} & \cellcolor{part!45}$\pm$ & \cellcolor{yes!35}+ & \textcolor{gray}{·} \\
EMCluster & \textcolor{gray}{·} & \textcolor{gray}{·} & \cellcolor{yes!35}+ & \cellcolor{yes!35}+ & \cellcolor{yes!35}+ \\
teigen & \textcolor{gray}{·} & \textcolor{gray}{·} & \cellcolor{yes!35}+ & \cellcolor{yes!35}+ & \textcolor{gray}{·} \\
CAMAN & \textcolor{gray}{·} & \textcolor{gray}{·} & \cellcolor{part!45}$\pm$ & \textcolor{gray}{·} & \cellcolor{part!45}$\pm$ \\
GaussianMixtures.jl & \cellcolor{part!45}$\pm$ & \textcolor{gray}{·} & \cellcolor{yes!35}+ & \cellcolor{part!45}$\pm$ & \textcolor{gray}{·} \\
fpc & \cellcolor{part!45}$\pm$ & \textcolor{gray}{·} & \cellcolor{yes!35}+ & \cellcolor{yes!35}+ & \textcolor{gray}{·} \\
ClusterR & \textcolor{gray}{·} & \textcolor{gray}{·} & \cellcolor{yes!35}+ & \cellcolor{yes!35}+ & \cellcolor{yes!35}+ \\
mlpack (mlpack::gmm) & \cellcolor{yes!35}+ & \cellcolor{part!45}$\pm$ & \cellcolor{yes!35}+ & \textcolor{gray}{·} & \textcolor{gray}{·} \\
student-mixture & \cellcolor{part!45}$\pm$ & \textcolor{gray}{·} & \cellcolor{yes!35}+ & \textcolor{gray}{·} & \cellcolor{yes!35}+ \\
kasakh Mixture-Models & \cellcolor{part!45}$\pm$ & \textcolor{gray}{·} & \cellcolor{yes!35}+ & \textcolor{gray}{·} & \cellcolor{part!45}$\pm$ \\
Armadillo C++ & \textcolor{gray}{·} & \textcolor{gray}{·} & \cellcolor{yes!35}+ & \textcolor{gray}{·} & \textcolor{gray}{·} \\
mixem & \cellcolor{part!45}$\pm$ & \cellcolor{yes!35}+ & \cellcolor{yes!35}+ & \cellcolor{yes!35}+ & \textcolor{gray}{·} \\
\midrule \textit{Выполнено, шт.} & 5 & 7 & 22 & 16 & 8 \\
\bottomrule
\end{longtable}}
{\scriptsize\setlength{\tabcolsep}{3pt}
\begin{tabular}{@{}l >{\raggedright\arraybackslash}p{7.4cm} l l >{\raggedright\arraybackslash}p{7.4cm} l@{}}
\texttt{TECH-01} & Композитный модульный пайплайн оценивания & обязательные & \texttt{TECH-04} & Иммутабельность структур параметров модели & желательные \\
\texttt{TECH-02} & Расширяемость новыми распределениями & обязательные & \texttt{TECH-05} & scikit-learn совместимый интерфейс & желательные \\
\texttt{TECH-03} & Логарифмическая стабильность вычислений & обязательные &  & &  \\
\end{tabular}}

\subsection{Инфраструктура}
{\scriptsize\setlength{\tabcolsep}{2.2pt}
\begin{longtable}{@{}l *{6}{c}@{}}
\toprule
Инструмент & \rotatebox{90}{\scriptsize INFRA-01} & \rotatebox{90}{\scriptsize INFRA-02} & \rotatebox{90}{\scriptsize INFRA-03} & \rotatebox{90}{\scriptsize INFRA-04} & \rotatebox{90}{\scriptsize INFRA-05} & \rotatebox{90}{\scriptsize INFRA-06} \\
\midrule
\endhead
brms & \cellcolor{yes!35}+ & \cellcolor{yes!35}+ & \textcolor{gray}{·} & \cellcolor{part!45}$\pm$ & \cellcolor{yes!35}+ & \cellcolor{yes!35}+ \\
FlexMix & \cellcolor{yes!35}+ & \cellcolor{yes!35}+ & \textcolor{gray}{·} & \cellcolor{part!45}$\pm$ & \cellcolor{yes!35}+ & \cellcolor{yes!35}+ \\
mclust (v5+) & \cellcolor{yes!35}+ & \cellcolor{yes!35}+ & \textcolor{gray}{·} & \cellcolor{part!45}$\pm$ & \cellcolor{yes!35}+ & \cellcolor{yes!35}+ \\
bnpy & \cellcolor{part!45}$\pm$ & \cellcolor{part!45}$\pm$ & \textcolor{gray}{·} & \textcolor{gray}{·} & \cellcolor{part!45}$\pm$ & \cellcolor{part!45}$\pm$ \\
mixR & \cellcolor{yes!35}+ & \cellcolor{yes!35}+ & \textcolor{gray}{·} & \cellcolor{part!45}$\pm$ & \cellcolor{yes!35}+ & \cellcolor{yes!35}+ \\
PySATL mpest & \cellcolor{yes!35}+ & \cellcolor{yes!35}+ & \cellcolor{yes!35}+ & \cellcolor{yes!35}+ & \cellcolor{yes!35}+ & \cellcolor{yes!35}+ \\
StepMix & \cellcolor{yes!35}+ & \cellcolor{yes!35}+ & \textcolor{gray}{·} & \cellcolor{part!45}$\pm$ & \cellcolor{yes!35}+ & \cellcolor{yes!35}+ \\
Rmixmod (mixmod) & \cellcolor{yes!35}+ & \cellcolor{yes!35}+ & \textcolor{gray}{·} & \cellcolor{part!45}$\pm$ & \cellcolor{yes!35}+ & \cellcolor{yes!35}+ \\
mixtools & \cellcolor{yes!35}+ & \cellcolor{part!45}$\pm$ & \textcolor{gray}{·} & \cellcolor{part!45}$\pm$ & \cellcolor{yes!35}+ & \cellcolor{yes!35}+ \\
rebmix & \cellcolor{yes!35}+ & \cellcolor{yes!35}+ & \textcolor{gray}{·} & \cellcolor{part!45}$\pm$ & \cellcolor{yes!35}+ & \cellcolor{yes!35}+ \\
pomegranate & \cellcolor{part!45}$\pm$ & \cellcolor{yes!35}+ & \textcolor{gray}{·} & \textcolor{gray}{·} & \cellcolor{yes!35}+ & \cellcolor{yes!35}+ \\
Distributions.jl & \cellcolor{yes!35}+ & \cellcolor{yes!35}+ & \cellcolor{part!45}$\pm$ & \cellcolor{yes!35}+ & \cellcolor{yes!35}+ & \cellcolor{yes!35}+ \\
distfit & \cellcolor{yes!35}+ & \cellcolor{yes!35}+ & \textcolor{gray}{·} & \cellcolor{part!45}$\pm$ & \cellcolor{yes!35}+ & \cellcolor{yes!35}+ \\
mixdist & \cellcolor{yes!35}+ & \cellcolor{part!45}$\pm$ & \textcolor{gray}{·} & \cellcolor{part!45}$\pm$ & \cellcolor{yes!35}+ & \cellcolor{yes!35}+ \\
scikit-learn (sklearn.mixture) & \cellcolor{yes!35}+ & \cellcolor{yes!35}+ & \cellcolor{part!45}$\pm$ & \cellcolor{yes!35}+ & \cellcolor{yes!35}+ & \cellcolor{yes!35}+ \\
sn (Skew-Normal/t) & \cellcolor{yes!35}+ & \cellcolor{yes!35}+ & \textcolor{gray}{·} & \cellcolor{part!45}$\pm$ & \cellcolor{yes!35}+ & \cellcolor{yes!35}+ \\
EMCluster & \cellcolor{yes!35}+ & \cellcolor{yes!35}+ & \textcolor{gray}{·} & \cellcolor{part!45}$\pm$ & \cellcolor{yes!35}+ & \cellcolor{yes!35}+ \\
teigen & \cellcolor{yes!35}+ & \cellcolor{part!45}$\pm$ & \textcolor{gray}{·} & \cellcolor{part!45}$\pm$ & \cellcolor{yes!35}+ & \cellcolor{yes!35}+ \\
CAMAN & \cellcolor{yes!35}+ & \textcolor{gray}{·} & \textcolor{gray}{·} & \textcolor{gray}{·} & \textcolor{gray}{·} & \cellcolor{yes!35}+ \\
GaussianMixtures.jl & \cellcolor{yes!35}+ & \cellcolor{yes!35}+ & \textcolor{gray}{·} & \cellcolor{part!45}$\pm$ & \cellcolor{yes!35}+ & \cellcolor{yes!35}+ \\
fpc & \cellcolor{yes!35}+ & \cellcolor{part!45}$\pm$ & \textcolor{gray}{·} & \cellcolor{part!45}$\pm$ & \cellcolor{yes!35}+ & \cellcolor{yes!35}+ \\
ClusterR & \cellcolor{yes!35}+ & \cellcolor{yes!35}+ & \textcolor{gray}{·} & \cellcolor{part!45}$\pm$ & \cellcolor{yes!35}+ & \cellcolor{yes!35}+ \\
mlpack (mlpack::gmm) & \cellcolor{yes!35}+ & \cellcolor{yes!35}+ & \textcolor{gray}{·} & \cellcolor{yes!35}+ & \cellcolor{yes!35}+ & \cellcolor{yes!35}+ \\
student-mixture & \textcolor{gray}{·} & \cellcolor{part!45}$\pm$ & \textcolor{gray}{·} & \textcolor{gray}{·} & \textcolor{gray}{·} & \cellcolor{part!45}$\pm$ \\
kasakh Mixture-Models & \textcolor{gray}{·} & \cellcolor{yes!35}+ & \textcolor{gray}{·} & \textcolor{gray}{·} & \cellcolor{yes!35}+ & \cellcolor{part!45}$\pm$ \\
Armadillo C++ & \cellcolor{yes!35}+ & \cellcolor{yes!35}+ & \textcolor{gray}{·} & \cellcolor{part!45}$\pm$ & \cellcolor{yes!35}+ & \cellcolor{yes!35}+ \\
mixem & \textcolor{gray}{·} & \cellcolor{yes!35}+ & \textcolor{gray}{·} & \textcolor{gray}{·} & \cellcolor{yes!35}+ & \cellcolor{yes!35}+ \\
\midrule \textit{Выполнено, шт.} & 22 & 20 & 1 & 4 & 24 & 24 \\
\bottomrule
\end{longtable}}
{\scriptsize\setlength{\tabcolsep}{3pt}
\begin{tabular}{@{}l >{\raggedright\arraybackslash}p{7.4cm} l l >{\raggedright\arraybackslash}p{7.4cm} l@{}}
\texttt{INFRA-01} & Декларативное управление зависимостями и сборкой & обязательные & \texttt{INFRA-04} & Автоматизированный статический анализ и линтинг & обязательные \\
\texttt{INFRA-02} & Юнит- и интеграционное тестирование с высоким покрытием & обязательные & \texttt{INFRA-05} & Автоматизированная непрерывная интеграция & обязательные \\
\texttt{INFRA-03} & Property-based тестирование математических инвариантов & обязательные & \texttt{INFRA-06} & Автоматическая генерация документации с математикой & обязательные \\
\end{tabular}}

\subsection{Визуализация}
{\scriptsize\setlength{\tabcolsep}{2.2pt}
\begin{longtable}{@{}l *{5}{c}@{}}
\toprule
Инструмент & \rotatebox{90}{\scriptsize VIS-01} & \rotatebox{90}{\scriptsize VIS-02} & \rotatebox{90}{\scriptsize VIS-03} & \rotatebox{90}{\scriptsize VIS-04} & \rotatebox{90}{\scriptsize VIS-05} \\
\midrule
\endhead
brms & \cellcolor{yes!35}+ & \cellcolor{yes!35}+ & \cellcolor{part!45}$\pm$ & \cellcolor{yes!35}+ & \cellcolor{part!45}$\pm$ \\
FlexMix & \cellcolor{yes!35}+ & \cellcolor{part!45}$\pm$ & \cellcolor{yes!35}+ & \textcolor{gray}{·} & \cellcolor{yes!35}+ \\
mclust (v5+) & \cellcolor{yes!35}+ & \textcolor{gray}{·} & \cellcolor{yes!35}+ & \textcolor{gray}{·} & \cellcolor{yes!35}+ \\
bnpy & \cellcolor{part!45}$\pm$ & \textcolor{gray}{·} & \cellcolor{part!45}$\pm$ & \cellcolor{yes!35}+ & \cellcolor{part!45}$\pm$ \\
mixR & \cellcolor{yes!35}+ & \cellcolor{yes!35}+ & \cellcolor{yes!35}+ & \textcolor{gray}{·} & \textcolor{gray}{·} \\
PySATL mpest & \textcolor{gray}{·} & \textcolor{gray}{·} & \textcolor{gray}{·} & \textcolor{gray}{·} & \textcolor{gray}{·} \\
StepMix & \cellcolor{yes!35}+ & \textcolor{gray}{·} & \cellcolor{yes!35}+ & \textcolor{gray}{·} & \cellcolor{yes!35}+ \\
Rmixmod (mixmod) & \cellcolor{yes!35}+ & \textcolor{gray}{·} & \cellcolor{yes!35}+ & \cellcolor{yes!35}+ & \textcolor{gray}{·} \\
mixtools & \cellcolor{yes!35}+ & \textcolor{gray}{·} & \textcolor{gray}{·} & \cellcolor{yes!35}+ & \cellcolor{yes!35}+ \\
rebmix & \cellcolor{yes!35}+ & \textcolor{gray}{·} & \cellcolor{part!45}$\pm$ & \textcolor{gray}{·} & \textcolor{gray}{·} \\
pomegranate & \textcolor{gray}{·} & \textcolor{gray}{·} & \textcolor{gray}{·} & \textcolor{gray}{·} & \textcolor{gray}{·} \\
Distributions.jl & \cellcolor{part!45}$\pm$ & \textcolor{gray}{·} & \textcolor{gray}{·} & \textcolor{gray}{·} & \textcolor{gray}{·} \\
distfit & \cellcolor{yes!35}+ & \cellcolor{yes!35}+ & \cellcolor{yes!35}+ & \textcolor{gray}{·} & \textcolor{gray}{·} \\
mixdist & \cellcolor{yes!35}+ & \textcolor{gray}{·} & \textcolor{gray}{·} & \textcolor{gray}{·} & \textcolor{gray}{·} \\
scikit-learn (sklearn.mixture) & \textcolor{gray}{·} & \textcolor{gray}{·} & \textcolor{gray}{·} & \textcolor{gray}{·} & \textcolor{gray}{·} \\
sn (Skew-Normal/t) & \cellcolor{yes!35}+ & \textcolor{gray}{·} & \textcolor{gray}{·} & \textcolor{gray}{·} & \textcolor{gray}{·} \\
EMCluster & \cellcolor{yes!35}+ & \textcolor{gray}{·} & \textcolor{gray}{·} & \textcolor{gray}{·} & \textcolor{gray}{·} \\
teigen & \cellcolor{yes!35}+ & \textcolor{gray}{·} & \cellcolor{yes!35}+ & \textcolor{gray}{·} & \cellcolor{yes!35}+ \\
CAMAN & \cellcolor{yes!35}+ & \textcolor{gray}{·} & \cellcolor{yes!35}+ & \cellcolor{yes!35}+ & \textcolor{gray}{·} \\
GaussianMixtures.jl & \textcolor{gray}{·} & \textcolor{gray}{·} & \textcolor{gray}{·} & \textcolor{gray}{·} & \textcolor{gray}{·} \\
fpc & \cellcolor{yes!35}+ & \textcolor{gray}{·} & \textcolor{gray}{·} & \textcolor{gray}{·} & \textcolor{gray}{·} \\
ClusterR & \textcolor{gray}{·} & \textcolor{gray}{·} & \cellcolor{yes!35}+ & \textcolor{gray}{·} & \textcolor{gray}{·} \\
mlpack (mlpack::gmm) & \textcolor{gray}{·} & \textcolor{gray}{·} & \textcolor{gray}{·} & \textcolor{gray}{·} & \textcolor{gray}{·} \\
student-mixture & \textcolor{gray}{·} & \textcolor{gray}{·} & \textcolor{gray}{·} & \textcolor{gray}{·} & \textcolor{gray}{·} \\
kasakh Mixture-Models & \textcolor{gray}{·} & \textcolor{gray}{·} & \textcolor{gray}{·} & \textcolor{gray}{·} & \textcolor{gray}{·} \\
Armadillo C++ & \textcolor{gray}{·} & \textcolor{gray}{·} & \textcolor{gray}{·} & \textcolor{gray}{·} & \textcolor{gray}{·} \\
mixem & \textcolor{gray}{·} & \textcolor{gray}{·} & \textcolor{gray}{·} & \textcolor{gray}{·} & \textcolor{gray}{·} \\
\midrule \textit{Выполнено, шт.} & 15 & 3 & 9 & 5 & 5 \\
\bottomrule
\end{longtable}}
{\scriptsize\setlength{\tabcolsep}{3pt}
\begin{tabular}{@{}l >{\raggedright\arraybackslash}p{7.4cm} l l >{\raggedright\arraybackslash}p{7.4cm} l@{}}
\texttt{VIS-01} & Декомпозиция плотности смеси & обязательные & \texttt{VIS-04} & Визуализация скорости сходимости & факультативные \\
\texttt{VIS-02} & Диагностические графики согласия (Q-Q, P-P) & желательные & \texttt{VIS-05} & Графики апостериорной неопределенности & факультативные \\
\texttt{VIS-03} & Профиль выбора числа компонент (AIC, BIC, ICL) & желательные &  & &  \\
\end{tabular}}

\subsection{Специальные методы}
{\scriptsize\setlength{\tabcolsep}{2.2pt}
\begin{longtable}{@{}l *{5}{c}@{}}
\toprule
Инструмент & \rotatebox{90}{\scriptsize PROP-01} & \rotatebox{90}{\scriptsize PROP-02} & \rotatebox{90}{\scriptsize PROP-03} & \rotatebox{90}{\scriptsize PROP-04} & \rotatebox{90}{\scriptsize PROP-05} \\
\midrule
\endhead
brms & \cellcolor{yes!35}+ & \textcolor{gray}{·} & \cellcolor{yes!35}+ & \textcolor{gray}{·} & \textcolor{gray}{·} \\
FlexMix & \cellcolor{yes!35}+ & \textcolor{gray}{·} & \cellcolor{part!45}$\pm$ & \textcolor{gray}{·} & \textcolor{gray}{·} \\
mclust (v5+) & \cellcolor{part!45}$\pm$ & \cellcolor{part!45}$\pm$ & \cellcolor{yes!35}+ & \textcolor{gray}{·} & \textcolor{gray}{·} \\
bnpy & \textcolor{gray}{·} & \textcolor{gray}{·} & \cellcolor{yes!35}+ & \textcolor{gray}{·} & \textcolor{gray}{·} \\
mixR & \cellcolor{part!45}$\pm$ & \cellcolor{part!45}$\pm$ & \textcolor{gray}{·} & \textcolor{gray}{·} & \textcolor{gray}{·} \\
PySATL mpest & \textcolor{gray}{·} & \textcolor{gray}{·} & \textcolor{gray}{·} & \textcolor{gray}{·} & \textcolor{gray}{·} \\
StepMix & \cellcolor{yes!35}+ & \textcolor{gray}{·} & \cellcolor{part!45}$\pm$ & \textcolor{gray}{·} & \textcolor{gray}{·} \\
Rmixmod (mixmod) & \cellcolor{part!45}$\pm$ & \textcolor{gray}{·} & \cellcolor{part!45}$\pm$ & \textcolor{gray}{·} & \textcolor{gray}{·} \\
mixtools & \cellcolor{part!45}$\pm$ & \cellcolor{part!45}$\pm$ & \cellcolor{part!45}$\pm$ & \textcolor{gray}{·} & \textcolor{gray}{·} \\
rebmix & \textcolor{gray}{·} & \textcolor{gray}{·} & \textcolor{gray}{·} & \textcolor{gray}{·} & \textcolor{gray}{·} \\
pomegranate & \cellcolor{part!45}$\pm$ & \textcolor{gray}{·} & \cellcolor{part!45}$\pm$ & \textcolor{gray}{·} & \textcolor{gray}{·} \\
Distributions.jl & \textcolor{gray}{·} & \textcolor{gray}{·} & \textcolor{gray}{·} & \textcolor{gray}{·} & \textcolor{gray}{·} \\
distfit & \textcolor{gray}{·} & \textcolor{gray}{·} & \textcolor{gray}{·} & \textcolor{gray}{·} & \cellcolor{part!45}$\pm$ \\
mixdist & \cellcolor{yes!35}+ & \textcolor{gray}{·} & \textcolor{gray}{·} & \textcolor{gray}{·} & \textcolor{gray}{·} \\
scikit-learn (sklearn.mixture) & \cellcolor{part!45}$\pm$ & \textcolor{gray}{·} & \cellcolor{yes!35}+ & \textcolor{gray}{·} & \textcolor{gray}{·} \\
sn (Skew-Normal/t) & \cellcolor{part!45}$\pm$ & \textcolor{gray}{·} & \textcolor{gray}{·} & \textcolor{gray}{·} & \textcolor{gray}{·} \\
EMCluster & \cellcolor{part!45}$\pm$ & \textcolor{gray}{·} & \cellcolor{part!45}$\pm$ & \textcolor{gray}{·} & \textcolor{gray}{·} \\
teigen & \cellcolor{part!45}$\pm$ & \textcolor{gray}{·} & \textcolor{gray}{·} & \textcolor{gray}{·} & \textcolor{gray}{·} \\
CAMAN & \textcolor{gray}{·} & \cellcolor{part!45}$\pm$ & \textcolor{gray}{·} & \textcolor{gray}{·} & \textcolor{gray}{·} \\
GaussianMixtures.jl & \textcolor{gray}{·} & \textcolor{gray}{·} & \textcolor{gray}{·} & \textcolor{gray}{·} & \textcolor{gray}{·} \\
fpc & \textcolor{gray}{·} & \textcolor{gray}{·} & \cellcolor{part!45}$\pm$ & \textcolor{gray}{·} & \textcolor{gray}{·} \\
ClusterR & \textcolor{gray}{·} & \textcolor{gray}{·} & \cellcolor{part!45}$\pm$ & \textcolor{gray}{·} & \textcolor{gray}{·} \\
mlpack (mlpack::gmm) & \textcolor{gray}{·} & \textcolor{gray}{·} & \cellcolor{part!45}$\pm$ & \textcolor{gray}{·} & \textcolor{gray}{·} \\
student-mixture & \textcolor{gray}{·} & \textcolor{gray}{·} & \cellcolor{part!45}$\pm$ & \textcolor{gray}{·} & \textcolor{gray}{·} \\
kasakh Mixture-Models & \textcolor{gray}{·} & \textcolor{gray}{·} & \textcolor{gray}{·} & \textcolor{gray}{·} & \textcolor{gray}{·} \\
Armadillo C++ & \textcolor{gray}{·} & \textcolor{gray}{·} & \textcolor{gray}{·} & \textcolor{gray}{·} & \textcolor{gray}{·} \\
mixem & \textcolor{gray}{·} & \textcolor{gray}{·} & \textcolor{gray}{·} & \textcolor{gray}{·} & \textcolor{gray}{·} \\
\midrule \textit{Выполнено, шт.} & 4 & 0 & 4 & 0 & 0 \\
\bottomrule
\end{longtable}}
{\scriptsize\setlength{\tabcolsep}{3pt}
\begin{tabular}{@{}l >{\raggedright\arraybackslash}p{7.4cm} l l >{\raggedright\arraybackslash}p{7.4cm} l@{}}
\texttt{PROP-01} & Декларативное задание параметрических ограничений (RESTRICT) & факультативные & \texttt{PROP-04} & Символьно-аналитическое оценивание смесей высокой точности & факультативные \\
\texttt{PROP-02} & Скорректированные тесты отношения правдоподобий (VLMR, LMR, MLRT) & факультативные & \texttt{PROP-05} & Автоматический эвристический перебор аналитических форм компонент & факультативные \\
\texttt{PROP-03} & Адаптивная регуляризация и стабилизация масштаба & факультативные &  & &  \\
\end{tabular}}

\subsection{Представление смеси}
{\scriptsize\setlength{\tabcolsep}{2.2pt}
\begin{longtable}{@{}l *{4}{c}@{}}
\toprule
Инструмент & \rotatebox{90}{\scriptsize REPR-01} & \rotatebox{90}{\scriptsize REPR-02} & \rotatebox{90}{\scriptsize REPR-03} & \rotatebox{90}{\scriptsize REPR-04} \\
\midrule
\endhead
brms & \cellcolor{part!45}$\pm$ & \cellcolor{yes!35}+ & \cellcolor{part!45}$\pm$ & \cellcolor{part!45}$\pm$ \\
FlexMix & \textcolor{gray}{·} & \cellcolor{yes!35}+ & \textcolor{gray}{·} & \textcolor{gray}{·} \\
mclust (v5+) & \cellcolor{part!45}$\pm$ & \cellcolor{yes!35}+ & \cellcolor{part!45}$\pm$ & \textcolor{gray}{·} \\
bnpy & \cellcolor{part!45}$\pm$ & \cellcolor{yes!35}+ & \textcolor{gray}{·} & \cellcolor{part!45}$\pm$ \\
mixR & \cellcolor{yes!35}+ & \cellcolor{yes!35}+ & \cellcolor{part!45}$\pm$ & \textcolor{gray}{·} \\
PySATL mpest & \cellcolor{yes!35}+ & \cellcolor{yes!35}+ & \cellcolor{part!45}$\pm$ & \textcolor{gray}{·} \\
StepMix & \textcolor{gray}{·} & \cellcolor{yes!35}+ & \textcolor{gray}{·} & \textcolor{gray}{·} \\
Rmixmod (mixmod) & \textcolor{gray}{·} & \cellcolor{yes!35}+ & \textcolor{gray}{·} & \cellcolor{yes!35}+ \\
mixtools & \textcolor{gray}{·} & \cellcolor{yes!35}+ & \textcolor{gray}{·} & \textcolor{gray}{·} \\
rebmix & \cellcolor{part!45}$\pm$ & \cellcolor{yes!35}+ & \cellcolor{part!45}$\pm$ & \textcolor{gray}{·} \\
pomegranate & \cellcolor{part!45}$\pm$ & \cellcolor{yes!35}+ & \textcolor{gray}{·} & \cellcolor{yes!35}+ \\
Distributions.jl & \cellcolor{yes!35}+ & \cellcolor{yes!35}+ & \cellcolor{yes!35}+ & \cellcolor{part!45}$\pm$ \\
distfit & \cellcolor{yes!35}+ & \cellcolor{yes!35}+ & \cellcolor{part!45}$\pm$ & \cellcolor{yes!35}+ \\
mixdist & \textcolor{gray}{·} & \textcolor{gray}{·} & \textcolor{gray}{·} & \textcolor{gray}{·} \\
scikit-learn (sklearn.mixture) & \textcolor{gray}{·} & \cellcolor{yes!35}+ & \textcolor{gray}{·} & \cellcolor{part!45}$\pm$ \\
sn (Skew-Normal/t) & \cellcolor{yes!35}+ & \cellcolor{yes!35}+ & \cellcolor{yes!35}+ & \textcolor{gray}{·} \\
EMCluster & \textcolor{gray}{·} & \cellcolor{yes!35}+ & \textcolor{gray}{·} & \textcolor{gray}{·} \\
teigen & \textcolor{gray}{·} & \textcolor{gray}{·} & \textcolor{gray}{·} & \textcolor{gray}{·} \\
CAMAN & \cellcolor{part!45}$\pm$ & \cellcolor{part!45}$\pm$ & \textcolor{gray}{·} & \cellcolor{part!45}$\pm$ \\
GaussianMixtures.jl & \textcolor{gray}{·} & \cellcolor{yes!35}+ & \textcolor{gray}{·} & \cellcolor{part!45}$\pm$ \\
fpc & \textcolor{gray}{·} & \textcolor{gray}{·} & \textcolor{gray}{·} & \textcolor{gray}{·} \\
ClusterR & \textcolor{gray}{·} & \textcolor{gray}{·} & \textcolor{gray}{·} & \textcolor{gray}{·} \\
mlpack (mlpack::gmm) & \textcolor{gray}{·} & \cellcolor{yes!35}+ & \textcolor{gray}{·} & \cellcolor{yes!35}+ \\
student-mixture & \textcolor{gray}{·} & \cellcolor{yes!35}+ & \textcolor{gray}{·} & \textcolor{gray}{·} \\
kasakh Mixture-Models & \textcolor{gray}{·} & \cellcolor{yes!35}+ & \textcolor{gray}{·} & \textcolor{gray}{·} \\
Armadillo C++ & \textcolor{gray}{·} & \cellcolor{yes!35}+ & \textcolor{gray}{·} & \cellcolor{part!45}$\pm$ \\
mixem & \cellcolor{part!45}$\pm$ & \textcolor{gray}{·} & \textcolor{gray}{·} & \textcolor{gray}{·} \\
\midrule \textit{Выполнено, шт.} & 5 & 21 & 2 & 4 \\
\bottomrule
\end{longtable}}
{\scriptsize\setlength{\tabcolsep}{3pt}
\begin{tabular}{@{}l >{\raggedright\arraybackslash}p{7.4cm} l l >{\raggedright\arraybackslash}p{7.4cm} l@{}}
\texttt{REPR-01} & Поддержка вероятностных функций (PDF, CDF, PPF) & обязательные & \texttt{REPR-03} & Аналитический расчет моментов смеси & желательные \\
\texttt{REPR-02} & Сэмплирование & обязательные & \texttt{REPR-04} & Сериализация смеси в JSON & желательные \\
\end{tabular}}

\subsection{Производительность}
{\scriptsize\setlength{\tabcolsep}{2.2pt}
\begin{longtable}{@{}l *{4}{c}@{}}
\toprule
Инструмент & \rotatebox{90}{\scriptsize PERF-01} & \rotatebox{90}{\scriptsize PERF-02} & \rotatebox{90}{\scriptsize PERF-03} & \rotatebox{90}{\scriptsize PERF-04} \\
\midrule
\endhead
brms & \cellcolor{yes!35}+ & \cellcolor{part!45}$\pm$ & \cellcolor{yes!35}+ & \cellcolor{yes!35}+ \\
FlexMix & \cellcolor{part!45}$\pm$ & \cellcolor{part!45}$\pm$ & \cellcolor{part!45}$\pm$ & \textcolor{gray}{·} \\
mclust (v5+) & \cellcolor{yes!35}+ & \cellcolor{part!45}$\pm$ & \cellcolor{yes!35}+ & \cellcolor{part!45}$\pm$ \\
bnpy & \cellcolor{yes!35}+ & \cellcolor{yes!35}+ & \cellcolor{yes!35}+ & \cellcolor{yes!35}+ \\
mixR & \textcolor{gray}{·} & \textcolor{gray}{·} & \cellcolor{yes!35}+ & \cellcolor{yes!35}+ \\
PySATL mpest & \textcolor{gray}{·} & \cellcolor{yes!35}+ & \cellcolor{yes!35}+ & \textcolor{gray}{·} \\
StepMix & \cellcolor{yes!35}+ & \cellcolor{part!45}$\pm$ & \textcolor{gray}{·} & \cellcolor{part!45}$\pm$ \\
Rmixmod (mixmod) & \cellcolor{yes!35}+ & \cellcolor{part!45}$\pm$ & \cellcolor{yes!35}+ & \cellcolor{yes!35}+ \\
mixtools & \textcolor{gray}{·} & \textcolor{gray}{·} & \cellcolor{part!45}$\pm$ & \textcolor{gray}{·} \\
rebmix & \textcolor{gray}{·} & \cellcolor{yes!35}+ & \cellcolor{yes!35}+ & \cellcolor{yes!35}+ \\
pomegranate & \cellcolor{part!45}$\pm$ & \textcolor{gray}{·} & \cellcolor{yes!35}+ & \cellcolor{yes!35}+ \\
Distributions.jl & \cellcolor{part!45}$\pm$ & \cellcolor{part!45}$\pm$ & \cellcolor{yes!35}+ & \cellcolor{yes!35}+ \\
distfit & \cellcolor{part!45}$\pm$ & \textcolor{gray}{·} & \cellcolor{part!45}$\pm$ & \cellcolor{part!45}$\pm$ \\
mixdist & \textcolor{gray}{·} & \textcolor{gray}{·} & \textcolor{gray}{·} & \cellcolor{yes!35}+ \\
scikit-learn (sklearn.mixture) & \cellcolor{yes!35}+ & \cellcolor{part!45}$\pm$ & \cellcolor{yes!35}+ & \cellcolor{yes!35}+ \\
sn (Skew-Normal/t) & \textcolor{gray}{·} & \textcolor{gray}{·} & \cellcolor{yes!35}+ & \textcolor{gray}{·} \\
EMCluster & \textcolor{gray}{·} & \textcolor{gray}{·} & \cellcolor{yes!35}+ & \cellcolor{yes!35}+ \\
teigen & \cellcolor{part!45}$\pm$ & \cellcolor{part!45}$\pm$ & \cellcolor{yes!35}+ & \textcolor{gray}{·} \\
CAMAN & \textcolor{gray}{·} & \textcolor{gray}{·} & \cellcolor{yes!35}+ & \textcolor{gray}{·} \\
GaussianMixtures.jl & \cellcolor{yes!35}+ & \cellcolor{yes!35}+ & \cellcolor{yes!35}+ & \cellcolor{yes!35}+ \\
fpc & \textcolor{gray}{·} & \textcolor{gray}{·} & \cellcolor{part!45}$\pm$ & \textcolor{gray}{·} \\
ClusterR & \cellcolor{yes!35}+ & \textcolor{gray}{·} & \cellcolor{yes!35}+ & \cellcolor{yes!35}+ \\
mlpack (mlpack::gmm) & \cellcolor{yes!35}+ & \textcolor{gray}{·} & \cellcolor{yes!35}+ & \cellcolor{yes!35}+ \\
student-mixture & \textcolor{gray}{·} & \textcolor{gray}{·} & \textcolor{gray}{·} & \textcolor{gray}{·} \\
kasakh Mixture-Models & \cellcolor{part!45}$\pm$ & \textcolor{gray}{·} & \cellcolor{yes!35}+ & \textcolor{gray}{·} \\
Armadillo C++ & \cellcolor{yes!35}+ & \textcolor{gray}{·} & \cellcolor{yes!35}+ & \cellcolor{yes!35}+ \\
mixem & \textcolor{gray}{·} & \textcolor{gray}{·} & \textcolor{gray}{·} & \textcolor{gray}{·} \\
\midrule \textit{Выполнено, шт.} & 10 & 4 & 19 & 14 \\
\bottomrule
\end{longtable}}
{\scriptsize\setlength{\tabcolsep}{3pt}
\begin{tabular}{@{}l >{\raggedright\arraybackslash}p{7.4cm} l l >{\raggedright\arraybackslash}p{7.4cm} l@{}}
\texttt{PERF-01} & Параллелизация независимых вычислений & желательные & \texttt{PERF-03} & Векторизованные вычисления (NumPy, SIMD) & желательные \\
\texttt{PERF-02} & Кэширование промежуточных структур и плотностей & желательные & \texttt{PERF-04} & Контроль аллокаций памяти и потоковые вычисления & факультативные \\
\end{tabular}}

\subsection{Диагностика}
{\scriptsize\setlength{\tabcolsep}{2.2pt}
\begin{longtable}{@{}l *{4}{c}@{}}
\toprule
Инструмент & \rotatebox{90}{\scriptsize DIAG-01} & \rotatebox{90}{\scriptsize DIAG-02} & \rotatebox{90}{\scriptsize DIAG-03} & \rotatebox{90}{\scriptsize DIAG-04} \\
\midrule
\endhead
brms & \cellcolor{part!45}$\pm$ & \cellcolor{yes!35}+ & \cellcolor{yes!35}+ & \textcolor{gray}{·} \\
FlexMix & \textcolor{gray}{·} & \cellcolor{yes!35}+ & \textcolor{gray}{·} & \cellcolor{part!45}$\pm$ \\
mclust (v5+) & \textcolor{gray}{·} & \cellcolor{yes!35}+ & \textcolor{gray}{·} & \cellcolor{part!45}$\pm$ \\
bnpy & \textcolor{gray}{·} & \textcolor{gray}{·} & \cellcolor{yes!35}+ & \cellcolor{part!45}$\pm$ \\
mixR & \textcolor{gray}{·} & \textcolor{gray}{·} & \textcolor{gray}{·} & \textcolor{gray}{·} \\
PySATL mpest & \textcolor{gray}{·} & \textcolor{gray}{·} & \textcolor{gray}{·} & \textcolor{gray}{·} \\
StepMix & \textcolor{gray}{·} & \cellcolor{yes!35}+ & \textcolor{gray}{·} & \cellcolor{part!45}$\pm$ \\
Rmixmod (mixmod) & \textcolor{gray}{·} & \textcolor{gray}{·} & \textcolor{gray}{·} & \textcolor{gray}{·} \\
mixtools & \textcolor{gray}{·} & \cellcolor{yes!35}+ & \textcolor{gray}{·} & \cellcolor{part!45}$\pm$ \\
rebmix & \textcolor{gray}{·} & \cellcolor{yes!35}+ & \textcolor{gray}{·} & \textcolor{gray}{·} \\
pomegranate & \textcolor{gray}{·} & \textcolor{gray}{·} & \textcolor{gray}{·} & \textcolor{gray}{·} \\
Distributions.jl & \cellcolor{part!45}$\pm$ & \textcolor{gray}{·} & \textcolor{gray}{·} & \textcolor{gray}{·} \\
distfit & \cellcolor{yes!35}+ & \textcolor{gray}{·} & \textcolor{gray}{·} & \textcolor{gray}{·} \\
mixdist & \cellcolor{yes!35}+ & \cellcolor{yes!35}+ & \textcolor{gray}{·} & \textcolor{gray}{·} \\
scikit-learn (sklearn.mixture) & \textcolor{gray}{·} & \textcolor{gray}{·} & \textcolor{gray}{·} & \textcolor{gray}{·} \\
sn (Skew-Normal/t) & \textcolor{gray}{·} & \cellcolor{yes!35}+ & \textcolor{gray}{·} & \textcolor{gray}{·} \\
EMCluster & \textcolor{gray}{·} & \cellcolor{yes!35}+ & \textcolor{gray}{·} & \cellcolor{part!45}$\pm$ \\
teigen & \textcolor{gray}{·} & \textcolor{gray}{·} & \textcolor{gray}{·} & \cellcolor{part!45}$\pm$ \\
CAMAN & \textcolor{gray}{·} & \cellcolor{yes!35}+ & \textcolor{gray}{·} & \textcolor{gray}{·} \\
GaussianMixtures.jl & \textcolor{gray}{·} & \textcolor{gray}{·} & \textcolor{gray}{·} & \textcolor{gray}{·} \\
fpc & \textcolor{gray}{·} & \cellcolor{part!45}$\pm$ & \textcolor{gray}{·} & \cellcolor{yes!35}+ \\
ClusterR & \textcolor{gray}{·} & \textcolor{gray}{·} & \cellcolor{yes!35}+ & \textcolor{gray}{·} \\
mlpack (mlpack::gmm) & \textcolor{gray}{·} & \textcolor{gray}{·} & \cellcolor{yes!35}+ & \textcolor{gray}{·} \\
student-mixture & \textcolor{gray}{·} & \textcolor{gray}{·} & \textcolor{gray}{·} & \textcolor{gray}{·} \\
kasakh Mixture-Models & \textcolor{gray}{·} & \textcolor{gray}{·} & \textcolor{gray}{·} & \textcolor{gray}{·} \\
Armadillo C++ & \textcolor{gray}{·} & \textcolor{gray}{·} & \textcolor{gray}{·} & \textcolor{gray}{·} \\
mixem & \textcolor{gray}{·} & \textcolor{gray}{·} & \textcolor{gray}{·} & \textcolor{gray}{·} \\
\midrule \textit{Выполнено, шт.} & 2 & 10 & 4 & 1 \\
\bottomrule
\end{longtable}}
{\scriptsize\setlength{\tabcolsep}{3pt}
\begin{tabular}{@{}l >{\raggedright\arraybackslash}p{7.4cm} l l >{\raggedright\arraybackslash}p{7.4cm} l@{}}
\texttt{DIAG-01} & Статистические критерии согласия (KS, AD, Chi2) & желательные & \texttt{DIAG-03} & Бенчмарки производительности (ASV, pytest-benchmark) & факультативные \\
\texttt{DIAG-02} & Стандартные ошибки и доверительные интервалы оценок & желательные & \texttt{DIAG-04} & Оценка отделимости и перекрытия компонент & факультативные \\
\end{tabular}}

\end{landscape}
\part{Прикладные сценарии}
\section{Характеристики сценариев}
Каждый сценарий описывается пятью признаками: типом смеси (однородная или разнородная), формой компонент, априорными сведениями о числе и разделимости компонент,
режимом наблюдения и объёмом выборки. В последних колонках --- число критических и желательных требований.
{\scriptsize\setlength{\tabcolsep}{3pt}
\begin{longtable}{@{}l >{\raggedright\arraybackslash}p{2.7cm} >{\raggedright\arraybackslash}p{1.5cm} >{\raggedright\arraybackslash}p{1.9cm} >{\raggedright\arraybackslash}p{3.3cm} >{\raggedright\arraybackslash}p{2.5cm} >{\raggedright\arraybackslash}p{1.6cm} r r@{}}
\toprule
ID & Область & Смесь & Компоненты & Априорные сведения & Режим наблюдения & Выборка & Крит. & Жел. \\
\midrule
\endhead
\textbf{UC-01} & Астрофизика & однородная & симметричные & число компонент неизвестно, компоненты разделены & прямые наблюдения & малая & 9 & 13 \\
\textbf{UC-02} & Экология и гидрохимия & однородная & симметричные & число компонент неизвестно, компоненты перекрываются & прямые наблюдения & умеренная & 9 & 12 \\
\textbf{UC-03} & Фармакогенетика и биомедицина & разнородная & асимметричные & число компонент известно, компоненты перекрываются & прямые наблюдения & умеренная & 7 & 10 \\
\textbf{UC-04} & Метрология и историческое материаловедение & однородная & симметричные & число компонент неизвестно, компоненты перекрываются & прямые наблюдения & умеренная & 10 & 10 \\
\textbf{UC-05} & Гидрогеология и геофизика & однородная & симметричные & число компонент известно, компоненты разделены & прямые наблюдения & умеренная & 12 & 10 \\
\textbf{UC-06} & Промышленная надежность & разнородная & положительный носитель & число компонент известно, механизмы конкурируют & цензурированные наблюдения & умеренная & 13 & 11 \\
\textbf{UC-07} & Метеорология и ветроэнергетика & разнородная & положительный носитель & число компонент неизвестно, компоненты перекрываются & прямые наблюдения & большая & 12 & 8 \\
\textbf{UC-08} & Актуарные расчеты и финансовые риски & разнородная & тяжёлый хвост & число компонент известно, семейства разнородны & несбалансированные веса компонент & умеренная & 13 & 8 \\
\textbf{UC-09} & Иммунология и проточная цитометрия & однородная & асимметричные & число компонент неизвестно, компоненты перекрываются & несбалансированные веса компонент & большая & 10 & 10 \\
\textbf{UC-10} & Сетевой трафик и телекоммуникации & разнородная & тяжёлый хвост & число компонент неизвестно, семейства разнородны & сгруппированные данные & очень большая & 16 & 11 \\
\bottomrule
\end{longtable}}
Состав: однородных сценариев --- 5, разнородных --- 5;
по форме компонент: симметричные --- 4, асимметричные --- 2, положительный носитель --- 2, тяжёлый хвост --- 2;
по выборке: умеренная --- 6, большая --- 2, малая --- 1, очень большая --- 1.

\begin{landscape}
\section{Требования сценариев}
Таблица связывает 48 требований, на которые ссылаются сценарии (из 86), со сценариями: К --- критическое, Ж --- желательное. Справа:
число сценариев с критическим и желательным требованием, число инструментов с полным выполнением и статус \texttt{pysatl-mpest}.
{\scriptsize\setlength{\tabcolsep}{3pt}
\begin{longtable}{@{}l >{\raggedright\arraybackslash}p{8.2cm} *{10}{c} r r r l@{}}
\toprule
ID & Требование & 01 & 02 & 03 & 04 & 05 & 06 & 07 & 08 & 09 & 10 & Крит. & Жел. & Инстр. & \texttt{mpest} \\
\midrule
\endhead
\texttt{EST-01} & Каркас и фасад классического EM-алгоритма & \cellcolor{hm!70}К & \cellcolor{hm!70}К & \cellcolor{hm!70}К & \cellcolor{hm!70}К & \cellcolor{hm!70}К & \cellcolor{hm!70}К & \cellcolor{hm!70}К & \cellcolor{hm!70}К & \cellcolor{hm!70}К & \cellcolor{hm!70}К & 10 & 0 & 22 & \cellcolor{part!40}част. \\
\texttt{EST-02} & Каркас и фасад ECM-алгоритма &  &  &  &  &  & \cellcolor{hm!70}К & \cellcolor{hm!70}К &  & \cellcolor{hm!70}К &  & 3 & 0 & 6 & \cellcolor{yes!30}да \\
\texttt{EST-09} & Замкнутые аналитические формулы параметров на M-шаге &  &  &  &  & \cellcolor{hm!70}К &  &  &  &  &  & 1 & 0 & 0 & \cellcolor{part!40}част. \\
\texttt{EST-13} & Явные параметрические границы &  &  & \cellcolor{hm!22}Ж &  &  & \cellcolor{hm!22}Ж &  & \cellcolor{hm!22}Ж &  &  & 0 & 3 & 5 & \textcolor{gray}{нет} \\
\texttt{EST-14} & Предотвращение вырождения масштаба и сингулярностей правдоподобия & \cellcolor{hm!70}К & \cellcolor{hm!70}К & \cellcolor{hm!70}К & \cellcolor{hm!70}К & \cellcolor{hm!70}К & \cellcolor{hm!70}К & \cellcolor{hm!70}К & \cellcolor{hm!70}К & \cellcolor{hm!70}К & \cellcolor{hm!70}К & 10 & 0 & 9 & \cellcolor{yes!30}да \\
\texttt{EST-15} & Сохранение истории сходимости и времени работы & \cellcolor{hm!22}Ж & \cellcolor{hm!22}Ж & \cellcolor{hm!22}Ж & \cellcolor{hm!22}Ж & \cellcolor{hm!22}Ж & \cellcolor{hm!22}Ж & \cellcolor{hm!22}Ж & \cellcolor{hm!22}Ж & \cellcolor{hm!22}Ж & \cellcolor{hm!22}Ж & 0 & 10 & 4 & \cellcolor{yes!30}да \\
\texttt{EST-16} & Оценивание по сгруппированным и интервальным данным &  &  &  &  &  &  &  &  &  & \cellcolor{hm!70}К & 1 & 0 & 3 & \textcolor{gray}{нет} \\
\texttt{EST-17} & Взвешенные наблюдения и веса выборки &  &  &  &  &  &  &  &  &  & \cellcolor{hm!22}Ж & 0 & 1 & 16 & \textcolor{gray}{нет} \\
\texttt{EST-18} & Оценивание по цензурированным и усеченным данным &  &  &  &  &  & \cellcolor{hm!70}К &  &  &  &  & 1 & 0 & 2 & \textcolor{gray}{нет} \\
\texttt{EST-19} & Ускорение сходимости EM-алгоритма (Aitken, SQUAREM, PX-EM) &  &  &  &  &  &  &  &  & \cellcolor{hm!22}Ж & \cellcolor{hm!70}К & 1 & 1 & 0 & \textcolor{gray}{нет} \\
\texttt{EST-20} & Рекурсивный потоковый EM-алгоритм &  &  &  &  &  &  &  &  &  & \cellcolor{hm!22}Ж & 0 & 1 & 1 & \textcolor{gray}{нет} \\
\texttt{EST-22} & Покомпонентный EM-алгоритм (CW-EM, SAGE) &  &  &  & \cellcolor{hm!22}Ж &  &  &  &  &  &  & 0 & 1 & 0 & \textcolor{gray}{нет} \\
\texttt{EST-24} & Байесовское сэмплирование методом расширения данных (Gibbs, MCMC) & \cellcolor{hm!22}Ж & \cellcolor{hm!22}Ж & \cellcolor{hm!22}Ж &  &  & \cellcolor{hm!22}Ж &  &  &  &  & 0 & 4 & 1 & \textcolor{gray}{нет} \\
\texttt{STRUC-01} & Однородные смеси распределений & \cellcolor{hm!70}К & \cellcolor{hm!70}К &  & \cellcolor{hm!70}К & \cellcolor{hm!70}К &  &  &  &  &  & 4 & 0 & 26 & \cellcolor{yes!30}да \\
\texttt{STRUC-02} & Разнородные смеси распределений &  &  &  &  &  & \cellcolor{hm!70}К & \cellcolor{hm!70}К & \cellcolor{hm!70}К &  & \cellcolor{hm!70}К & 4 & 0 & 6 & \cellcolor{yes!30}да \\
\texttt{STRUC-03} & Нормальное семейство распределений & \cellcolor{hm!70}К & \cellcolor{hm!70}К &  & \cellcolor{hm!70}К & \cellcolor{hm!70}К &  &  &  &  &  & 4 & 0 & 25 & \cellcolor{yes!30}да \\
\texttt{STRUC-04} & Экспоненциальное семейство распределений &  &  &  &  &  & \cellcolor{hm!22}Ж &  &  &  & \cellcolor{hm!70}К & 1 & 1 & 11 & \cellcolor{yes!30}да \\
\texttt{STRUC-05} & Асимметричные положительные семейства (Weibull, Gamma, Lognormal) &  &  &  &  &  & \cellcolor{hm!70}К & \cellcolor{hm!70}К & \cellcolor{hm!70}К &  & \cellcolor{hm!70}К & 4 & 0 & 11 & \cellcolor{yes!30}да \\
\texttt{STRUC-06} & Тяжелохвостые семейства (Pareto, Cauchy) &  &  &  &  &  &  &  & \cellcolor{hm!70}К &  & \cellcolor{hm!22}Ж & 1 & 1 & 8 & \cellcolor{yes!30}да \\
\texttt{STRUC-08} & Оценивание распределений с неизвестным порогом &  &  &  &  &  &  &  & \cellcolor{hm!70}К &  &  & 1 & 0 & 1 & \cellcolor{yes!30}да \\
\texttt{STRUC-09} & Информационные критерии выбора числа компонент (AIC, BIC, CAIC) & \cellcolor{hm!70}К & \cellcolor{hm!70}К & \cellcolor{hm!70}К & \cellcolor{hm!70}К & \cellcolor{hm!70}К & \cellcolor{hm!70}К & \cellcolor{hm!70}К & \cellcolor{hm!70}К & \cellcolor{hm!70}К & \cellcolor{hm!70}К & 10 & 0 & 15 & \cellcolor{part!40}част. \\
\texttt{STRUC-10} & Критерий интегрированного правдоподобия классификации (ICL) &  &  &  &  &  &  &  &  & \cellcolor{hm!70}К &  & 1 & 0 & 8 & \textcolor{gray}{нет} \\
\texttt{STRUC-11} & Параметрический бутстреп отношения правдоподобий (LRT) & \cellcolor{hm!22}Ж &  &  &  & \cellcolor{hm!22}Ж &  &  &  &  &  & 0 & 2 & 7 & \textcolor{gray}{нет} \\
\texttt{STRUC-13} & Сравнение и отбор конфигураций семейств &  &  &  &  &  & \cellcolor{hm!70}К & \cellcolor{hm!70}К & \cellcolor{hm!70}К &  & \cellcolor{hm!22}Ж & 3 & 1 & 1 & \textcolor{gray}{нет} \\
\texttt{STRUC-16} & Объединение перекрывающихся компонент по критерию унимодальности (DIP-тест) &  &  &  &  &  &  &  &  & \cellcolor{hm!22}Ж &  & 0 & 1 & 1 & \textcolor{gray}{нет} \\
\texttt{STRUC-21} & Равенство и линейные ограничения параметров компонент & \cellcolor{hm!22}Ж & \cellcolor{hm!22}Ж &  & \cellcolor{hm!70}К & \cellcolor{hm!70}К & \cellcolor{hm!22}Ж &  &  &  &  & 2 & 3 & 8 & \textcolor{gray}{нет} \\
\texttt{STRUC-22} & Асимметричные и тяжелохвостые семейства (Skew-Normal, Student-t) &  &  & \cellcolor{hm!70}К &  &  &  &  &  & \cellcolor{hm!70}К &  & 2 & 0 & 4 & \textcolor{gray}{нет} \\
\texttt{INIT-02} & Рандомизированный мультистарт & \cellcolor{hm!22}Ж & \cellcolor{hm!70}К &  & \cellcolor{hm!22}Ж &  &  &  &  &  &  & 1 & 2 & 9 & \textcolor{gray}{нет} \\
\texttt{INIT-03} & Кластеризация данных &  & \cellcolor{hm!22}Ж &  & \cellcolor{hm!70}К & \cellcolor{hm!22}Ж &  &  &  & \cellcolor{hm!70}К &  & 2 & 2 & 16 & \textcolor{gray}{нет} \\
\texttt{INIT-04} & Квантильное детерминированное разбиение & \cellcolor{hm!70}К & \cellcolor{hm!22}Ж & \cellcolor{hm!22}Ж &  &  & \cellcolor{hm!22}Ж & \cellcolor{hm!70}К & \cellcolor{hm!70}К &  & \cellcolor{hm!70}К & 4 & 3 & 4 & \textcolor{gray}{нет} \\
\texttt{INIT-05} & Инициализация методом моментов по сегментам &  &  &  &  & \cellcolor{hm!70}К &  & \cellcolor{hm!22}Ж &  &  &  & 1 & 1 & 2 & \textcolor{gray}{нет} \\
\texttt{INIT-06} & Случайный выбор наблюдений &  &  &  & \cellcolor{hm!22}Ж &  &  &  &  &  &  & 0 & 1 & 19 & \textcolor{gray}{нет} \\
\texttt{TECH-01} & Композитный модульный пайплайн оценивания & \cellcolor{hm!70}К & \cellcolor{hm!70}К & \cellcolor{hm!70}К & \cellcolor{hm!70}К & \cellcolor{hm!70}К & \cellcolor{hm!70}К & \cellcolor{hm!70}К & \cellcolor{hm!70}К & \cellcolor{hm!70}К & \cellcolor{hm!70}К & 10 & 0 & 5 & \cellcolor{yes!30}да \\
\texttt{TECH-03} & Логарифмическая стабильность вычислений & \cellcolor{hm!70}К & \cellcolor{hm!70}К & \cellcolor{hm!70}К & \cellcolor{hm!70}К & \cellcolor{hm!70}К & \cellcolor{hm!70}К & \cellcolor{hm!70}К & \cellcolor{hm!70}К & \cellcolor{hm!70}К & \cellcolor{hm!70}К & 10 & 0 & 22 & \cellcolor{yes!30}да \\
\texttt{VIS-01} & Декомпозиция плотности смеси & \cellcolor{hm!22}Ж & \cellcolor{hm!22}Ж & \cellcolor{hm!22}Ж & \cellcolor{hm!22}Ж & \cellcolor{hm!22}Ж & \cellcolor{hm!22}Ж & \cellcolor{hm!22}Ж & \cellcolor{hm!22}Ж & \cellcolor{hm!22}Ж & \cellcolor{hm!22}Ж & 0 & 10 & 15 & \textcolor{gray}{нет} \\
\texttt{VIS-02} & Диагностические графики согласия (Q-Q, P-P) & \cellcolor{hm!22}Ж & \cellcolor{hm!22}Ж & \cellcolor{hm!22}Ж & \cellcolor{hm!22}Ж & \cellcolor{hm!22}Ж & \cellcolor{hm!22}Ж & \cellcolor{hm!22}Ж & \cellcolor{hm!22}Ж & \cellcolor{hm!22}Ж & \cellcolor{hm!22}Ж & 0 & 10 & 3 & \textcolor{gray}{нет} \\
\texttt{VIS-03} & Профиль выбора числа компонент (AIC, BIC, ICL) & \cellcolor{hm!22}Ж & \cellcolor{hm!22}Ж &  & \cellcolor{hm!22}Ж & \cellcolor{hm!22}Ж & \cellcolor{hm!22}Ж & \cellcolor{hm!22}Ж & \cellcolor{hm!22}Ж & \cellcolor{hm!22}Ж & \cellcolor{hm!22}Ж & 0 & 9 & 9 & \textcolor{gray}{нет} \\
\texttt{VIS-04} & Визуализация скорости сходимости &  &  &  &  & \cellcolor{hm!22}Ж &  &  &  &  & \cellcolor{hm!22}Ж & 0 & 2 & 5 & \textcolor{gray}{нет} \\
\texttt{VIS-05} & Графики апостериорной неопределенности & \cellcolor{hm!22}Ж & \cellcolor{hm!22}Ж & \cellcolor{hm!22}Ж &  &  &  &  &  & \cellcolor{hm!22}Ж &  & 0 & 4 & 5 & \textcolor{gray}{нет} \\
\texttt{REPR-01} & Поддержка вероятностных функций (PDF, CDF, PPF) & \cellcolor{hm!70}К & \cellcolor{hm!70}К & \cellcolor{hm!70}К & \cellcolor{hm!70}К & \cellcolor{hm!70}К & \cellcolor{hm!70}К & \cellcolor{hm!70}К & \cellcolor{hm!70}К & \cellcolor{hm!70}К & \cellcolor{hm!70}К & 10 & 0 & 5 & \cellcolor{yes!30}да \\
\texttt{REPR-02} & Сэмплирование & \cellcolor{hm!22}Ж & \cellcolor{hm!22}Ж & \cellcolor{hm!22}Ж & \cellcolor{hm!22}Ж & \cellcolor{hm!22}Ж & \cellcolor{hm!22}Ж & \cellcolor{hm!22}Ж & \cellcolor{hm!22}Ж & \cellcolor{hm!22}Ж & \cellcolor{hm!22}Ж & 0 & 10 & 21 & \cellcolor{yes!30}да \\
\texttt{REPR-03} & Аналитический расчет моментов смеси &  &  &  &  & \cellcolor{hm!70}К & \cellcolor{hm!70}К & \cellcolor{hm!70}К & \cellcolor{hm!70}К &  & \cellcolor{hm!70}К & 5 & 0 & 2 & \cellcolor{part!40}част. \\
\texttt{PERF-01} & Параллелизация независимых вычислений &  &  &  &  &  &  &  &  &  & \cellcolor{hm!70}К & 1 & 0 & 10 & \textcolor{gray}{нет} \\
\texttt{PERF-03} & Векторизованные вычисления (NumPy, SIMD) &  &  &  &  &  &  &  &  &  & \cellcolor{hm!70}К & 1 & 0 & 19 & \cellcolor{yes!30}да \\
\texttt{PERF-04} & Контроль аллокаций памяти и потоковые вычисления &  &  &  &  &  &  &  &  &  & \cellcolor{hm!70}К & 1 & 0 & 14 & \textcolor{gray}{нет} \\
\texttt{DIAG-01} & Статистические критерии согласия (KS, AD, Chi2) & \cellcolor{hm!22}Ж & \cellcolor{hm!22}Ж & \cellcolor{hm!22}Ж & \cellcolor{hm!22}Ж & \cellcolor{hm!22}Ж & \cellcolor{hm!22}Ж & \cellcolor{hm!22}Ж & \cellcolor{hm!22}Ж & \cellcolor{hm!22}Ж & \cellcolor{hm!22}Ж & 0 & 10 & 2 & \textcolor{gray}{нет} \\
\texttt{DIAG-02} & Стандартные ошибки и доверительные интервалы оценок & \cellcolor{hm!22}Ж &  & \cellcolor{hm!22}Ж &  & \cellcolor{hm!22}Ж & \cellcolor{hm!70}К & \cellcolor{hm!22}Ж & \cellcolor{hm!22}Ж &  &  & 1 & 5 & 10 & \textcolor{gray}{нет} \\
\texttt{DIAG-04} & Оценка отделимости и перекрытия компонент & \cellcolor{hm!22}Ж & \cellcolor{hm!22}Ж &  & \cellcolor{hm!22}Ж &  &  &  &  & \cellcolor{hm!22}Ж &  & 0 & 4 & 1 & \textcolor{gray}{нет} \\
\bottomrule
\end{longtable}}
Столбцы 01--10 --- номера сценариев UC-01--UC-10.

\end{landscape}
\section{Структура критических требований}
Распределение критических требований сценариев по категориям показывает, какие группы возможностей определяют пригодность инструмента для конкретной задачи.
{\small
\begin{longtable}{@{}l rrrrrrrrrr@{}}
\toprule
Сценарий & \rotatebox{90}{\scriptsize Оценивание} & \rotatebox{90}{\scriptsize Структура смеси} & \rotatebox{90}{\scriptsize Инициализация} & \rotatebox{90}{\scriptsize Архитектура} & \rotatebox{90}{\scriptsize Инфраструктура} & \rotatebox{90}{\scriptsize Визуализация} & \rotatebox{90}{\scriptsize Специальные методы} & \rotatebox{90}{\scriptsize Представление смеси} & \rotatebox{90}{\scriptsize Производительность} & \rotatebox{90}{\scriptsize Диагностика} \\
\midrule
\endhead
UC-01 & 2 & 3 & 1 & 2 &  &  &  & 1 &  &  \\
UC-02 & 2 & 3 & 1 & 2 &  &  &  & 1 &  &  \\
UC-03 & 2 & 2 &  & 2 &  &  &  & 1 &  &  \\
UC-04 & 2 & 4 & 1 & 2 &  &  &  & 1 &  &  \\
UC-05 & 3 & 4 & 1 & 2 &  &  &  & 2 &  &  \\
UC-06 & 4 & 4 &  & 2 &  &  &  & 2 &  & 1 \\
UC-07 & 3 & 4 & 1 & 2 &  &  &  & 2 &  &  \\
UC-08 & 2 & 6 & 1 & 2 &  &  &  & 2 &  &  \\
UC-09 & 3 & 3 & 1 & 2 &  &  &  & 1 &  &  \\
UC-10 & 4 & 4 & 1 & 2 &  &  &  & 2 & 3 &  \\
\bottomrule
\end{longtable}}

\section{Покрытие сценариев инструментами}
\subsection{Покрытие критических требований}
Значение --- доля критических требований сценария, выполненных инструментом (частично выполненные с весом $\tfrac12$).
{\scriptsize\setlength{\tabcolsep}{3pt}
\begin{longtable}{@{}l *{11}{r}@{}}
\toprule
Инструмент & UC-01 & UC-02 & UC-03 & UC-04 & UC-05 & UC-06 & UC-07 & UC-08 & UC-09 & UC-10 & Ср. \\
\midrule
\endhead
brms & \cellcolor{hm!56}72 & \cellcolor{hm!63}83 & \cellcolor{hm!60}79 & \cellcolor{hm!58}75 & \cellcolor{hm!52}67 & \cellcolor{hm!53}69 & \cellcolor{hm!46}58 & \cellcolor{hm!51}65 & \cellcolor{hm!44}55 & \cellcolor{hm!55}72 & \cellcolor{hm!54}70 \\
PySATL mpest & \cellcolor{hm!59}78 & \cellcolor{hm!59}78 & \cellcolor{hm!55}71 & \cellcolor{hm!54}70 & \cellcolor{hm!52}67 & \cellcolor{hm!51}65 & \cellcolor{hm!55}71 & \cellcolor{hm!56}73 & \cellcolor{hm!47}60 & \cellcolor{hm!47}59 & \cellcolor{hm!53}69 \\
mclust (v5+) & \cellcolor{hm!67}89 & \cellcolor{hm!59}78 & \cellcolor{hm!55}71 & \cellcolor{hm!68}90 & \cellcolor{hm!58}75 & \cellcolor{hm!40}50 & \cellcolor{hm!43}54 & \cellcolor{hm!40}50 & \cellcolor{hm!54}70 & \cellcolor{hm!47}59 & \cellcolor{hm!53}69 \\
FlexMix & \cellcolor{hm!56}72 & \cellcolor{hm!63}83 & \cellcolor{hm!55}71 & \cellcolor{hm!58}75 & \cellcolor{hm!52}67 & \cellcolor{hm!53}69 & \cellcolor{hm!49}62 & \cellcolor{hm!43}54 & \cellcolor{hm!54}70 & \cellcolor{hm!42}53 & \cellcolor{hm!52}68 \\
mixR & \cellcolor{hm!59}78 & \cellcolor{hm!59}78 & \cellcolor{hm!50}64 & \cellcolor{hm!58}75 & \cellcolor{hm!55}71 & \cellcolor{hm!40}50 & \cellcolor{hm!43}54 & \cellcolor{hm!40}50 & \cellcolor{hm!44}55 & \cellcolor{hm!49}62 & \cellcolor{hm!50}64 \\
StepMix & \cellcolor{hm!56}72 & \cellcolor{hm!63}83 & \cellcolor{hm!50}64 & \cellcolor{hm!64}85 & \cellcolor{hm!52}67 & \cellcolor{hm!43}54 & \cellcolor{hm!40}50 & \cellcolor{hm!32}38 & \cellcolor{hm!58}75 & \cellcolor{hm!33}41 & \cellcolor{hm!49}63 \\
Distributions.jl & \cellcolor{hm!52}67 & \cellcolor{hm!52}67 & \cellcolor{hm!50}64 & \cellcolor{hm!47}60 & \cellcolor{hm!49}62 & \cellcolor{hm!45}58 & \cellcolor{hm!46}58 & \cellcolor{hm!48}62 & \cellcolor{hm!36}45 & \cellcolor{hm!51}66 & \cellcolor{hm!48}61 \\
Rmixmod (mixmod) & \cellcolor{hm!56}72 & \cellcolor{hm!63}83 & \cellcolor{hm!50}64 & \cellcolor{hm!64}85 & \cellcolor{hm!52}67 & \cellcolor{hm!32}38 & \cellcolor{hm!34}42 & \cellcolor{hm!32}38 & \cellcolor{hm!50}65 & \cellcolor{hm!40}50 & \cellcolor{hm!47}61 \\
rebmix & \cellcolor{hm!56}72 & \cellcolor{hm!56}72 & \cellcolor{hm!50}64 & \cellcolor{hm!50}65 & \cellcolor{hm!46}58 & \cellcolor{hm!43}54 & \cellcolor{hm!40}50 & \cellcolor{hm!37}46 & \cellcolor{hm!44}55 & \cellcolor{hm!49}62 & \cellcolor{hm!47}60 \\
pomegranate & \cellcolor{hm!52}67 & \cellcolor{hm!52}67 & \cellcolor{hm!50}64 & \cellcolor{hm!58}75 & \cellcolor{hm!46}58 & \cellcolor{hm!37}46 & \cellcolor{hm!40}50 & \cellcolor{hm!40}50 & \cellcolor{hm!44}55 & \cellcolor{hm!47}59 & \cellcolor{hm!46}59 \\
scikit-learn (sklearn.mixture) & \cellcolor{hm!56}72 & \cellcolor{hm!63}83 & \cellcolor{hm!50}64 & \cellcolor{hm!61}80 & \cellcolor{hm!49}62 & \cellcolor{hm!29}35 & \cellcolor{hm!31}38 & \cellcolor{hm!29}35 & \cellcolor{hm!44}55 & \cellcolor{hm!38}47 & \cellcolor{hm!45}57 \\
student-mixture & \cellcolor{hm!48}61 & \cellcolor{hm!56}72 & \cellcolor{hm!60}79 & \cellcolor{hm!50}65 & \cellcolor{hm!40}50 & \cellcolor{hm!35}42 & \cellcolor{hm!37}46 & \cellcolor{hm!35}42 & \cellcolor{hm!58}75 & \cellcolor{hm!25}28 & \cellcolor{hm!44}56 \\
bnpy & \cellcolor{hm!52}67 & \cellcolor{hm!59}78 & \cellcolor{hm!50}64 & \cellcolor{hm!54}70 & \cellcolor{hm!43}54 & \cellcolor{hm!27}31 & \cellcolor{hm!28}33 & \cellcolor{hm!27}31 & \cellcolor{hm!44}55 & \cellcolor{hm!38}47 & \cellcolor{hm!42}53 \\
sn (Skew-Normal/t) & \cellcolor{hm!44}56 & \cellcolor{hm!44}56 & \cellcolor{hm!45}57 & \cellcolor{hm!44}55 & \cellcolor{hm!52}67 & \cellcolor{hm!35}42 & \cellcolor{hm!31}38 & \cellcolor{hm!35}42 & \cellcolor{hm!33}40 & \cellcolor{hm!29}34 & \cellcolor{hm!39}49 \\
mixem & \cellcolor{hm!48}61 & \cellcolor{hm!48}61 & \cellcolor{hm!40}50 & \cellcolor{hm!44}55 & \cellcolor{hm!40}50 & \cellcolor{hm!35}42 & \cellcolor{hm!37}46 & \cellcolor{hm!35}42 & \cellcolor{hm!30}35 & \cellcolor{hm!33}41 & \cellcolor{hm!39}48 \\
kasakh Mixture-Models & \cellcolor{hm!48}61 & \cellcolor{hm!48}61 & \cellcolor{hm!40}50 & \cellcolor{hm!50}65 & \cellcolor{hm!37}46 & \cellcolor{hm!24}27 & \cellcolor{hm!25}29 & \cellcolor{hm!29}35 & \cellcolor{hm!44}55 & \cellcolor{hm!27}31 & \cellcolor{hm!37}46 \\
ClusterR & \cellcolor{hm!48}61 & \cellcolor{hm!48}61 & \cellcolor{hm!40}50 & \cellcolor{hm!50}65 & \cellcolor{hm!40}50 & \cellcolor{hm!24}27 & \cellcolor{hm!25}29 & \cellcolor{hm!24}27 & \cellcolor{hm!36}45 & \cellcolor{hm!33}41 & \cellcolor{hm!37}46 \\
GaussianMixtures.jl & \cellcolor{hm!48}61 & \cellcolor{hm!48}61 & \cellcolor{hm!40}50 & \cellcolor{hm!50}65 & \cellcolor{hm!40}50 & \cellcolor{hm!24}27 & \cellcolor{hm!25}29 & \cellcolor{hm!24}27 & \cellcolor{hm!36}45 & \cellcolor{hm!33}41 & \cellcolor{hm!37}46 \\
EMCluster & \cellcolor{hm!44}56 & \cellcolor{hm!52}67 & \cellcolor{hm!35}43 & \cellcolor{hm!54}70 & \cellcolor{hm!43}54 & \cellcolor{hm!27}31 & \cellcolor{hm!22}25 & \cellcolor{hm!21}23 & \cellcolor{hm!40}50 & \cellcolor{hm!27}31 & \cellcolor{hm!36}45 \\
teigen & \cellcolor{hm!36}44 & \cellcolor{hm!36}44 & \cellcolor{hm!45}57 & \cellcolor{hm!47}60 & \cellcolor{hm!37}46 & \cellcolor{hm!27}31 & \cellcolor{hm!28}33 & \cellcolor{hm!27}31 & \cellcolor{hm!54}70 & \cellcolor{hm!27}31 & \cellcolor{hm!36}45 \\
mixtools & \cellcolor{hm!48}61 & \cellcolor{hm!40}50 & \cellcolor{hm!30}36 & \cellcolor{hm!40}50 & \cellcolor{hm!37}46 & \cellcolor{hm!35}42 & \cellcolor{hm!37}46 & \cellcolor{hm!32}38 & \cellcolor{hm!33}40 & \cellcolor{hm!31}38 & \cellcolor{hm!36}45 \\
mlpack (mlpack::gmm) & \cellcolor{hm!44}56 & \cellcolor{hm!48}61 & \cellcolor{hm!35}43 & \cellcolor{hm!47}60 & \cellcolor{hm!37}46 & \cellcolor{hm!21}23 & \cellcolor{hm!22}25 & \cellcolor{hm!21}23 & \cellcolor{hm!33}40 & \cellcolor{hm!31}38 & \cellcolor{hm!34}41 \\
mixdist & \cellcolor{hm!44}56 & \cellcolor{hm!36}44 & \cellcolor{hm!25}29 & \cellcolor{hm!40}50 & \cellcolor{hm!37}46 & \cellcolor{hm!29}35 & \cellcolor{hm!31}38 & \cellcolor{hm!32}38 & \cellcolor{hm!19}20 & \cellcolor{hm!40}50 & \cellcolor{hm!33}40 \\
fpc & \cellcolor{hm!44}56 & \cellcolor{hm!44}56 & \cellcolor{hm!35}43 & \cellcolor{hm!47}60 & \cellcolor{hm!37}46 & \cellcolor{hm!24}27 & \cellcolor{hm!22}25 & \cellcolor{hm!21}23 & \cellcolor{hm!33}40 & \cellcolor{hm!20}22 & \cellcolor{hm!33}40 \\
CAMAN & \cellcolor{hm!44}56 & \cellcolor{hm!44}56 & \cellcolor{hm!35}43 & \cellcolor{hm!40}50 & \cellcolor{hm!37}46 & \cellcolor{hm!27}31 & \cellcolor{hm!22}25 & \cellcolor{hm!21}23 & \cellcolor{hm!26}30 & \cellcolor{hm!29}34 & \cellcolor{hm!33}39 \\
Armadillo C++ & \cellcolor{hm!40}50 & \cellcolor{hm!40}50 & \cellcolor{hm!30}36 & \cellcolor{hm!44}55 & \cellcolor{hm!34}42 & \cellcolor{hm!18}19 & \cellcolor{hm!20}21 & \cellcolor{hm!18}19 & \cellcolor{hm!30}35 & \cellcolor{hm!29}34 & \cellcolor{hm!30}36 \\
distfit & \cellcolor{hm!28}33 & \cellcolor{hm!28}33 & \cellcolor{hm!30}36 & \cellcolor{hm!26}30 & \cellcolor{hm!25}29 & \cellcolor{hm!29}35 & \cellcolor{hm!31}38 & \cellcolor{hm!37}46 & \cellcolor{hm!22}25 & \cellcolor{hm!31}38 & \cellcolor{hm!29}34 \\
\bottomrule
\end{longtable}}
Лучший инструмент для сценария в среднем закрывает 78\,\% его критических требований; полное покрытие отсутствует у всех инструментов во всех сценариях.

\subsection{Сколько инструментов нужно, чтобы закрыть сценарий}
Если ни один инструмент не закрывает критические требования сценария полностью, возникает вопрос, сколько инструментов пришлось бы совмещать.
Жадный подбор: на каждом шаге выбирается инструмент, полностью выполняющий наибольшее число ещё не закрытых критических требований. Результат --- верхняя оценка
размера набора и показатель разрозненности возможностей.
{\footnotesize
\begin{longtable}{@{}l r >{\raggedright\arraybackslash}p{8.8cm} >{\raggedright\arraybackslash}p{3.6cm}@{}}
\toprule
Сценарий & Набор & Инструменты (сколько критических требований добавляет) & Не закрыто \\
\midrule
\endhead
UC-01 & 2 & mixR (7), PySATL mpest (2) & --- \\
UC-02 & 2 & StepMix (7), PySATL mpest (2) & --- \\
UC-03 & 2 & student-mixture (5), PySATL mpest (2) & --- \\
UC-04 & 2 & StepMix (8), PySATL mpest (2) & --- \\
UC-05 & 3 & StepMix (7), sn (Skew-Normal/t) (3), PySATL mpest (1) & EST-09 \\
UC-06 & 5 & FlexMix (8), PySATL mpest (2), sn (Skew-Normal/t) (1), mixR (1), distfit (1) & --- \\
UC-07 & 4 & PySATL mpest (7), mixR (3), sn (Skew-Normal/t) (1), distfit (1) & --- \\
UC-08 & 4 & PySATL mpest (8), mixR (3), sn (Skew-Normal/t) (1), distfit (1) & --- \\
UC-09 & 3 & student-mixture (7), PySATL mpest (2), teigen (1) & --- \\
UC-10 & 3 & brms (10), mixR (4), sn (Skew-Normal/t) (1) & EST-19 \\
\bottomrule
\end{longtable}}
Среднее число инструментов в наборе: 3.0. Сценарии, у которых после подбора остаются полностью не выполненные никем требования:
UC-05, UC-10.
Это значит, что для сценария либо нужно совмещать несколько пакетов (и вручную согласовывать форматы моделей), либо дописывать недостающее.

\part{Разнородные смеси}
\begin{landscape}
\section{Полное сопоставление средств для разнородных смесей}
Разнородные смеси, где компоненты принадлежат разным семействам, --- основной сценарий библиотеки. Полностью допускают их 5 инструментов кроме \texttt{pysatl-mpest}:
\texttt{pomegranate}, \texttt{mixem}, \texttt{FlexMix}, \texttt{brms}, \texttt{Distributions.jl}; частично --- \texttt{StepMix}, \texttt{Rmixmod (mixmod)}, \texttt{mixdist}. Таблица сопоставляет все 86 требований для этих 9 средств.
{\scriptsize\setlength{\tabcolsep}{3pt}
\begin{longtable}{@{}l >{\raggedright\arraybackslash}p{8.0cm} *{9}{c}@{}}
\toprule
ID & Требование & \rotatebox{90}{\scriptsize pomegranate} & \rotatebox{90}{\scriptsize mixem} & \rotatebox{90}{\scriptsize FlexMix} & \rotatebox{90}{\scriptsize brms} & \rotatebox{90}{\scriptsize Distributions.jl} & \rotatebox{90}{\scriptsize PySATL mpest} & \rotatebox{90}{\scriptsize StepMix} & \rotatebox{90}{\scriptsize Rmixmod (mixmod)} & \rotatebox{90}{\scriptsize mixdist} \\
\midrule
\endhead
\midrule \multicolumn{11}{@{}l}{\textbf{Оценивание}} \\
\texttt{EST-01} & Каркас и фасад классического EM-алгоритма & \cellcolor{yes!30}да & \cellcolor{yes!30}да & \cellcolor{yes!30}да & \textcolor{gray}{нет} & \cellcolor{yes!30}да & \cellcolor{part!40}част. & \cellcolor{yes!30}да & \cellcolor{yes!30}да & \cellcolor{yes!30}да \\
\texttt{EST-02} & Каркас и фасад ECM-алгоритма & \textcolor{gray}{нет} & \textcolor{gray}{нет} & \cellcolor{yes!30}да & \textcolor{gray}{нет} & \textcolor{gray}{нет} & \cellcolor{yes!30}да & \cellcolor{yes!30}да & \textcolor{gray}{нет} & \textcolor{gray}{нет} \\
\texttt{EST-03} & Каркас и фасад ECME-алгоритма & \textcolor{gray}{нет} & \textcolor{gray}{нет} & \textcolor{gray}{нет} & \textcolor{gray}{нет} & \textcolor{gray}{нет} & \cellcolor{yes!30}да & \textcolor{gray}{нет} & \textcolor{gray}{нет} & \textcolor{gray}{нет} \\
\texttt{EST-04} & Direct MLE & \textcolor{gray}{нет} & \textcolor{gray}{нет} & \textcolor{gray}{нет} & \cellcolor{yes!30}да & \cellcolor{part!40}част. & \textcolor{gray}{нет} & \textcolor{gray}{нет} & \textcolor{gray}{нет} & \cellcolor{yes!30}да \\
\texttt{EST-05} & Стратегия суррогатной Q-функции на M-шаге & \cellcolor{yes!30}да & \cellcolor{yes!30}да & \cellcolor{yes!30}да & \textcolor{gray}{нет} & \cellcolor{yes!30}да & \cellcolor{yes!30}да & \cellcolor{yes!30}да & \cellcolor{yes!30}да & \cellcolor{yes!30}да \\
\texttt{EST-06} & Стратегия наблюдаемого правдоподобия на M-шаге & \textcolor{gray}{нет} & \textcolor{gray}{нет} & \textcolor{gray}{нет} & \textcolor{gray}{нет} & \textcolor{gray}{нет} & \cellcolor{yes!30}да & \textcolor{gray}{нет} & \textcolor{gray}{нет} & \textcolor{gray}{нет} \\
\texttt{EST-07} & Стратегия взвешенного метода моментов на M-шаге & \textcolor{gray}{нет} & \textcolor{gray}{нет} & \textcolor{gray}{нет} & \textcolor{gray}{нет} & \textcolor{gray}{нет} & \cellcolor{part!40}част. & \textcolor{gray}{нет} & \textcolor{gray}{нет} & \textcolor{gray}{нет} \\
\texttt{EST-08} & Стратегия метода L-моментов на M-шаге & \textcolor{gray}{нет} & \textcolor{gray}{нет} & \textcolor{gray}{нет} & \textcolor{gray}{нет} & \textcolor{gray}{нет} & \cellcolor{part!40}част. & \textcolor{gray}{нет} & \textcolor{gray}{нет} & \textcolor{gray}{нет} \\
\texttt{EST-09} & Замкнутые аналитические формулы параметров на M-шаге & \cellcolor{part!40}част. & \cellcolor{part!40}част. & \cellcolor{part!40}част. & \textcolor{gray}{нет} & \cellcolor{part!40}част. & \cellcolor{part!40}част. & \cellcolor{part!40}част. & \cellcolor{part!40}част. & \cellcolor{part!40}част. \\
\texttt{EST-10} & Стохастический и классификационный EM-алгоритм & \cellcolor{part!40}част. & \textcolor{gray}{нет} & \cellcolor{yes!30}да & \textcolor{gray}{нет} & \textcolor{gray}{нет} & \textcolor{gray}{нет} & \textcolor{gray}{нет} & \cellcolor{yes!30}да & \textcolor{gray}{нет} \\
\texttt{EST-11} & Алгоритм Rough-Enhanced-Bayes & \textcolor{gray}{нет} & \textcolor{gray}{нет} & \textcolor{gray}{нет} & \textcolor{gray}{нет} & \textcolor{gray}{нет} & \textcolor{gray}{нет} & \textcolor{gray}{нет} & \textcolor{gray}{нет} & \textcolor{gray}{нет} \\
\texttt{EST-12} & Фиксация подмножеств параметров & \cellcolor{yes!30}да & \textcolor{gray}{нет} & \cellcolor{yes!30}да & \cellcolor{yes!30}да & \textcolor{gray}{нет} & \cellcolor{yes!30}да & \cellcolor{yes!30}да & \cellcolor{part!40}част. & \cellcolor{yes!30}да \\
\texttt{EST-13} & Явные параметрические границы & \cellcolor{part!40}част. & \textcolor{gray}{нет} & \textcolor{gray}{нет} & \cellcolor{yes!30}да & \textcolor{gray}{нет} & \textcolor{gray}{нет} & \textcolor{gray}{нет} & \textcolor{gray}{нет} & \cellcolor{yes!30}да \\
\texttt{EST-14} & Предотвращение вырождения масштаба и сингулярностей правдоподобия & \cellcolor{yes!30}да & \cellcolor{part!40}част. & \cellcolor{part!40}част. & \cellcolor{yes!30}да & \textcolor{gray}{нет} & \cellcolor{yes!30}да & \cellcolor{yes!30}да & \cellcolor{part!40}част. & \textcolor{gray}{нет} \\
\texttt{EST-15} & Сохранение истории сходимости и времени работы & \cellcolor{part!40}част. & \cellcolor{part!40}част. & \cellcolor{part!40}част. & \cellcolor{yes!30}да & \textcolor{gray}{нет} & \cellcolor{yes!30}да & \cellcolor{part!40}част. & \cellcolor{part!40}част. & \cellcolor{part!40}част. \\
\texttt{EST-16} & Оценивание по сгруппированным и интервальным данным & \textcolor{gray}{нет} & \textcolor{gray}{нет} & \textcolor{gray}{нет} & \cellcolor{part!40}част. & \textcolor{gray}{нет} & \textcolor{gray}{нет} & \textcolor{gray}{нет} & \textcolor{gray}{нет} & \cellcolor{yes!30}да \\
\texttt{EST-17} & Взвешенные наблюдения и веса выборки & \cellcolor{yes!30}да & \textcolor{gray}{нет} & \cellcolor{yes!30}да & \cellcolor{yes!30}да & \cellcolor{yes!30}да & \textcolor{gray}{нет} & \cellcolor{yes!30}да & \cellcolor{yes!30}да & \cellcolor{yes!30}да \\
\texttt{EST-18} & Оценивание по цензурированным и усеченным данным & \textcolor{gray}{нет} & \textcolor{gray}{нет} & \cellcolor{part!40}част. & \cellcolor{yes!30}да & \cellcolor{part!40}част. & \textcolor{gray}{нет} & \textcolor{gray}{нет} & \textcolor{gray}{нет} & \textcolor{gray}{нет} \\
\texttt{EST-19} & Ускорение сходимости EM-алгоритма (Aitken, SQUAREM, PX-EM) & \textcolor{gray}{нет} & \textcolor{gray}{нет} & \textcolor{gray}{нет} & \textcolor{gray}{нет} & \textcolor{gray}{нет} & \textcolor{gray}{нет} & \textcolor{gray}{нет} & \textcolor{gray}{нет} & \cellcolor{part!40}част. \\
\texttt{EST-20} & Рекурсивный потоковый EM-алгоритм & \cellcolor{part!40}част. & \textcolor{gray}{нет} & \textcolor{gray}{нет} & \textcolor{gray}{нет} & \textcolor{gray}{нет} & \textcolor{gray}{нет} & \textcolor{gray}{нет} & \textcolor{gray}{нет} & \textcolor{gray}{нет} \\
\texttt{EST-21} & Оценивание методом минимального расстояния (Hellinger, Wasserstein) & \textcolor{gray}{нет} & \textcolor{gray}{нет} & \textcolor{gray}{нет} & \textcolor{gray}{нет} & \textcolor{gray}{нет} & \textcolor{gray}{нет} & \textcolor{gray}{нет} & \textcolor{gray}{нет} & \cellcolor{part!40}част. \\
\texttt{EST-22} & Покомпонентный EM-алгоритм (CW-EM, SAGE) & \textcolor{gray}{нет} & \textcolor{gray}{нет} & \textcolor{gray}{нет} & \textcolor{gray}{нет} & \textcolor{gray}{нет} & \textcolor{gray}{нет} & \textcolor{gray}{нет} & \textcolor{gray}{нет} & \textcolor{gray}{нет} \\
\texttt{EST-23} & Робастное оценивание по усеченному правдоподобию (TCLUST, C-шаг) & \textcolor{gray}{нет} & \textcolor{gray}{нет} & \textcolor{gray}{нет} & \textcolor{gray}{нет} & \textcolor{gray}{нет} & \textcolor{gray}{нет} & \textcolor{gray}{нет} & \textcolor{gray}{нет} & \textcolor{gray}{нет} \\
\texttt{EST-24} & Байесовское сэмплирование методом расширения данных (Gibbs, MCMC) & \textcolor{gray}{нет} & \textcolor{gray}{нет} & \textcolor{gray}{нет} & \cellcolor{yes!30}да & \textcolor{gray}{нет} & \textcolor{gray}{нет} & \textcolor{gray}{нет} & \textcolor{gray}{нет} & \textcolor{gray}{нет} \\
\midrule \multicolumn{11}{@{}l}{\textbf{Структура смеси}} \\
\texttt{STRUC-01} & Однородные смеси распределений & \cellcolor{yes!30}да & \cellcolor{yes!30}да & \cellcolor{yes!30}да & \cellcolor{yes!30}да & \cellcolor{yes!30}да & \cellcolor{yes!30}да & \cellcolor{yes!30}да & \cellcolor{yes!30}да & \cellcolor{yes!30}да \\
\texttt{STRUC-02} & Разнородные смеси распределений & \cellcolor{yes!30}да & \cellcolor{yes!30}да & \cellcolor{yes!30}да & \cellcolor{yes!30}да & \cellcolor{yes!30}да & \cellcolor{yes!30}да & \cellcolor{part!40}част. & \cellcolor{part!40}част. & \cellcolor{part!40}част. \\
\texttt{STRUC-03} & Нормальное семейство распределений & \cellcolor{yes!30}да & \cellcolor{yes!30}да & \cellcolor{yes!30}да & \cellcolor{yes!30}да & \cellcolor{yes!30}да & \cellcolor{yes!30}да & \cellcolor{yes!30}да & \cellcolor{yes!30}да & \cellcolor{yes!30}да \\
\texttt{STRUC-04} & Экспоненциальное семейство распределений & \cellcolor{yes!30}да & \cellcolor{yes!30}да & \cellcolor{yes!30}да & \cellcolor{yes!30}да & \cellcolor{yes!30}да & \cellcolor{yes!30}да & \textcolor{gray}{нет} & \textcolor{gray}{нет} & \cellcolor{yes!30}да \\
\texttt{STRUC-05} & Асимметричные положительные семейства (Weibull, Gamma, Lognormal) & \cellcolor{yes!30}да & \cellcolor{yes!30}да & \cellcolor{yes!30}да & \cellcolor{yes!30}да & \cellcolor{yes!30}да & \cellcolor{yes!30}да & \textcolor{gray}{нет} & \textcolor{gray}{нет} & \cellcolor{yes!30}да \\
\texttt{STRUC-06} & Тяжелохвостые семейства (Pareto, Cauchy) & \cellcolor{part!40}част. & \textcolor{gray}{нет} & \cellcolor{part!40}част. & \cellcolor{yes!30}да & \cellcolor{yes!30}да & \cellcolor{yes!30}да & \textcolor{gray}{нет} & \textcolor{gray}{нет} & \textcolor{gray}{нет} \\
\texttt{STRUC-07} & Семейства на ограниченном носителе (Beta, Uniform) & \cellcolor{yes!30}да & \textcolor{gray}{нет} & \cellcolor{yes!30}да & \cellcolor{yes!30}да & \cellcolor{yes!30}да & \cellcolor{yes!30}да & \textcolor{gray}{нет} & \textcolor{gray}{нет} & \textcolor{gray}{нет} \\
\texttt{STRUC-08} & Оценивание распределений с неизвестным порогом & \textcolor{gray}{нет} & \textcolor{gray}{нет} & \textcolor{gray}{нет} & \cellcolor{part!40}част. & \textcolor{gray}{нет} & \cellcolor{yes!30}да & \textcolor{gray}{нет} & \textcolor{gray}{нет} & \cellcolor{part!40}част. \\
\texttt{STRUC-09} & Информационные критерии выбора числа компонент (AIC, BIC, CAIC) & \textcolor{gray}{нет} & \textcolor{gray}{нет} & \cellcolor{yes!30}да & \cellcolor{yes!30}да & \cellcolor{part!40}част. & \cellcolor{part!40}част. & \cellcolor{yes!30}да & \cellcolor{yes!30}да & \cellcolor{part!40}част. \\
\texttt{STRUC-10} & Критерий интегрированного правдоподобия классификации (ICL) & \textcolor{gray}{нет} & \textcolor{gray}{нет} & \cellcolor{yes!30}да & \textcolor{gray}{нет} & \textcolor{gray}{нет} & \textcolor{gray}{нет} & \cellcolor{yes!30}да & \cellcolor{yes!30}да & \textcolor{gray}{нет} \\
\texttt{STRUC-11} & Параметрический бутстреп отношения правдоподобий (LRT) & \textcolor{gray}{нет} & \textcolor{gray}{нет} & \cellcolor{part!40}част. & \textcolor{gray}{нет} & \textcolor{gray}{нет} & \textcolor{gray}{нет} & \cellcolor{yes!30}да & \textcolor{gray}{нет} & \textcolor{gray}{нет} \\
\texttt{STRUC-12} & Динамическое удаление вырожденных компонент & \cellcolor{part!40}част. & \textcolor{gray}{нет} & \cellcolor{yes!30}да & \textcolor{gray}{нет} & \textcolor{gray}{нет} & \cellcolor{yes!30}да & \textcolor{gray}{нет} & \textcolor{gray}{нет} & \textcolor{gray}{нет} \\
\texttt{STRUC-13} & Сравнение и отбор конфигураций семейств & \textcolor{gray}{нет} & \textcolor{gray}{нет} & \textcolor{gray}{нет} & \textcolor{gray}{нет} & \textcolor{gray}{нет} & \textcolor{gray}{нет} & \textcolor{gray}{нет} & \textcolor{gray}{нет} & \textcolor{gray}{нет} \\
\texttt{STRUC-14} & Жадное инкрементальное обучение порядка смеси & \textcolor{gray}{нет} & \textcolor{gray}{нет} & \textcolor{gray}{нет} & \textcolor{gray}{нет} & \textcolor{gray}{нет} & \textcolor{gray}{нет} & \textcolor{gray}{нет} & \textcolor{gray}{нет} & \textcolor{gray}{нет} \\
\texttt{STRUC-15} & Динамические переходы слияния и расщепления (SMEM) & \textcolor{gray}{нет} & \textcolor{gray}{нет} & \textcolor{gray}{нет} & \textcolor{gray}{нет} & \textcolor{gray}{нет} & \textcolor{gray}{нет} & \textcolor{gray}{нет} & \textcolor{gray}{нет} & \textcolor{gray}{нет} \\
\texttt{STRUC-16} & Объединение перекрывающихся компонент по критерию унимодальности (DIP-тест) & \textcolor{gray}{нет} & \textcolor{gray}{нет} & \textcolor{gray}{нет} & \textcolor{gray}{нет} & \textcolor{gray}{нет} & \textcolor{gray}{нет} & \textcolor{gray}{нет} & \textcolor{gray}{нет} & \textcolor{gray}{нет} \\
\texttt{STRUC-17} & Непараметрическое градиентное нахождение числа компонент (NPMLE) & \textcolor{gray}{нет} & \textcolor{gray}{нет} & \textcolor{gray}{нет} & \textcolor{gray}{нет} & \textcolor{gray}{нет} & \textcolor{gray}{нет} & \textcolor{gray}{нет} & \textcolor{gray}{нет} & \textcolor{gray}{нет} \\
\texttt{STRUC-18} & Предсказание порядка смеси через мета-обучение (Meta-Learning) & \textcolor{gray}{нет} & \textcolor{gray}{нет} & \textcolor{gray}{нет} & \textcolor{gray}{нет} & \textcolor{gray}{нет} & \textcolor{gray}{нет} & \textcolor{gray}{нет} & \textcolor{gray}{нет} & \textcolor{gray}{нет} \\
\texttt{STRUC-19} & Классификация параметрических семейств с помощью нейросетей & \textcolor{gray}{нет} & \textcolor{gray}{нет} & \textcolor{gray}{нет} & \textcolor{gray}{нет} & \textcolor{gray}{нет} & \textcolor{gray}{нет} & \textcolor{gray}{нет} & \textcolor{gray}{нет} & \textcolor{gray}{нет} \\
\texttt{STRUC-20} & Амортизированный нейросетевой вывод параметров смеси (Deep Sets) & \textcolor{gray}{нет} & \textcolor{gray}{нет} & \textcolor{gray}{нет} & \textcolor{gray}{нет} & \textcolor{gray}{нет} & \textcolor{gray}{нет} & \textcolor{gray}{нет} & \textcolor{gray}{нет} & \textcolor{gray}{нет} \\
\texttt{STRUC-21} & Равенство и линейные ограничения параметров компонент & \cellcolor{part!40}част. & \textcolor{gray}{нет} & \cellcolor{yes!30}да & \cellcolor{yes!30}да & \textcolor{gray}{нет} & \textcolor{gray}{нет} & \cellcolor{yes!30}да & \cellcolor{yes!30}да & \cellcolor{yes!30}да \\
\texttt{STRUC-22} & Асимметричные и тяжелохвостые семейства (Skew-Normal, Student-t) & \cellcolor{part!40}част. & \textcolor{gray}{нет} & \cellcolor{part!40}част. & \cellcolor{yes!30}да & \cellcolor{part!40}част. & \textcolor{gray}{нет} & \textcolor{gray}{нет} & \textcolor{gray}{нет} & \textcolor{gray}{нет} \\
\texttt{STRUC-23} & Вариационный байесовский вывод (VBEM) & \textcolor{gray}{нет} & \textcolor{gray}{нет} & \textcolor{gray}{нет} & \cellcolor{part!40}част. & \textcolor{gray}{нет} & \textcolor{gray}{нет} & \textcolor{gray}{нет} & \textcolor{gray}{нет} & \textcolor{gray}{нет} \\
\midrule \multicolumn{11}{@{}l}{\textbf{Инициализация}} \\
\texttt{INIT-01} & Пользовательская явная инициализация & \cellcolor{yes!30}да & \cellcolor{yes!30}да & \cellcolor{yes!30}да & \cellcolor{yes!30}да & \cellcolor{yes!30}да & \cellcolor{yes!30}да & \cellcolor{part!40}част. & \cellcolor{part!40}част. & \cellcolor{yes!30}да \\
\texttt{INIT-02} & Рандомизированный мультистарт & \textcolor{gray}{нет} & \textcolor{gray}{нет} & \cellcolor{yes!30}да & \cellcolor{yes!30}да & \textcolor{gray}{нет} & \textcolor{gray}{нет} & \cellcolor{yes!30}да & \cellcolor{yes!30}да & \textcolor{gray}{нет} \\
\texttt{INIT-03} & Кластеризация данных & \cellcolor{yes!30}да & \textcolor{gray}{нет} & \textcolor{gray}{нет} & \textcolor{gray}{нет} & \textcolor{gray}{нет} & \textcolor{gray}{нет} & \cellcolor{yes!30}да & \cellcolor{yes!30}да & \textcolor{gray}{нет} \\
\texttt{INIT-04} & Квантильное детерминированное разбиение & \textcolor{gray}{нет} & \textcolor{gray}{нет} & \textcolor{gray}{нет} & \textcolor{gray}{нет} & \textcolor{gray}{нет} & \textcolor{gray}{нет} & \textcolor{gray}{нет} & \textcolor{gray}{нет} & \cellcolor{yes!30}да \\
\texttt{INIT-05} & Инициализация методом моментов по сегментам & \textcolor{gray}{нет} & \textcolor{gray}{нет} & \textcolor{gray}{нет} & \textcolor{gray}{нет} & \textcolor{gray}{нет} & \textcolor{gray}{нет} & \textcolor{gray}{нет} & \textcolor{gray}{нет} & \textcolor{gray}{нет} \\
\texttt{INIT-06} & Случайный выбор наблюдений & \cellcolor{yes!30}да & \textcolor{gray}{нет} & \cellcolor{yes!30}да & \cellcolor{yes!30}да & \textcolor{gray}{нет} & \textcolor{gray}{нет} & \cellcolor{yes!30}да & \cellcolor{yes!30}да & \textcolor{gray}{нет} \\
\midrule \multicolumn{11}{@{}l}{\textbf{Архитектура}} \\
\texttt{TECH-01} & Композитный модульный пайплайн оценивания & \cellcolor{part!40}част. & \cellcolor{part!40}част. & \cellcolor{yes!30}да & \cellcolor{yes!30}да & \cellcolor{part!40}част. & \cellcolor{yes!30}да & \cellcolor{part!40}част. & \cellcolor{yes!30}да & \textcolor{gray}{нет} \\
\texttt{TECH-02} & Расширяемость новыми распределениями & \cellcolor{yes!30}да & \cellcolor{yes!30}да & \cellcolor{yes!30}да & \cellcolor{yes!30}да & \cellcolor{yes!30}да & \cellcolor{yes!30}да & \cellcolor{part!40}част. & \textcolor{gray}{нет} & \textcolor{gray}{нет} \\
\texttt{TECH-03} & Логарифмическая стабильность вычислений & \cellcolor{yes!30}да & \cellcolor{yes!30}да & \cellcolor{yes!30}да & \cellcolor{yes!30}да & \cellcolor{yes!30}да & \cellcolor{yes!30}да & \cellcolor{yes!30}да & \cellcolor{yes!30}да & \cellcolor{part!40}част. \\
\texttt{TECH-04} & Иммутабельность структур параметров модели & \textcolor{gray}{нет} & \cellcolor{yes!30}да & \cellcolor{yes!30}да & \cellcolor{yes!30}да & \cellcolor{yes!30}да & \cellcolor{yes!30}да & \textcolor{gray}{нет} & \cellcolor{yes!30}да & \cellcolor{yes!30}да \\
\texttt{TECH-05} & scikit-learn совместимый интерфейс & \cellcolor{part!40}част. & \textcolor{gray}{нет} & \textcolor{gray}{нет} & \cellcolor{yes!30}да & \textcolor{gray}{нет} & \textcolor{gray}{нет} & \cellcolor{yes!30}да & \cellcolor{yes!30}да & \textcolor{gray}{нет} \\
\midrule \multicolumn{11}{@{}l}{\textbf{Инфраструктура}} \\
\texttt{INFRA-01} & Декларативное управление зависимостями и сборкой & \cellcolor{part!40}част. & \textcolor{gray}{нет} & \cellcolor{yes!30}да & \cellcolor{yes!30}да & \cellcolor{yes!30}да & \cellcolor{yes!30}да & \cellcolor{yes!30}да & \cellcolor{yes!30}да & \cellcolor{yes!30}да \\
\texttt{INFRA-02} & Юнит- и интеграционное тестирование с высоким покрытием & \cellcolor{yes!30}да & \cellcolor{yes!30}да & \cellcolor{yes!30}да & \cellcolor{yes!30}да & \cellcolor{yes!30}да & \cellcolor{yes!30}да & \cellcolor{yes!30}да & \cellcolor{yes!30}да & \cellcolor{part!40}част. \\
\texttt{INFRA-03} & Property-based тестирование математических инвариантов & \textcolor{gray}{нет} & \textcolor{gray}{нет} & \textcolor{gray}{нет} & \textcolor{gray}{нет} & \cellcolor{part!40}част. & \cellcolor{yes!30}да & \textcolor{gray}{нет} & \textcolor{gray}{нет} & \textcolor{gray}{нет} \\
\texttt{INFRA-04} & Автоматизированный статический анализ и линтинг & \textcolor{gray}{нет} & \textcolor{gray}{нет} & \cellcolor{part!40}част. & \cellcolor{part!40}част. & \cellcolor{yes!30}да & \cellcolor{yes!30}да & \cellcolor{part!40}част. & \cellcolor{part!40}част. & \cellcolor{part!40}част. \\
\texttt{INFRA-05} & Автоматизированная непрерывная интеграция & \cellcolor{yes!30}да & \cellcolor{yes!30}да & \cellcolor{yes!30}да & \cellcolor{yes!30}да & \cellcolor{yes!30}да & \cellcolor{yes!30}да & \cellcolor{yes!30}да & \cellcolor{yes!30}да & \cellcolor{yes!30}да \\
\texttt{INFRA-06} & Автоматическая генерация документации с математикой & \cellcolor{yes!30}да & \cellcolor{yes!30}да & \cellcolor{yes!30}да & \cellcolor{yes!30}да & \cellcolor{yes!30}да & \cellcolor{yes!30}да & \cellcolor{yes!30}да & \cellcolor{yes!30}да & \cellcolor{yes!30}да \\
\midrule \multicolumn{11}{@{}l}{\textbf{Визуализация}} \\
\texttt{VIS-01} & Декомпозиция плотности смеси & \textcolor{gray}{нет} & \textcolor{gray}{нет} & \cellcolor{yes!30}да & \cellcolor{yes!30}да & \cellcolor{part!40}част. & \textcolor{gray}{нет} & \cellcolor{yes!30}да & \cellcolor{yes!30}да & \cellcolor{yes!30}да \\
\texttt{VIS-02} & Диагностические графики согласия (Q-Q, P-P) & \textcolor{gray}{нет} & \textcolor{gray}{нет} & \cellcolor{part!40}част. & \cellcolor{yes!30}да & \textcolor{gray}{нет} & \textcolor{gray}{нет} & \textcolor{gray}{нет} & \textcolor{gray}{нет} & \textcolor{gray}{нет} \\
\texttt{VIS-03} & Профиль выбора числа компонент (AIC, BIC, ICL) & \textcolor{gray}{нет} & \textcolor{gray}{нет} & \cellcolor{yes!30}да & \cellcolor{part!40}част. & \textcolor{gray}{нет} & \textcolor{gray}{нет} & \cellcolor{yes!30}да & \cellcolor{yes!30}да & \textcolor{gray}{нет} \\
\texttt{VIS-04} & Визуализация скорости сходимости & \textcolor{gray}{нет} & \textcolor{gray}{нет} & \textcolor{gray}{нет} & \cellcolor{yes!30}да & \textcolor{gray}{нет} & \textcolor{gray}{нет} & \textcolor{gray}{нет} & \cellcolor{yes!30}да & \textcolor{gray}{нет} \\
\texttt{VIS-05} & Графики апостериорной неопределенности & \textcolor{gray}{нет} & \textcolor{gray}{нет} & \cellcolor{yes!30}да & \cellcolor{part!40}част. & \textcolor{gray}{нет} & \textcolor{gray}{нет} & \cellcolor{yes!30}да & \textcolor{gray}{нет} & \textcolor{gray}{нет} \\
\midrule \multicolumn{11}{@{}l}{\textbf{Специальные методы}} \\
\texttt{PROP-01} & Декларативное задание параметрических ограничений (RESTRICT) & \cellcolor{part!40}част. & \textcolor{gray}{нет} & \cellcolor{yes!30}да & \cellcolor{yes!30}да & \textcolor{gray}{нет} & \textcolor{gray}{нет} & \cellcolor{yes!30}да & \cellcolor{part!40}част. & \cellcolor{yes!30}да \\
\texttt{PROP-02} & Скорректированные тесты отношения правдоподобий (VLMR, LMR, MLRT) & \textcolor{gray}{нет} & \textcolor{gray}{нет} & \textcolor{gray}{нет} & \textcolor{gray}{нет} & \textcolor{gray}{нет} & \textcolor{gray}{нет} & \textcolor{gray}{нет} & \textcolor{gray}{нет} & \textcolor{gray}{нет} \\
\texttt{PROP-03} & Адаптивная регуляризация и стабилизация масштаба & \cellcolor{part!40}част. & \textcolor{gray}{нет} & \cellcolor{part!40}част. & \cellcolor{yes!30}да & \textcolor{gray}{нет} & \textcolor{gray}{нет} & \cellcolor{part!40}част. & \cellcolor{part!40}част. & \textcolor{gray}{нет} \\
\texttt{PROP-04} & Символьно-аналитическое оценивание смесей высокой точности & \textcolor{gray}{нет} & \textcolor{gray}{нет} & \textcolor{gray}{нет} & \textcolor{gray}{нет} & \textcolor{gray}{нет} & \textcolor{gray}{нет} & \textcolor{gray}{нет} & \textcolor{gray}{нет} & \textcolor{gray}{нет} \\
\texttt{PROP-05} & Автоматический эвристический перебор аналитических форм компонент & \textcolor{gray}{нет} & \textcolor{gray}{нет} & \textcolor{gray}{нет} & \textcolor{gray}{нет} & \textcolor{gray}{нет} & \textcolor{gray}{нет} & \textcolor{gray}{нет} & \textcolor{gray}{нет} & \textcolor{gray}{нет} \\
\midrule \multicolumn{11}{@{}l}{\textbf{Представление смеси}} \\
\texttt{REPR-01} & Поддержка вероятностных функций (PDF, CDF, PPF) & \cellcolor{part!40}част. & \cellcolor{part!40}част. & \textcolor{gray}{нет} & \cellcolor{part!40}част. & \cellcolor{yes!30}да & \cellcolor{yes!30}да & \textcolor{gray}{нет} & \textcolor{gray}{нет} & \textcolor{gray}{нет} \\
\texttt{REPR-02} & Сэмплирование & \cellcolor{yes!30}да & \textcolor{gray}{нет} & \cellcolor{yes!30}да & \cellcolor{yes!30}да & \cellcolor{yes!30}да & \cellcolor{yes!30}да & \cellcolor{yes!30}да & \cellcolor{yes!30}да & \textcolor{gray}{нет} \\
\texttt{REPR-03} & Аналитический расчет моментов смеси & \textcolor{gray}{нет} & \textcolor{gray}{нет} & \textcolor{gray}{нет} & \cellcolor{part!40}част. & \cellcolor{yes!30}да & \cellcolor{part!40}част. & \textcolor{gray}{нет} & \textcolor{gray}{нет} & \textcolor{gray}{нет} \\
\texttt{REPR-04} & Сериализация смеси в JSON & \cellcolor{yes!30}да & \textcolor{gray}{нет} & \textcolor{gray}{нет} & \cellcolor{part!40}част. & \cellcolor{part!40}част. & \textcolor{gray}{нет} & \textcolor{gray}{нет} & \cellcolor{yes!30}да & \textcolor{gray}{нет} \\
\midrule \multicolumn{11}{@{}l}{\textbf{Производительность}} \\
\texttt{PERF-01} & Параллелизация независимых вычислений & \cellcolor{part!40}част. & \textcolor{gray}{нет} & \cellcolor{part!40}част. & \cellcolor{yes!30}да & \cellcolor{part!40}част. & \textcolor{gray}{нет} & \cellcolor{yes!30}да & \cellcolor{yes!30}да & \textcolor{gray}{нет} \\
\texttt{PERF-02} & Кэширование промежуточных структур и плотностей & \textcolor{gray}{нет} & \textcolor{gray}{нет} & \cellcolor{part!40}част. & \cellcolor{part!40}част. & \cellcolor{part!40}част. & \cellcolor{yes!30}да & \cellcolor{part!40}част. & \cellcolor{part!40}част. & \textcolor{gray}{нет} \\
\texttt{PERF-03} & Векторизованные вычисления (NumPy, SIMD) & \cellcolor{yes!30}да & \textcolor{gray}{нет} & \cellcolor{part!40}част. & \cellcolor{yes!30}да & \cellcolor{yes!30}да & \cellcolor{yes!30}да & \textcolor{gray}{нет} & \cellcolor{yes!30}да & \textcolor{gray}{нет} \\
\texttt{PERF-04} & Контроль аллокаций памяти и потоковые вычисления & \cellcolor{yes!30}да & \textcolor{gray}{нет} & \textcolor{gray}{нет} & \cellcolor{yes!30}да & \cellcolor{yes!30}да & \textcolor{gray}{нет} & \cellcolor{part!40}част. & \cellcolor{yes!30}да & \cellcolor{yes!30}да \\
\midrule \multicolumn{11}{@{}l}{\textbf{Диагностика}} \\
\texttt{DIAG-01} & Статистические критерии согласия (KS, AD, Chi2) & \textcolor{gray}{нет} & \textcolor{gray}{нет} & \textcolor{gray}{нет} & \cellcolor{part!40}част. & \cellcolor{part!40}част. & \textcolor{gray}{нет} & \textcolor{gray}{нет} & \textcolor{gray}{нет} & \cellcolor{yes!30}да \\
\texttt{DIAG-02} & Стандартные ошибки и доверительные интервалы оценок & \textcolor{gray}{нет} & \textcolor{gray}{нет} & \cellcolor{yes!30}да & \cellcolor{yes!30}да & \textcolor{gray}{нет} & \textcolor{gray}{нет} & \cellcolor{yes!30}да & \textcolor{gray}{нет} & \cellcolor{yes!30}да \\
\texttt{DIAG-03} & Бенчмарки производительности (ASV, pytest-benchmark) & \textcolor{gray}{нет} & \textcolor{gray}{нет} & \textcolor{gray}{нет} & \cellcolor{yes!30}да & \textcolor{gray}{нет} & \textcolor{gray}{нет} & \textcolor{gray}{нет} & \textcolor{gray}{нет} & \textcolor{gray}{нет} \\
\texttt{DIAG-04} & Оценка отделимости и перекрытия компонент & \textcolor{gray}{нет} & \textcolor{gray}{нет} & \cellcolor{part!40}част. & \textcolor{gray}{нет} & \textcolor{gray}{нет} & \textcolor{gray}{нет} & \cellcolor{part!40}част. & \textcolor{gray}{нет} & \textcolor{gray}{нет} \\
\bottomrule
\end{longtable}}

\end{landscape}
\section{Выводы по разнородным смесям}
{\small
\begin{center}
\begin{tabular}{@{}l l r r r r@{}}
\toprule
Инструмент & Разнородность & Все, \% & Обяз., \% & \shortstack[r]{\scriptsize В тройке лучших\\\scriptsize по 5 разнор. сценариям} & \shortstack[r]{\scriptsize Среднее покрытие\\\scriptsize разнор. сценариев, \%} \\
\midrule
pomegranate & полностью & \cellcolor{hm!30}36 & \cellcolor{hm!47}60 & 0 & \cellcolor{hm!43}54 \\
mixem & полностью & \cellcolor{hm!18}19 & \cellcolor{hm!36}44 & 0 & \cellcolor{hm!36}44 \\
FlexMix & полностью & \cellcolor{hm!38}47 & \cellcolor{hm!55}72 & 3 & \cellcolor{hm!48}62 \\
brms & полностью & \cellcolor{hm!43}54 & \cellcolor{hm!51}66 & 4 & \cellcolor{hm!53}69 \\
Distributions.jl & полностью & \cellcolor{hm!29}35 & \cellcolor{hm!47}60 & 3 & \cellcolor{hm!48}61 \\
PySATL mpest & полностью & \cellcolor{hm!33}40 & \cellcolor{hm!62}81 & 3 & \cellcolor{hm!53}68 \\
StepMix & частично & \cellcolor{hm!32}39 & \cellcolor{hm!48}62 & 0 & \cellcolor{hm!40}49 \\
Rmixmod (mixmod) & частично & \cellcolor{hm!32}38 & \cellcolor{hm!42}53 & 0 & \cellcolor{hm!38}47 \\
mixdist & частично & \cellcolor{hm!28}33 & \cellcolor{hm!39}49 & 0 & \cellcolor{hm!31}38 \\
\bottomrule
\end{tabular}
\end{center}}
\begin{itemize}
\item Схему ECME среди разнородных средств реализует только \texttt{PySATL mpest}; блочную схему ECM --- \texttt{FlexMix}, \texttt{PySATL mpest}.
Остальные средства либо оптимизируют все параметры сразу (\texttt{pomegranate}, \texttt{mixem}), либо вообще не используют итерационное оценивание
(\texttt{brms} применяет байесовское сэмплирование).
\item Мультистарт, информационные критерии, фиксацию параметров и явные границы одновременно выполняет \texttt{brms}.
Этот инструмент не оценивает смесь итерационно по максимуму правдоподобия.
\item Сравнение конфигураций семейств (\texttt{STRUC-13}) не выполняет ни одно из этих средств; для разнородных смесей остаётся
комбинаторным перебором, который нужно ограничивать порядком и бюджетом.
\item Бутстреп отношения правдоподобий не реализован ни в одном из этих средств.
\end{itemize}

\part{Положение \texttt{pysatl-mpest}}
\begin{landscape}
\section{Пробелы}
\texttt{pysatl-mpest} выполняет полностью 31 из 86 требований, частично --- 6, не выполняет --- 49.
Покрытие: общее 40\,\%, обязательные 81\,\%, желательные 20\,\%, факультативные 0\,\%.
Возможности, которые больше не выполняет полностью ни один инструмент: \texttt{EST-06} Стратегия наблюдаемого правдоподобия на M-шаге; \texttt{STRUC-08} Оценивание распределений с неизвестным порогом; \texttt{INFRA-03} Property-based тестирование математических инвариантов.

Таблица содержит все требования, не выполненные полностью, в порядке приоритета, числа сценариев с критическим требованием и числа сценариев с желательным.
<<Др.>> --- сколько других инструментов выполняют требование полностью; <<Доноры>> --- первые из них (готовые решения, которые можно изучить).
{\scriptsize\setlength{\tabcolsep}{3pt}
\begin{longtable}{@{}l >{\raggedright\arraybackslash}p{8.2cm} l l r r r >{\raggedright\arraybackslash}p{6.2cm}@{}}
\toprule
ID & Требование & Приоритет & Статус & Др. & Крит. & Жел. & Доноры \\
\midrule
\endhead
\texttt{EST-01} & Каркас и фасад классического EM-алгоритма & обязательные & \cellcolor{part!40}част. & 22 & 10 & --- & scikit-learn (sklearn.mixture), pomegranate, mixem, StepMix \\
\texttt{STRUC-09} & Информационные критерии выбора числа компонент (AIC, BIC, CAIC) & обязательные & \cellcolor{part!40}част. & 15 & 10 & --- & scikit-learn (sklearn.mixture), StepMix, student-mixture, kasakh Mixture-Models \\
\texttt{INIT-03} & Кластеризация данных & обязательные & \textcolor{gray}{нет} & 16 & 2 & 2 & scikit-learn (sklearn.mixture), pomegranate, StepMix, student-mixture \\
\texttt{INIT-02} & Рандомизированный мультистарт & обязательные & \textcolor{gray}{нет} & 9 & 1 & 2 & scikit-learn (sklearn.mixture), StepMix, student-mixture, FlexMix \\
\texttt{EST-09} & Замкнутые аналитические формулы параметров на M-шаге & обязательные & \cellcolor{part!40}част. & 0 & 1 & --- & --- \\
\texttt{VIS-01} & Декомпозиция плотности смеси & обязательные & \textcolor{gray}{нет} & 15 & --- & 10 & StepMix, distfit, mclust (v5+), mixtools \\
\texttt{EST-13} & Явные параметрические границы & обязательные & \textcolor{gray}{нет} & 5 & --- & 3 & kasakh Mixture-Models, distfit, mixR, mixdist \\
\texttt{STRUC-11} & Параметрический бутстреп отношения правдоподобий (LRT) & обязательные & \textcolor{gray}{нет} & 7 & --- & 2 & StepMix, mclust (v5+), mixtools, mixR \\
\texttt{REPR-03} & Аналитический расчет моментов смеси & желательные & \cellcolor{part!40}част. & 2 & 5 & --- & sn (Skew-Normal/t), Distributions.jl \\
\texttt{INIT-04} & Квантильное детерминированное разбиение & желательные & \textcolor{gray}{нет} & 4 & 4 & 3 & mclust (v5+), mixtools, mixR, mixdist \\
\texttt{STRUC-13} & Сравнение и отбор конфигураций семейств & желательные & \textcolor{gray}{нет} & 1 & 3 & 1 & distfit \\
\texttt{STRUC-21} & Равенство и линейные ограничения параметров компонент & желательные & \textcolor{gray}{нет} & 8 & 2 & 3 & StepMix, mclust (v5+), FlexMix, Rmixmod (mixmod) \\
\texttt{DIAG-02} & Стандартные ошибки и доверительные интервалы оценок & желательные & \textcolor{gray}{нет} & 10 & 1 & 5 & StepMix, mclust (v5+), mixtools, FlexMix \\
\texttt{INIT-05} & Инициализация методом моментов по сегментам & желательные & \textcolor{gray}{нет} & 2 & 1 & 1 & mixR, sn (Skew-Normal/t) \\
\texttt{PERF-01} & Параллелизация независимых вычислений & желательные & \textcolor{gray}{нет} & 10 & 1 & --- & scikit-learn (sklearn.mixture), StepMix, mclust (v5+), Rmixmod (mixmod) \\
\texttt{STRUC-10} & Критерий интегрированного правдоподобия классификации (ICL) & желательные & \textcolor{gray}{нет} & 8 & 1 & --- & StepMix, kasakh Mixture-Models, mclust (v5+), FlexMix \\
\texttt{EST-16} & Оценивание по сгруппированным и интервальным данным & желательные & \textcolor{gray}{нет} & 3 & 1 & --- & mixR, rebmix, mixdist \\
\texttt{VIS-02} & Диагностические графики согласия (Q-Q, P-P) & желательные & \textcolor{gray}{нет} & 3 & --- & 10 & distfit, mixR, brms \\
\texttt{DIAG-01} & Статистические критерии согласия (KS, AD, Chi2) & желательные & \textcolor{gray}{нет} & 2 & --- & 10 & distfit, mixdist \\
\texttt{VIS-03} & Профиль выбора числа компонент (AIC, BIC, ICL) & желательные & \textcolor{gray}{нет} & 9 & --- & 9 & StepMix, distfit, mclust (v5+), FlexMix \\
\texttt{INIT-06} & Случайный выбор наблюдений & желательные & \textcolor{gray}{нет} & 19 & --- & 1 & scikit-learn (sklearn.mixture), pomegranate, StepMix, student-mixture \\
\texttt{EST-17} & Взвешенные наблюдения и веса выборки & желательные & \textcolor{gray}{нет} & 16 & --- & 1 & scikit-learn (sklearn.mixture), pomegranate, StepMix, mclust (v5+) \\
\texttt{STRUC-16} & Объединение перекрывающихся компонент по критерию унимодальности (DIP-тест) & желательные & \textcolor{gray}{нет} & 1 & --- & 1 & fpc \\
\texttt{EST-22} & Покомпонентный EM-алгоритм (CW-EM, SAGE) & желательные & \textcolor{gray}{нет} & 0 & --- & 1 & --- \\
\texttt{TECH-05} & scikit-learn совместимый интерфейс & желательные & \textcolor{gray}{нет} & 8 & --- & --- & scikit-learn (sklearn.mixture), StepMix, student-mixture, rebmix \\
\texttt{EST-04} & Direct MLE & желательные & \textcolor{gray}{нет} & 6 & --- & --- & kasakh Mixture-Models, distfit, mixdist, sn (Skew-Normal/t) \\
\texttt{REPR-04} & Сериализация смеси в JSON & желательные & \textcolor{gray}{нет} & 4 & --- & --- & pomegranate, distfit, Rmixmod (mixmod), mlpack (mlpack::gmm) \\
\texttt{STRUC-15} & Динамические переходы слияния и расщепления (SMEM) & желательные & \textcolor{gray}{нет} & 3 & --- & --- & rebmix, GaussianMixtures.jl, bnpy \\
\texttt{EST-07} & Стратегия взвешенного метода моментов на M-шаге & желательные & \cellcolor{part!40}част. & 1 & --- & --- & sn (Skew-Normal/t) \\
\texttt{STRUC-17} & Непараметрическое градиентное нахождение числа компонент (NPMLE) & желательные & \textcolor{gray}{нет} & 1 & --- & --- & CAMAN \\
\texttt{EST-08} & Стратегия метода L-моментов на M-шаге & желательные & \cellcolor{part!40}част. & 0 & --- & --- & --- \\
\texttt{STRUC-14} & Жадное инкрементальное обучение порядка смеси & желательные & \textcolor{gray}{нет} & 0 & --- & --- & --- \\
\texttt{STRUC-18} & Предсказание порядка смеси через мета-обучение (Meta-Learning) & желательные & \textcolor{gray}{нет} & 0 & --- & --- & --- \\
\texttt{STRUC-19} & Классификация параметрических семейств с помощью нейросетей & желательные & \textcolor{gray}{нет} & 0 & --- & --- & --- \\
\texttt{STRUC-20} & Амортизированный нейросетевой вывод параметров смеси (Deep Sets) & желательные & \textcolor{gray}{нет} & 0 & --- & --- & --- \\
\texttt{STRUC-22} & Асимметричные и тяжелохвостые семейства (Skew-Normal, Student-t) & факультативные & \textcolor{gray}{нет} & 4 & 2 & --- & student-mixture, teigen, sn (Skew-Normal/t), brms \\
\texttt{EST-19} & Ускорение сходимости EM-алгоритма (Aitken, SQUAREM, PX-EM) & факультативные & \textcolor{gray}{нет} & 0 & 1 & 1 & --- \\
\texttt{PERF-04} & Контроль аллокаций памяти и потоковые вычисления & факультативные & \textcolor{gray}{нет} & 14 & 1 & --- & scikit-learn (sklearn.mixture), pomegranate, mixR, rebmix \\
\texttt{EST-18} & Оценивание по цензурированным и усеченным данным & факультативные & \textcolor{gray}{нет} & 2 & 1 & --- & mixR, brms \\
\texttt{VIS-05} & Графики апостериорной неопределенности & факультативные & \textcolor{gray}{нет} & 5 & --- & 4 & StepMix, mclust (v5+), mixtools, FlexMix \\
\texttt{DIAG-04} & Оценка отделимости и перекрытия компонент & факультативные & \textcolor{gray}{нет} & 1 & --- & 4 & fpc \\
\texttt{EST-24} & Байесовское сэмплирование методом расширения данных (Gibbs, MCMC) & факультативные & \textcolor{gray}{нет} & 1 & --- & 4 & brms \\
\texttt{VIS-04} & Визуализация скорости сходимости & факультативные & \textcolor{gray}{нет} & 5 & --- & 2 & mixtools, Rmixmod (mixmod), brms, bnpy \\
\texttt{EST-20} & Рекурсивный потоковый EM-алгоритм & факультативные & \textcolor{gray}{нет} & 1 & --- & 1 & bnpy \\
\texttt{DIAG-03} & Бенчмарки производительности (ASV, pytest-benchmark) & факультативные & \textcolor{gray}{нет} & 4 & --- & --- & ClusterR, brms, mlpack (mlpack::gmm), bnpy \\
\texttt{PROP-01} & Декларативное задание параметрических ограничений (RESTRICT) & факультативные & \textcolor{gray}{нет} & 4 & --- & --- & StepMix, FlexMix, mixdist, brms \\
\texttt{PROP-03} & Адаптивная регуляризация и стабилизация масштаба & факультативные & \textcolor{gray}{нет} & 4 & --- & --- & scikit-learn (sklearn.mixture), mclust (v5+), brms, bnpy \\
\texttt{EST-10} & Стохастический и классификационный EM-алгоритм & факультативные & \textcolor{gray}{нет} & 3 & --- & --- & mixtools, FlexMix, Rmixmod (mixmod) \\
\texttt{STRUC-23} & Вариационный байесовский вывод (VBEM) & факультативные & \textcolor{gray}{нет} & 3 & --- & --- & scikit-learn (sklearn.mixture), student-mixture, bnpy \\
\texttt{EST-11} & Алгоритм Rough-Enhanced-Bayes & факультативные & \textcolor{gray}{нет} & 1 & --- & --- & rebmix \\
\texttt{EST-23} & Робастное оценивание по усеченному правдоподобию (TCLUST, C-шаг) & факультативные & \textcolor{gray}{нет} & 1 & --- & --- & fpc \\
\texttt{EST-21} & Оценивание методом минимального расстояния (Hellinger, Wasserstein) & факультативные & \textcolor{gray}{нет} & 0 & --- & --- & --- \\
\texttt{PROP-02} & Скорректированные тесты отношения правдоподобий (VLMR, LMR, MLRT) & факультативные & \textcolor{gray}{нет} & 0 & --- & --- & --- \\
\texttt{PROP-04} & Символьно-аналитическое оценивание смесей высокой точности & факультативные & \textcolor{gray}{нет} & 0 & --- & --- & --- \\
\texttt{PROP-05} & Автоматический эвристический перебор аналитических форм компонент & факультативные & \textcolor{gray}{нет} & 0 & --- & --- & --- \\
\bottomrule
\end{longtable}}

\end{landscape}
\section{Приоритеты развития}\label{sec:roadmap}
Чтобы упорядочить пробелы по пользе для прикладных задач, для каждого невыполненного требования рассчитан средний прирост покрытия критических требований по 10 сценариям:
прирост равен сумме по сценариям, где требование критическое, величины (недостающая доля требования)/(число критических требований сценария), делённой на 10.
Таким образом, требование ценнее, если оно критично для многих сценариев с коротким списком критических требований. Расчёт не учитывает трудоёмкость реализации.

В таблице: <<Сцен.>> --- число сценариев с критическим требованием, <<Прирост>> --- вклад требования в среднее покрытие (процентные пункты), <<Накопл.>> --- среднее покрытие после закрытия этого и всех предыдущих пробелов (\%).

Исходное среднее покрытие критических требований по сценариям: 69\,\%. Если довести до полного выполнения все обязательные требования (в том числе выполненные частично), оно станет 82\,\%,
общее покрытие --- 47\,\%; если закрыть все критические требования сценариев --- 100\,\%.
{\footnotesize
\begin{longtable}{@{}r l >{\raggedright\arraybackslash}p{5.8cm} l r r r@{}}
\toprule
\# & ID & Требование & Приоритет & Сцен. & Прирост & Накопл. \\
\midrule
\endhead
1 & \texttt{EST-01} & Каркас и фасад классического EM-алгоритма & обязательные & 10 & 4.7 & 74.0 \\
2 & \texttt{STRUC-09} & Информационные критерии выбора числа компонент (AIC, BIC, CAIC) & обязательные & 10 & 4.7 & 78.7 \\
3 & \texttt{INIT-04} & Квантильное детерминированное разбиение & желательные & 4 & 3.3 & 82.1 \\
4 & \texttt{STRUC-22} & Асимметричные и тяжелохвостые семейства (Skew-Normal, Student-t) & факультативные & 2 & 2.4 & 84.5 \\
5 & \texttt{STRUC-13} & Сравнение и отбор конфигураций семейств & желательные & 3 & 2.4 & 86.9 \\
6 & \texttt{INIT-03} & Кластеризация данных & обязательные & 2 & 2.0 & 88.9 \\
7 & \texttt{REPR-03} & Аналитический расчет моментов смеси & желательные & 5 & 1.9 & 90.8 \\
8 & \texttt{STRUC-21} & Равенство и линейные ограничения параметров компонент & желательные & 2 & 1.8 & 92.6 \\
9 & \texttt{INIT-02} & Рандомизированный мультистарт & обязательные & 1 & 1.1 & 93.7 \\
10 & \texttt{STRUC-10} & Критерий интегрированного правдоподобия классификации (ICL) & желательные & 1 & 1.0 & 94.7 \\
11 & \texttt{INIT-05} & Инициализация методом моментов по сегментам & желательные & 1 & 0.8 & 95.5 \\
12 & \texttt{DIAG-02} & Стандартные ошибки и доверительные интервалы оценок & желательные & 1 & 0.8 & 96.3 \\
\bottomrule
\end{longtable}}
Первые двенадцать пробелов поднимают среднее покрытие с 69\,\% до 96\,\%.
Порядок следует читать как ориентир: он основан на критичности для выбранных сценариев, а не на оценке стоимости разработки.

\section{Что ещё следует из таблиц}
\begin{itemize}
\item Обязательные требования, не выполненные \texttt{pysatl-mpest}: кластеризация данных; рандомизированный мультистарт; декомпозиция плотности смеси; явные параметрические границы; параметрический бутстреп отношения правдоподобий (LRT). Для каждого из них в таблице пробелов указаны инструменты-доноры.
\item Критические требования не менее чем двух сценариев, выполненные не полностью: каркас и фасад классического EM-алгоритма; кластеризация данных; квантильное детерминированное разбиение; аналитический расчет моментов смеси; информационные критерии выбора числа компонент (AIC, BIC, CAIC); сравнение и отбор конфигураций семейств; равенство и линейные ограничения параметров компонент; асимметричные и тяжелохвостые семейства (Skew-Normal, Student-t).
\item Связь методов EM, ECM и ECME с объёмом данных и числом компонент в базе не оценивалась: сценарии различаются выборкой
(малая, умеренная, большая, очень большая) и знанием числа компонент, но оценки фиксируют наличие метода, а не его поведение.
Выбор метода под класс сценария проверяется экспериментально.
\end{itemize}

\end{document}
